% Amélioration de la santé reproductive — SVT 1ère A — Cours signé OnBuch+
% Programme officiel MINESEC. Compiler avec : tectonic NA06-amelioration-sante-reproductive.tex
\newcommand{\DOCMATIERE}{SVT}
\newcommand{\DOCNIVEAU}{1ère A}
\newcommand{\DOCMODULE}{Module 2 : L'Éducation à la Santé}
\newcommand{\DOCLECON}{NA06}
\newcommand{\DOCTITRE}{Amélioration de la santé reproductive}
% OnBuch+ — préambule « cours signés » ENRICHI (compilé avec tectonic / XeLaTeX)
% Système visuel « Pop » d'OnBuch+ : fond crème (jamais blanc), bordures encre
% épaisses, ombres « dures » décalées, titres Archivo Black en capitales.
% Version MATHÉMATIQUES : graphes (pgfplots/tikz), boîtes pédagogiques
% (définition, propriété, méthode, attention, le savais-tu, exercice résolu,
% exercice, TP, bilan), page de garde et signature OnBuch+ ancrée en bas.
%
% Macros attendues dans main.tex AVANT \input{preamble} :
%   \DOCMATIERE  \DOCNIVEAU  \DOCMODULE  \DOCLECON (numéro)  \DOCTITRE
\documentclass[11pt]{article}

\usepackage{fontspec}
\usepackage[french]{babel}
\usepackage[a4paper,margin=2cm,top=3.4cm,bottom=2.8cm,headheight=1.4cm,footskip=1.3cm]{geometry}
\usepackage{amsmath,amssymb,amsfonts}
\usepackage{mathtools} % \xmapsto, \coloneqq... souvent utilisés par les modèles
\usepackage{cancel} % simplifications barrées (bilans de forces)
% \abs{z} (module, valeur absolue) et \norm{u} (norme), utilisés par les modèles
\providecommand{\abs}{}\let\abs\relax\DeclarePairedDelimiter{\abs}{\lvert}{\rvert}
\providecommand{\norm}{}\let\norm\relax\DeclarePairedDelimiter{\norm}{\lVert}{\rVert}
\usepackage[table]{xcolor} % [table] : fournit aussi \rowcolors (alternance de lignes)
\usepackage{pagecolor}
\usepackage{tikz}
\usetikzlibrary{arrows.meta,positioning,calc,shapes.geometric,shapes.misc,shapes.arrows,decorations.pathmorphing,decorations.pathreplacing,decorations.markings,patterns,shadows,mindmap,trees,fit,backgrounds,matrix,intersections,angles,quotes}
\usepackage{pgfplots}
\usepackage[version=4]{mhchem} % équations nucléaires \ce{^{235}_{92}U}
\usepackage[european,siunitx]{circuitikz} % schémas électriques
\pgfplotsset{compat=1.18}
\usepgfplotslibrary{fillbetween,groupplots} % aires entre courbes (calcul intégral)
\usepackage{enumitem}
\usepackage{fancyhdr}
% [most] charge toutes les bibliothèques tcolorbox (listings, minted, raster,
% documentation...) et gonfle inutilement la mémoire TeX sur un document long
% (~300 boîtes) : on ne charge que ce qui est réellement utilisé.
\usepackage[skins,breakable]{tcolorbox}
\usepackage{graphicx}
\usepackage{colortbl}
\usepackage{booktabs}
\usepackage{tabularx}
\usepackage{diagbox} % case d'en-tête diagonale des tableaux à double entrée
% Colonnes extensibles centrée / gauche / droite (C, L, R), souvent utilisées par les modèles
\newcolumntype{C}{>{\centering\arraybackslash}X}
\newcolumntype{L}{>{\raggedright\arraybackslash}X}
\newcolumntype{R}{>{\raggedleft\arraybackslash}X}
\usepackage{array}
\usepackage{multicol}
\usepackage{siunitx}
% parse-numbers=false : accepte aussi \SI{2,5 \times 10^{-3}}{...}
\sisetup{locale=FR,output-decimal-marker={,},group-minimum-digits=4,parse-numbers=false}
\DeclareSIUnit\ohm{\ensuremath{\symup{\Omega}}} % U+2126 absent de la police
\DeclareSIUnit\division{div} % réglages d'oscilloscope (ms/div, V/div)
\DeclareSIUnit\tr{tr} % tours (vitesse de rotation, ex. \SI{1500}{\tr\per\minute})
\DeclareSIUnit\ch{ch} % cheval-vapeur (puissance moteur, unité usuelle hors SI)
\DeclareSIUnit\franccfa{FCFA} % franc CFA (coûts, contexte camerounais)
\DeclareSIUnit\dioptre{\ensuremath{\delta}} % vergence d'une lentille (dioptrie)
\usepackage{qrcode}
% \degree : commande fréquemment utilisée par les modèles pour un angle ou une
% température en dehors de \SI{}{} (non fournie par les paquets chargés).
\providecommand{\degree}{^{\circ}}
% \qed (fin de démonstration), utilisé par les modèles sans amsthm
\providecommand{\qed}{\hfill$\square$}
% \barre : double barre des valeurs interdites dans un tableau de variations (tkz-tab)
\providecommand{\barre}{\ensuremath{\|}}
% environnement proof (amsthm non chargé) utilisé par les modèles
\ifdefined\proof\else\newenvironment{proof}[1][Démonstration]{\par\noindent\textit{#1.}\ }{\qed\par}\fi
\usepackage{lastpage}
\usepackage{hyperref}
\hypersetup{hidelinks}

% ============================================================
% Palette « Pop » OnBuch+ — mêmes hex que la classe P (plus_ui.dart)
% ============================================================
\definecolor{poporange}{HTML}{F59321}
\definecolor{popdark}{HTML}{E07A0C}
\definecolor{popcream}{HTML}{FFF6E8}
\definecolor{popink}{HTML}{1C1714}
\definecolor{poppurple}{HTML}{7A5AE0}
\definecolor{popgreen}{HTML}{2E9E6A}
\definecolor{poppink}{HTML}{E0457B}
\definecolor{popblue}{HTML}{2F7DE1}
\definecolor{popgold}{HTML}{C98A12}
\definecolor{popmuted}{HTML}{8A7F76}
\definecolor{popline}{HTML}{EADFD2}
\definecolor{popgray}{HTML}{8A7F76}
\colorlet{popgrey}{popgray}
% Alias tolérants (le modèle nomme parfois les couleurs différemment)
\colorlet{popwhite}{white}
\colorlet{popblack}{popink}
\colorlet{popred}{poppink}
\colorlet{popyellow}{popgold}
% Teintes claires pour fonds de tableaux / zones
\colorlet{popblueL}{popblue!10!white}
\colorlet{poporangeL}{poporange!14!white}
\colorlet{popgreenL}{popgreen!12!white}
\colorlet{poppurpleL}{poppurple!12!white}
\colorlet{poppinkL}{poppink!12!white}
\colorlet{popgoldL}{popgold!14!white}

\pagecolor{popcream}

% ============================================================
% Typographie
% ============================================================
\defaultfontfeatures{Ligatures=TeX}
\setmainfont{PlusJakartaSans-Medium.ttf}[Path=../../../fonts/,BoldFont=PlusJakartaSans-Bold.ttf]
% Maths en sans-serif (Fira Math) avec chiffres et lettres droites de Plus
% Jakarta, pour que les formules se fondent dans le texte.
\usepackage[math-style=ISO,bold-style=ISO]{unicode-math}
\setmathfont{FiraMath-Regular.otf}[Path=../../../fonts/]
\setmathfont{PlusJakartaSans-Medium.ttf}[Path=../../../fonts/,range={up/num,up/{Latin,latin},\mathup}]
\usepackage{icomma} % virgule décimale sans espace en mode math
% Caractères mathématiques Unicode absents des polices de texte (Plus Jakarta,
% Archivo Black) : rendus par leur équivalent mathématique, en texte comme en
% formule — sinon ils s'affichent en carré vide (ex. « Arithmétique dans ℤ »).
\usepackage{newunicodechar}
\newunicodechar{ℝ}{\ensuremath{\mathbb{R}}}
\newunicodechar{ℤ}{\ensuremath{\mathbb{Z}}}
\newunicodechar{ℂ}{\ensuremath{\mathbb{C}}}
\newunicodechar{ℕ}{\ensuremath{\mathbb{N}}}
\newunicodechar{ℚ}{\ensuremath{\mathbb{Q}}}
\newunicodechar{𝔻}{\ensuremath{\mathbb{D}}}
\newunicodechar{α}{\ensuremath{\alpha}}
\newunicodechar{β}{\ensuremath{\beta}}
\newunicodechar{γ}{\ensuremath{\gamma}}
\newunicodechar{δ}{\ensuremath{\delta}}
\newunicodechar{ε}{\ensuremath{\varepsilon}}
\newunicodechar{θ}{\ensuremath{\theta}}
\newunicodechar{λ}{\ensuremath{\lambda}}
\newunicodechar{μ}{\ensuremath{\mu}}
\newunicodechar{σ}{\ensuremath{\sigma}}
\newunicodechar{τ}{\ensuremath{\tau}}
\newunicodechar{φ}{\ensuremath{\varphi}}
\newunicodechar{ψ}{\ensuremath{\psi}}
\newunicodechar{ω}{\ensuremath{\omega}}
\newunicodechar{Σ}{\ensuremath{\Sigma}}
% majuscules produites par \MakeUppercase dans les titres (α-aminés -> Α)
\newunicodechar{Α}{\ensuremath{\alpha}}
\newunicodechar{Β}{\ensuremath{\beta}}
\newunicodechar{Γ}{\ensuremath{\Gamma}}
\newunicodechar{Δ}{\ensuremath{\Delta}}
\newunicodechar{Ε}{\ensuremath{\varepsilon}}
\newunicodechar{Θ}{\ensuremath{\Theta}}
\newunicodechar{Λ}{\ensuremath{\Lambda}}
\newunicodechar{Μ}{\ensuremath{\mu}}
\newunicodechar{Τ}{\ensuremath{\tau}}
\newunicodechar{Φ}{\ensuremath{\Phi}}
\newunicodechar{Ψ}{\ensuremath{\Psi}}
\newunicodechar{Ω}{\ensuremath{\Omega}}
\newunicodechar{↦}{\ensuremath{\mapsto}}
\newunicodechar{∈}{\ensuremath{\in}}
\newunicodechar{∉}{\ensuremath{\notin}}
\newunicodechar{∧}{\ensuremath{\wedge}}
\newunicodechar{⇔}{\ensuremath{\Leftrightarrow}}
\newunicodechar{⟺}{\ensuremath{\Longleftrightarrow}}
\newunicodechar{⇒}{\ensuremath{\Rightarrow}}
\newunicodechar{⟹}{\ensuremath{\Longrightarrow}}
\newunicodechar{∘}{\ensuremath{\circ}}
\newunicodechar{∪}{\ensuremath{\cup}}
\newunicodechar{∩}{\ensuremath{\cap}}
\newunicodechar{≡}{\ensuremath{\equiv}}
\newunicodechar{⊂}{\ensuremath{\subset}}
\newunicodechar{⊥}{\ensuremath{\perp}}
\newunicodechar{∃}{\ensuremath{\exists}}
\newunicodechar{∀}{\ensuremath{\forall}}
\newunicodechar{∛}{\ensuremath{\sqrt[3]{\,}}}
\newunicodechar{∜}{\ensuremath{\sqrt[4]{\,}}}
\newunicodechar{⃗}{\ensuremath{{}^{\to}}}
\newunicodechar{ⁿ}{\ifmmode^{n}\else\textsuperscript{\ensuremath{n}}\fi}
\newunicodechar{ˣ}{\ifmmode^{x}\else\textsuperscript{\ensuremath{x}}\fi}
\newunicodechar{ᵇ}{\ifmmode^{b}\else\textsuperscript{\ensuremath{b}}\fi}
\newunicodechar{ʳ}{\ifmmode^{r}\else\textsuperscript{\ensuremath{r}}\fi}
\newunicodechar{ᵘ}{\ifmmode^{u}\else\textsuperscript{\ensuremath{u}}\fi}
\newunicodechar{ʸ}{\ifmmode^{y}\else\textsuperscript{\ensuremath{y}}\fi}
\newunicodechar{ᵃ}{\ifmmode^{a}\else\textsuperscript{\ensuremath{a}}\fi}
\newunicodechar{ᵉ}{\ifmmode^{e}\else\textsuperscript{\ensuremath{e}}\fi}
\newunicodechar{ᵏ}{\ifmmode^{k}\else\textsuperscript{\ensuremath{k}}\fi}
\newunicodechar{ᵖ}{\ifmmode^{p}\else\textsuperscript{\ensuremath{p}}\fi}
\newunicodechar{ᴺ}{\ifmmode^{N}\else\textsuperscript{\ensuremath{N}}\fi}
\newunicodechar{ᶜ}{\ifmmode^{c}\else\textsuperscript{\ensuremath{c}}\fi}
\newunicodechar{ʰ}{\ifmmode^{h}\else\textsuperscript{\ensuremath{h}}\fi}
\newunicodechar{ᵠ}{\ifmmode^{q}\else\textsuperscript{\ensuremath{q}}\fi}
\newunicodechar{ₙ}{\ifmmode_{n}\else\textsubscript{\ensuremath{n}}\fi}
\newunicodechar{ₐ}{\ifmmode_{a}\else\textsubscript{\ensuremath{a}}\fi}
\newunicodechar{ᵢ}{\ifmmode_{i}\else\textsubscript{\ensuremath{i}}\fi}
\newunicodechar{ᵣ}{\ifmmode_{r}\else\textsubscript{\ensuremath{r}}\fi}
\newunicodechar{ₖ}{\ifmmode_{k}\else\textsubscript{\ensuremath{k}}\fi}
\newunicodechar{ᵦ}{\ifmmode_{\beta}\else\textsubscript{\ensuremath{\beta}}\fi}
\newunicodechar{ₓ}{\ifmmode_{x}\else\textsubscript{\ensuremath{x}}\fi}
\newfontfamily\popdisplay{ArchivoBlack-Regular.ttf}[Path=../../../fonts/]
\newfontfamily\poplogo{SpaceGrotesk-Bold.ttf}[Path=../../../fonts/]
\newfontfamily\popsemi{PlusJakartaSans-SemiBold.ttf}[Path=../../../fonts/]
\newfontfamily\popxbold{PlusJakartaSans-ExtraBold.ttf}[Path=../../../fonts/]
\newfontfamily\popxboldls{PlusJakartaSans-ExtraBold.ttf}[Path=../../../fonts/,LetterSpace=3.0]

\setlist[enumerate]{leftmargin=1.7em,itemsep=5pt,topsep=4pt}
\setlist[itemize]{leftmargin=1.7em,itemsep=4pt,topsep=4pt,label={\color{poporange}\textbullet}}
\setlength{\parindent}{0pt}
\setlength{\parskip}{5pt}
\renewcommand{\arraystretch}{1.25}

% pgfplots : style Pop par défaut
\pgfplotsset{
  every axis/.append style={axis line style={popink,line width=1pt},tick style={popink},
    grid style={popline,line width=0.5pt},label style={font=\small},tick label style={font=\footnotesize}},
  popcurve/.style={poporange,line width=1.8pt,no markers,smooth},
}
% Même style disponible dans un \draw TikZ simple (hors pgfplots)
\tikzset{popcurve/.style={poporange,line width=1.8pt}}
\tikzset{popgrid/.style={popline,very thin}}
\colorlet{popcurve}{poporange} % le nom du style est parfois utilisé comme couleur

% ============================================================
% Pill / squiggle
% ============================================================
\newcommand{\poppill}[3]{%
  \tikz[baseline=-0.55ex]\node[draw=popink,line width=1.2pt,rounded corners=2.6pt,%
    fill=#2,inner xsep=6.5pt,inner ysep=3pt]{%
    \popxboldls\fontsize{7}{7}\selectfont\color{#3}\MakeUppercase{#1}};%
}
\newcommand{\popsquiggle}[1]{%
  \tikz[baseline=-0.4ex]{
    \draw[#1,line width=2.6pt,line cap=round]
      plot [smooth] coordinates {(0,0.09) (0.34,0.01) (0.68,0.21) (1.02,0.01) (1.36,0.10)};
  }%
}
% Mot-clé surligné (marqueur orange sous le texte)
\newcommand{\cle}[1]{{\popxbold\color{popdark}#1}}

% ============================================================
% En-tête / pied
% ============================================================
\pagestyle{fancy}
\fancyhf{}
\renewcommand{\headrulewidth}{0pt}
\renewcommand{\footrulewidth}{0pt}
\lhead{\raisebox{-0.15ex}{\poplogo\fontsize{13}{13}\selectfont\color{popink}OnBuch\color{poporange}+}}
\rhead{\poppill{\DOCMATIERE\ \textbullet\ \DOCNIVEAU\ \textbullet\ Leçon \DOCLECON}{white}{popink}}
\lfoot{\popxbold\fontsize{7.6}{7.6}\selectfont\color{popmuted}\MakeUppercase{Cours signé OnBuch+}\ \textbullet\ {\popsemi\DOCTITRE}}
\rfoot{\tikz[baseline=-0.6ex]\node[rounded corners=8.5pt,fill=popink,minimum height=17pt,minimum width=17pt,inner xsep=5pt,inner sep=0]{\popxbold\fontsize{8}{8}\selectfont\color{popcream}\thepage\,/\,\pageref*{LastPage}};}

% ============================================================
% Titres
% ============================================================
\newcommand{\chaptitre}[2]{%
  \begin{tcolorbox}[enhanced,colback=poporange,colframe=popink,boxrule=2.6pt,arc=11pt,
      drop shadow=popink,left=15pt,right=15pt,top=12pt,bottom=11pt,before skip=3pt,after skip=18pt]
    \poppill{#2}{popink}{popcream}\par\vspace{9pt}
    {\raggedright\hyphenpenalty=10000\exhyphenpenalty=10000\popdisplay\fontsize{22}{24}\selectfont\color{popcream}\MakeUppercase{#1}\par}
  \end{tcolorbox}
}
% Section numérotée : pastille orange avec le numéro + titre Archivo
\newcounter{popsec}
\newcommand{\coursec}[1]{%
  \stepcounter{popsec}\setcounter{popsub}{0}%
  \par\vspace{16pt}\noindent
  \begin{minipage}{\linewidth}
  \tikz[baseline=-0.7ex]\node[draw=popink,line width=1.6pt,fill=poporange,rounded corners=4pt,minimum size=22pt,inner sep=0]{\popdisplay\fontsize{12}{12}\selectfont\color{popcream}\thepopsec};%
  \hspace{8pt}{\popdisplay\fontsize{15}{17}\selectfont\color{popink}\MakeUppercase{#1}}\par
  \vspace{4pt}\noindent{\color{poporange}\rule{36pt}{3.6pt}}
  \end{minipage}\par\nopagebreak\vspace{8pt}
}
\newcounter{popsub}
\newcommand{\courssub}[1]{%
  \stepcounter{popsub}%
  \par\vspace{9pt}\noindent{\popxbold\fontsize{11.5}{13}\selectfont\color{popink}{\color{poporange}\thepopsec.\thepopsub}\hspace{6pt}#1}\par\nopagebreak\vspace{4pt}
}

% ============================================================
% Boîtes pédagogiques — même recette : fond blanc, bordure épaisse colorée,
% ombre dure encre, badge plein. Titre optionnel [#1].
% ============================================================
\tcbset{popbase/.style={breakable,enhanced,colback=white,boxrule=2.3pt,arc=9pt,
  shadow={3pt}{-3pt}{0pt}{popink},left=14pt,right=14pt,top=11pt,bottom=11pt,before skip=11pt,after skip=15pt,
  fontupper=\popsemi\fontsize{10.6}{14.5}\selectfont\color{popink},
  fontlower=\popsemi\fontsize{10.6}{14.5}\selectfont\color{popink}}}
% Coche dessinée (le glyphe ✓ n'existe pas dans les polices du document)
\newcommand{\popcheck}[1]{\tikz[baseline=-0.5ex]\draw[#1,line width=1.6pt,line cap=round,line join=round](-0.13,0.0)--(-0.04,-0.09)--(0.13,0.1);}
\providecommand{\coche}{\popcheck{popgreen}}
\providecommand{\resultat}[1]{\par\smallskip\noindent\fcolorbox{popgreen}{popgreenL}{\parbox{\dimexpr\linewidth-2\fboxsep-2\fboxrule\relax}{\textbf{Résultat.} #1}}\par\smallskip}
\newcommand{\popboxhead}[3]{\poppill{#1}{#2}{white}\ifx\relax#3\relax\else\hspace{6pt}{\popxbold\fontsize{10.6}{12}\selectfont\color{popink}#3}\fi\par\vspace{7pt}}

\newtcolorbox{aretenir}[1][]{popbase,colframe=popgreen,before upper={\popboxhead{À retenir}{popgreen}{#1}}}
\newtcolorbox{exemplebox}[1][]{popbase,colframe=poporange,before upper={\popboxhead{Exemple}{poporange}{#1}}}
\newtcolorbox{definition}[1][]{popbase,colframe=popblue,before upper={\popboxhead{Définition}{popblue}{#1}}}
\newtcolorbox{propriete}[1][]{popbase,colframe=poppurple,before upper={\popboxhead{Propriété}{poppurple}{#1}}}
\newtcolorbox{methode}[1][]{popbase,colframe=popgold,before upper={\popboxhead{Méthode}{popgold}{#1}}}
\newtcolorbox{attention}[1][]{popbase,colframe=poppink,before upper={\popboxhead{Attention}{poppink}{#1}}}
\newtcolorbox{experience}[1][]{popbase,colframe=popblue,colback=popblueL,before upper={\popboxhead{Expérience}{popblue}{#1}}}
% « Le savais-tu ? » : carte violette pleine (contexte, culture, Cameroun)
\newtcolorbox{savaistu}[1][]{popbase,colback=poppurple,colframe=popink,
  fontupper=\popsemi\fontsize{10.4}{14.2}\selectfont\color{white},
  before upper={\poppill{Le savais-tu\,?}{white}{poppurple}\ifx\relax#1\relax\else\hspace{6pt}{\popxbold\color{white}#1}\fi\par\vspace{7pt}}}
% Exercice résolu : énoncé en haut, \tcblower puis la solution sur fond crème
\newcounter{popexo}\newcounter{popexr}
\newtcolorbox{exoresolu}[1][]{popbase,colframe=poporange,colbacklower=popcream,
  segmentation style={popink,line width=1.2pt,dashed},
  before upper={\stepcounter{popexr}\popboxhead{Exercice résolu \thepopexr}{poporange}{#1}},
  before lower={\poppill{Solution}{popink}{popcream}\par\vspace{6pt}}}
% Exercice (non résolu en place) — la correction est dans la partie Corrigés
\newtcolorbox{exercice}[1][]{popbase,colframe=popink,
  before upper={\stepcounter{popexo}\popboxhead{Exercice \thepopexo}{popink}{#1}}}
% Corrigé
\newtcolorbox{corrige}[1][]{popbase,colframe=popgreen,colback=popgreenL,
  before upper={\popboxhead{Corrigé}{popgreen}{#1}}}
% Objectifs / prérequis / bilan
\newtcolorbox{objectifs}{popbase,colframe=popink,colback=poporangeL,
  before upper={\popboxhead{Objectifs de la leçon}{popink}{}}}
\newtcolorbox{prerequis}{popbase,colframe=popink,colback=popblueL,
  before upper={\popboxhead{Prérequis}{popblue}{}}}
\newtcolorbox{bilan}{popbase,colframe=popink,colback=white,boxrule=2.8pt,
  before upper={\popboxhead{Fiche bilan}{poporange}{L'essentiel à connaître pour l'examen}}}
% Figure encadrée avec légende
\newtcolorbox{popfigure}[1][]{enhanced,breakable=false,colback=white,colframe=popline,boxrule=1.4pt,arc=8pt,
  left=10pt,right=10pt,top=10pt,bottom=8pt,before skip=10pt,after skip=12pt,halign=center,
  lower separated=false,halign lower=center,
  fontlower=\popsemi\fontsize{9}{11.5}\selectfont\color{popmuted},
  #1}
\newcommand{\legende}[1]{\par\vspace{6pt}{\popsemi\fontsize{9}{11.5}\selectfont\color{popmuted}#1}}

% ============================================================
% Signature OnBuch+
% ============================================================
\newcommand{\poplogoinline}[1]{{\poplogo\fontsize{#1}{#1}\selectfont\color{popink}OnBuch\color{poporange}+}}
\newcommand{\signatureonbuch}{%
  \begin{tcolorbox}[enhanced,colback=popink,colframe=popink,boxrule=2.2pt,arc=10pt,
      drop shadow=poporange,left=14pt,right=14pt,top=10pt,bottom=10pt,before skip=0pt,after skip=0pt]
    \tikz[baseline=-0.7ex]\node[circle,fill=popgreen,draw=popcream,line width=1.3pt,minimum size=22pt,inner sep=0]{\popcheck{white}};%
    \hspace{9pt}\begin{minipage}[c]{0.62\linewidth}
      {\popxbold\fontsize{12}{13}\selectfont\color{popcream}Signé\ \textbullet\ {\poplogo OnBuch\color{poporange}+}}\par\vspace{2pt}
      {\raggedright\popsemi\fontsize{8}{10.5}\selectfont\color{popline}\MakeUppercase{Équipe pédagogique OnBuch+}\\\MakeUppercase{Conforme au programme officiel MINESEC}\par}
    \end{minipage}\hfill
    {\poplogo\fontsize{22}{22}\selectfont\color{popcream}OnBuch\color{poporange}+}
  \end{tcolorbox}%
}

% ============================================================
% Page de garde
% #1 titre · #2 matière · #3 niveau · #4 module · #5 n° leçon · #6 durée ·
% #7 description / objectifs (texte ou itemize)
% ============================================================
\newcommand{\pagegarde}[7]{%
  \thispagestyle{empty}
  \newgeometry{a4paper,margin=1.7cm,top=1.6cm,bottom=1.4cm}%
  \begin{tikzpicture}[remember picture,overlay]
    \fill[poporange!18] ([xshift=-3.2cm,yshift=-3.6cm]current page.north east) circle (4.6cm);
    \fill[poppurple!14] ([xshift=2.4cm,yshift=7.6cm]current page.south west) circle (3.8cm);
  \end{tikzpicture}%
  {\poplogo\fontsize{30}{30}\selectfont\color{popink}OnBuch\color{poporange}+}\hfill
  \raisebox{0.6ex}{\poppill{Cours signé \textbullet\ Premium}{popink}{popcream}}\par
  \vspace{1pt}\popsquiggle{poppurple}\par
  \vspace{8pt}
  \begin{tcolorbox}[enhanced,colback=poporange,colframe=popink,boxrule=2.8pt,arc=13pt,
      drop shadow=popink,left=18pt,right=18pt,top=12pt,bottom=13pt,after skip=10pt]
    \poppill{#2 \textbullet\ #3}{popink}{popcream}\hspace{5pt}\poppill{#4}{white}{popink}\par\vspace{8pt}
    {\popdisplay\fontsize{14}{14}\selectfont\color{popink}LEÇON #5}\par\vspace{4pt}
    {\raggedright\hyphenpenalty=10000\exhyphenpenalty=10000\popdisplay\fontsize{25}{27}\selectfont\color{popcream}\MakeUppercase{#1}\par}
    \vspace{8pt}
    {\popxbold\fontsize{9.5}{9.5}\selectfont\color{popink}\MakeUppercase{Durée conseillée : #6}}
  \end{tcolorbox}
  \begin{tcolorbox}[enhanced,colback=white,colframe=popink,boxrule=2.2pt,arc=10pt,
      drop shadow=popink,left=16pt,right=16pt,top=10pt,bottom=10pt,after skip=8pt]
    \poppill{Dans cette leçon}{poppurple}{white}\par\vspace{5pt}
    \popsemi\fontsize{9.3}{12.6}\selectfont\color{popink}#7
  \end{tcolorbox}
  \noindent\begin{minipage}[t]{0.487\linewidth}
    \begin{tcolorbox}[enhanced,colback=popgreen,colframe=popink,boxrule=2.2pt,arc=10pt,
        drop shadow=popink,left=11pt,right=11pt,top=10pt,bottom=10pt,height=3.1cm,valign=center]
      \tcbox[colback=white,colframe=popink,boxrule=1.3pt,arc=5pt,boxsep=0pt,nobeforeafter,
        left=3pt,right=3pt,top=3pt,bottom=3pt,valign=center]{\qrcode[height=1.9cm]{https://play.google.com/store/apps/details?id=com.luvvix.onbuch&hl=fr}}%
      \hspace{8pt}\begin{minipage}[c]{0.5\linewidth}
        {\popdisplay\fontsize{11}{12.5}\selectfont\color{white}\MakeUppercase{Télécharge OnBuch}}\par\vspace{4pt}
        {\raggedright\popsemi\fontsize{8.4}{11}\selectfont\color{white}Gratuit \textbullet\ Google Play\\Tuteur IA, annales, concours, résultats\par}
      \end{minipage}
    \end{tcolorbox}
  \end{minipage}\hfill
  \begin{minipage}[t]{0.487\linewidth}
    \begin{tcolorbox}[enhanced,colback=poppurple,colframe=popink,boxrule=2.2pt,arc=10pt,
        drop shadow=popink,left=11pt,right=11pt,top=10pt,bottom=10pt,height=3.1cm,valign=center]
      \tikz[baseline=-0.7ex]\node[circle,fill=white,draw=popink,line width=1.3pt,minimum size=1.6cm,inner sep=0]{\popdisplay\fontsize{24}{24}\selectfont\color{poppurple}+};%
      \hspace{8pt}\begin{minipage}[c]{0.6\linewidth}
        {\popdisplay\fontsize{11}{12.5}\selectfont\color{white}\MakeUppercase{Passe à OnBuch+}}\par\vspace{4pt}
        {\raggedright\popsemi\fontsize{8.4}{11}\selectfont\color{white}Tous les cours signés, Tuteur IA illimité, fiches PDF — onglet OnBuch+ de l'app\par}
      \end{minipage}
    \end{tcolorbox}
  \end{minipage}
  \vspace{8pt}
  \popcontenu
  \vfill
  \signatureonbuch
  \restoregeometry
  \newpage
}

% Rangée « Ce cours contient » (page de garde)
\newcommand{\poptile}[3]{% #1 couleur #2 gros texte #3 légende
  \begin{tcolorbox}[enhanced,colback=#1,colframe=popink,boxrule=2pt,arc=9pt,shadow={2.5pt}{-2.5pt}{0pt}{popink},
      left=6pt,right=6pt,top=9pt,bottom=9pt,halign=center,valign=center,height=1.75cm,width=\linewidth,nobeforeafter]
    {\popdisplay\fontsize{12.5}{13}\selectfont\color{white}\MakeUppercase{#2}}\par\vspace{3pt}
    {\centering\popsemi\fontsize{7.8}{9.5}\selectfont\color{white}#3\par}
  \end{tcolorbox}}
\newcommand{\popcontenu}{%
  \par\noindent{\popxboldls\fontsize{7.5}{7.5}\selectfont\color{popmuted}CE COURS SIGNÉ CONTIENT}\par\vspace{7pt}
  \noindent\begin{minipage}[t]{0.235\linewidth}\poptile{popblue}{Cours}{complet, illustré, pas à pas}\end{minipage}\hfill
  \begin{minipage}[t]{0.235\linewidth}\poptile{poporange}{Méthodes}{et exercices résolus}\end{minipage}\hfill
  \begin{minipage}[t]{0.235\linewidth}\poptile{poppink}{Exercices}{type examen + corrigés}\end{minipage}\hfill
  \begin{minipage}[t]{0.235\linewidth}\poptile{popgreen}{Fiche bilan}{l'essentiel à retenir}\end{minipage}\par}

% Sommaire
\newtcolorbox{sommaire}{enhanced,colback=white,colframe=popink,boxrule=2.2pt,arc=10pt,
  drop shadow=popink,left=18pt,right=18pt,top=13pt,bottom=13pt,before skip=6pt,after skip=20pt,
  before upper={\poppill{Sommaire}{poppurple}{white}\par\vspace{9pt}\popsemi\fontsize{10.6}{14.5}\selectfont\color{popink}}}

% Signature de fin de cours — poussée tout en bas de la dernière page
\newcommand{\findecours}{%
  \par\vfill\nopagebreak
  \begin{center}\popsquiggle{poporange}\end{center}\vspace{4pt}
  \signatureonbuch
}
\AtBeginDocument{\def\llbracket{\mathopen{[\mkern-3.5mu[}}\def\rrbracket{\mathclose{]\mkern-3.5mu]}}} % intervalles d'entiers (glyphe ⟦ absent de la police)
% Glyphes absents de Fira Math : redéfinis à partir de glyphes disponibles
\AtBeginDocument{%
  \def\setminus{\mathbin{\backslash}}\let\smallsetminus\setminus
  \def\vdots{\mathord{\vbox{\baselineskip3.2pt\lineskiplimit0pt\kern4pt\hbox{.}\hbox{.}\hbox{.}}}}%
  \def\ddots{\mathinner{\mkern1mu\raise6.5pt\hbox{.}\mkern2mu\raise3.6pt\hbox{.}\mkern2mu\raise0.7pt\hbox{.}\mkern1mu}}%
  \def\triangle{\mathord{\scalebox{0.9}{$\symup{\Delta}$}}}%
  \def\nabla{\mathord{\raisebox{\depth}{\scalebox{1}[-1]{$\symup{\Delta}$}}}}%
  \def\checkmark{\popcheck{popdark}}%
  \def\star{\mathbin{\ast}}\def\bigstar{\mathord{\ast}}%
  \def\diamond{\mathbin{\scalebox{0.6}{$\Diamond$}}}%
  \def\bigcup{\mathop{\mathchoice{\raisebox{-0.3em}{\scalebox{1.7}{$\cup$}}}{\scalebox{1.25}{$\cup$}}{\cup}{\cup}}\displaylimits}%
  \def\bigcap{\mathop{\mathchoice{\raisebox{-0.3em}{\scalebox{1.7}{$\cap$}}}{\scalebox{1.25}{$\cap$}}{\cap}{\cap}}\displaylimits}%
  \def\lesseqqgtr{\mathrel{\lessgtr}}\def\bigoplus{\mathop{\oplus}\displaylimits}%
}
\providecommand{\ssi}{si et seulement si} % abréviation fréquente
\AtBeginDocument{\providecommand{\wideparen}{\overparen}} % arc de cercle

%% FIGFIX-BEGIN v5 — correctifs globaux des figures (docs/onbuch-generation/tools/figfix)
\usepackage{adjustbox}\usepackage{etoolbox}\tcbuselibrary{raster}
% toute figure TikZ/pgfplots plus large que la ligne est réduite à la largeur disponible
\BeforeBeginEnvironment{tikzpicture}{\begin{adjustbox}{max width=\linewidth}}
\AfterEndEnvironment{tikzpicture}{\end{adjustbox}}
% légendes pgfplots lisibles par-dessus les courbes
\pgfplotsset{every axis/.append style={legend style={fill=white,fill opacity=0.92,text opacity=1,draw=black!15,inner sep=3pt}}}
% étiquettes d'arêtes (syntaxe edge["..."]) avec fond blanc
\tikzset{every edge quotes/.append style={fill=white,inner sep=1.2pt}}
%% FIGFIX-END
\begin{document}

\pagegarde{Amélioration de la santé reproductive}{SVT}{1ère A}{Module 2 : L'Éducation à la Santé}{NA06}{3 h}{Cette leçon présente les notions de base de la reproduction humaine et de la puberté, puis analyse les causes et conséquences des grossesses précoces et non désirées. Elle expose enfin les infections sexuellement transmissibles, leurs moyens de prévention et les notions générales de planification familiale, afin de conduire l'élève vers un comportement responsable en matière de santé reproductive.
\vspace{4pt}
\begin{multicols}{2}\small
\begin{itemize}
\item À la fin de cette leçon, je sais définir la puberté et décrire les transformations qu'elle entraîne chez le garçon et chez la fille
\item À la fin de cette leçon, je sais décrire les grandes étapes de la reproduction humaine (ovulation, fécondation, nidation)
\item À la fin de cette leçon, je sais expliquer les causes des grossesses précoces et non désirées chez les adolescents
\item À la fin de cette leçon, je sais expliquer les risques liés aux grossesses précoces pour la mère, l'enfant et la société
\item À la fin de cette leçon, je sais citer les principales infections sexuellement transmissibles et leurs modes de transmission
\item À la fin de cette leçon, je sais justifier l'utilisation des moyens de prévention des IST (abstinence, fidélité, préservatif, dépistage)
\end{itemize}\end{multicols}}

\begin{sommaire}
\begin{multicols}{2}
\begin{enumerate}
\item Puberté et bases de la reproduction humaine
\item Grossesses précoces et non désirées : causes
\item Grossesses précoces : conséquences
\item Infections sexuellement transmissibles (IST)
\item Prévention des IST et comportement responsable
\item Planification familiale (notions générales)
\item Activité d'intégration
\item Méthodes et exercices résolus
\item Exercices
\item Corrigés des exercices
\item Fiche bilan
\end{enumerate}
\end{multicols}
\end{sommaire}

\begin{objectifs}
\begin{itemize}
\item À la fin de cette leçon, je sais définir la puberté et décrire les transformations qu'elle entraîne chez le garçon et chez la fille
\item À la fin de cette leçon, je sais décrire les grandes étapes de la reproduction humaine (ovulation, fécondation, nidation)
\item À la fin de cette leçon, je sais expliquer les causes des grossesses précoces et non désirées chez les adolescents
\item À la fin de cette leçon, je sais expliquer les risques liés aux grossesses précoces pour la mère, l'enfant et la société
\item À la fin de cette leçon, je sais citer les principales infections sexuellement transmissibles et leurs modes de transmission
\item À la fin de cette leçon, je sais justifier l'utilisation des moyens de prévention des IST (abstinence, fidélité, préservatif, dépistage)
\item À la fin de cette leçon, je sais présenter les notions générales de planification familiale et citer quelques méthodes contraceptives
\item À la fin de cette leçon, je sais adopter et argumenter un comportement responsable en matière de santé reproductive
\end{itemize}
\end{objectifs}

\begin{prerequis}
\begin{itemize}
\item Notions sur l'appareil reproducteur humain vues en classe de 3e (organes, gamètes)
\item Notion de fécondation et de développement embryonnaire (niveau collège)
\item Notions générales sur les microbes et les maladies infectieuses (hygiène, vaccination)
\item Lecture et interprétation de graphiques et de tableaux statistiques
\item Notion de cycle menstruel abordée en début d'année ou en 3e
\end{itemize}
\end{prerequis}

\newpage

\coursec{Puberté et bases de la reproduction humaine}

Dans un quartier de Yaoundé, Amina, 16 ans, élève en classe de Seconde, découvre qu'elle est enceinte et doit interrompre sa scolarité. Son amie Clarisse, elle, a appris au centre de santé du quartier comment se protéger des infections sexuellement transmissibles. Pourquoi certains adolescents sont-ils mieux armés que d'autres face aux risques de la vie sexuelle ? Quelles connaissances scientifiques permettent de faire des choix responsables ? C'est à ces questions que cette leçon va répondre.

\courssub{La puberté : une transition majeure}

\begin{definition}[Puberté]
La \cle{puberté} est la période de transition entre l'enfance et l'âge adulte, marquée par l'acquisition de la capacité de se reproduire. Elle se traduit par une maturation des organes reproducteurs (caractères sexuels primaires) et l'apparition de caractères sexuels secondaires sous l'action des hormones sexuelles.
\end{definition}

\begin{definition}[Caractères sexuels primaires et secondaires]
Les \cle{caractères sexuels primaires} sont les organes reproducteurs présents dès la naissance (testicules, pénis chez le garçon ; ovaires, utérus, vagin chez la fille). Les \cle{caractères sexuels secondaires} sont les traits morphologiques et physiologiques qui apparaissent à la puberté et distinguent les sexes (pilosité, voix, seins, bassin, musculature).
\end{definition}

\begin{savaistu}[Puberté au Cameroun]
Au Cameroun, l'âge moyen des premières règles (ménarche) se situe autour de 13 ans, mais la puberté peut survenir plus tôt (dès 9-10 ans) ou plus tard (jusqu'à 15-16 ans) selon l'alimentation, l'hérédité, l'état nutritionnel et l'environnement. Une étude menée dans les lycées de Douala et Yaoundé montre que 15 \% des filles ont leurs premières règles avant 12 ans.
\end{savaistu}

\begin{propriete}[Âge moyen de la puberté]
\begin{itemize}
    \item \textbf{Fille :} 10 à 14 ans (ménarche vers 12-13 ans en moyenne).
    \item \textbf{Garçon :} 12 à 16 ans (premières éjaculations vers 13-14 ans en moyenne).
\end{itemize}
La variabilité individuelle est grande : la puberté précoce (avant 8 ans chez la fille, 9 ans chez le garçon) ou tardive (absence de signes à 13 ans fille, 14 ans garçon) justifie une consultation médicale.
\end{propriete}

\begin{popfigure}
\begin{tikzpicture}[scale=0.9, every node/.style={font=\small}]
\draw[->] (0,0) -- (14,0) node[right] {Âge (ans)};
\draw[->] (0,0) -- (0,8) node[above] {Fréquence cumulative (\%)};
\draw[poppink, thick, smooth] plot coordinates {
    (8,0) (9,2) (10,10) (11,30) (12,60) (12.5,80) (13,90) (14,97) (15,100)
};
\draw[popblue, thick, smooth, dashed] plot coordinates {
    (9,0) (10,1) (11,5) (12,15) (13,40) (14,70) (15,90) (16,98) (17,100)
};
\node[poppink] at (5,7) {\textbf{Filles}};
\node[popblue] at (11,3) {\textbf{Garçons}};
\draw[poppink, thin, dotted] (12.5,0) -- (12.5,80);
\draw[popblue, thin, dotted] (14,0) -- (14,70);
\node[poppink, anchor=north] at (12.5,0) {Médiane $\approx$ 12,5 ans};
\node[popblue, anchor=north] at (14,0) {Médiane $\approx$ 14 ans};
\end{tikzpicture}
\legende{Âge d'apparition de la puberté (ménarche chez la fille, première éjaculation chez le garçon) : distribution cumulative dans une population camerounaise type. La variabilité individuelle est normale.}
\end{popfigure}

\begin{attention}[Piège fréquent]
Ne pas confondre \emph{âge chronologique} et \emph{âge pubertaire}. Deux élèves de même classe (ex. 14 ans) peuvent être à des stades pubertaires très différents. L'absence de règles à 14 ans chez une fille ou de mue de la voix à 15 ans chez un garçon n'est pas systématiquement pathologique, mais mérite un avis médical.
\end{attention}

\courssub{Les transformations hormonales et physiques}

\begin{methode}[Démarche d'investigation : rôle des hormones à la puberté]
\begin{enumerate}
    \item \textbf{Observation :} À la puberté, les testicules augmentent de volume chez le garçon ; les seins se développent et les règles apparaissent chez la fille.
    \item \textbf{Hypothèse :} Des substances chimiques (hormones) sécrétées par les gonades (testicules, ovaires) pilotent ces transformations.
    \item \textbf{Expérience / Données :} Dosages sanguins montrent une élévation de la testostérone chez le garçon (de $<0,5$ à $10-30$ nmol/L) et des œstrogènes/progestérone chez la fille (cycles mensuels). L'ablation des gonades avant la puberté (castration expérimentale animale) bloque les caractères secondaires ; l'injection d'hormones les restaure.
    \item \textbf{Interprétation :} Les gonades, stimulées par l'hypophyse (LH, FSH), sécrètent les hormones sexuelles qui agissent sur les organes cibles.
    \item \textbf{Conclusion :} La testostérone masculinise (pilosité, voix, musculature, spermatogenèse) ; les œstrogènes féminisent (seins, bassin, cycle menstruel, ovulation) ; la progestérone prépare l'utérus à la grossesse.
\end{enumerate}
\end{methode}

\begin{propriete}[Hormones clés de la puberté]
\begin{center}
\begin{tabularx}{\linewidth}{|l|X|X|}\hline
\textbf{Hormone} & \textbf{Source principale} & \textbf{Rôles principaux} \\ \hline
Testostérone & Testicules (cellules de Leydig) & Croissance testicules/pénis, spermatogenèse, pilosité, mue, masse musculaire, libido \\ \hline
Œstrogènes (estradiol) & Ovaires (follicules) & Croissance utérus/trompes, seins, bassin, cycle menstruel, ostéogenèse, métabolisme lipidique \\ \hline
Progestérone & Corps jaune (ovaire) & Préparation endomètre (nidation), maintien grossesse, thermogenèse \\ \hline
\end{tabularx}
\end{center}
\end{propriete}

\begin{exemplebox}[Transformations visibles]
\begin{itemize}
    \item \textbf{Chez la fille :} Bouton mammaire (thélarchie) $\rightarrow$ pilosité pubienne (pubarchie) $\rightarrow$ poussée de croissance $\rightarrow$ ménarche (premières règles) $\rightarrow$ élargissement du bassin.
    \item \textbf{Chez le garçon :} Augmentation volume testiculaire $\rightarrow$ pilosité pubienne $\rightarrow$ croissance pénis $\rightarrow$ mue (voix) $\rightarrow$ pilosité faciale/axillaire $\rightarrow$ poussée de croissance tardive $\rightarrow$ premières éjaculations (spermarche).
\end{itemize}
\end{exemplebox}

\begin{popfigure}
\begin{tikzpicture}[scale=0.85, every node/.style={font=\small, align=center}]
\foreach \i/\label/\val/\col in {
    1/Thélarchie/10.5/poppink,
    2/Pubarchie/11.5/poppink,
    3/Poussée croissance/12/poppink,
    4/Ménarche/12.8/poppink,
    5/Élargissement bassin/13.5/poppink} {
    \fill[\col!70] (0.2,\i*1.2-0.4) rectangle (\val*0.35,\i*1.2+0.4);
    \node[left] at (0,\i*1.2) {\label};
    \node[right] at (\val*0.35+0.3,\i*1.2) {\val~ans};
}
\foreach \i/\label/\val/\col in {
    7/Testicules/11.5/popblue,
    8/Pubarchie/12.5/popblue,
    9/Pénis/13/popblue,
    10/Mue/13.8/popblue,
    11/Pilosité faciale/14.5/popblue,
    12/Poussée croissance/14/popblue,
    13/Spermarche/13.5/popblue} {
    \fill[\col!70] (6,\i*1.2-0.4) rectangle (6+\val*0.35,\i*1.2+0.4);
    \node[left] at (5.8,\i*1.2) {\label};
    \node[right] at (6+\val*0.35+0.3,\i*1.2) {\val~ans};
}
\node[poppink, font=\bfseries] at (2.5,7.5) {FILLE};
\node[popblue, font=\bfseries] at (9.5,7.5) {GARÇON};
\draw[<->] (0,-0.5) -- (11,-0.5) node[midway,below] {Âge (ans)};
\foreach \x in {8,9,10,11,12,13,14,15,16} {
    \draw (\x*0.35,-0.6) -- (\x*0.35,-0.4) node[below] {\x};
    \draw (6+\x*0.35,-0.6) -- (6+\x*0.35,-0.4) node[below] {\x};
}
\end{tikzpicture}
\legende{Âges moyens d'apparition des caractères sexuels secondaires chez la fille (rose) et le garçon (bleu) au Cameroun. Les barres indiquent l'âge médian ; l'intervalle interindividuel est d'environ $\pm 1,5$ an.}
\end{popfigure}

\courssub{Appareils reproducteurs et gamétogenèse}

\begin{popfigure}
\begin{tikzpicture}[scale=0.75, every node/.style={font=\small, align=center}]
% --- APPAREIL MASCULIN (gauche) ---
\begin{scope}[xshift=0cm]
\fill[popblueL] (0,5) ellipse (1.2 and 0.4);
\node at (0,5) {Vessie};
\fill[popblue!50] (-1.5,4) ellipse (0.6 and 0.8);
\fill[popblue!50] (1.5,4) ellipse (0.6 and 0.8);
\node[align=center] at (-1.5,4){Vésicules\\ séminales};
\node[align=center] at (1.5,4){Vésicules\\ séminales};
\fill[popblue!70] (0,2.8) ellipse (1 and 0.5);
\node at (0,2.8) {Prostate};
\draw[popblue, thick] (0,2.3) -- (0,0.5);
\draw[popblue, thick] (-0.5,0.5) -- (-0.5,-1) -- (0.5,-1) -- (0.5,0.5) -- cycle;
\node at (0,-0.2) {Pénis};
\draw[popblue, thick] (-1.5,3.2) .. controls (-1,3) and (-0.5,2.5) .. (0,2.8);
\draw[popblue, thick] (1.5,3.2) .. controls (1,3) and (0.5,2.5) .. (0,2.8);
\fill[popblue!30] (-2.2,2.5) ellipse (0.4 and 0.7);
\fill[popblue!30] (2.2,2.5) ellipse (0.4 and 0.7);
\node at (-2.2,2.5) {Épididyme};
\node at (2.2,2.5) {Épididyme};
\fill[popblue!20] (-2.2,1) ellipse (0.6 and 1);
\fill[popblue!20] (2.2,1) ellipse (0.6 and 1);
\node at (-2.2,1) {Testicule};
\node at (2.2,1) {Testicule};
\draw[->, poporange, thick] (-2.2,0) -- (-2.2,-0.8) node[below] {Spermatogenèse};
\draw[->, poporange, thick] (2.2,0) -- (2.2,-0.8);
\node[popblue, font=\bfseries] at (0,6) {APP. REPRODUCTEUR MASCULIN};
\end{scope}

% --- APPAREIL FÉMININ (droite) ---
\begin{scope}[xshift=10cm]
\fill[poppink!30] (-2.5,3.5) ellipse (0.5 and 1);
\fill[poppink!30] (2.5,3.5) ellipse (0.5 and 1);
\node at (-2.5,3.5) {Ovaire};
\node at (2.5,3.5) {Ovaire};
\draw[poppink, thick] (-2.5,4) .. controls (-1.5,4.5) and (0,4.5) .. (0,3.5);
\draw[poppink, thick] (2.5,4) .. controls (1.5,4.5) and (0,4.5) .. (0,3.5);
\node at (-3.5,4) {Pavillon};
\node at (3.5,4) {Pavillon};
\fill[poppink!50] (0,1.5) ellipse (1.5 and 2);
\node at (0,1.5) {Utérus};
\fill[poppink!70] (0,-0.2) ellipse (0.8 and 0.6);
\node at (0,-0.2) {Col};
\draw[poppink, thick] (-0.8,-0.8) -- (-0.8,-2.5) -- (0.8,-2.5) -- (0.8,-0.8);
\node at (0,-1.6) {Vagin};
\draw[->, poporange, thick] (-2.5,4.5) -- (-2.5,5.2) node[above] {Ovulation};
\draw[->, poporange, thick] (2.5,4.5) -- (2.5,5.2);
\node[poppink, font=\bfseries] at (0,6) {APP. REPRODUCTEUR FÉMININ};
\end{scope}
\end{tikzpicture}
\legende{Coupes sagittales simplifiées des appareils reproducteurs masculin (gauche) et féminin (droite). En orange : sites de production des gamètes (spermatogenèse dans les testicules, ovulation dans les ovaires) et trajet des gamètes.}
\end{popfigure}

\begin{definition}[Gamétogenèse]
La \cle{gamétogenèse} est la formation des gamètes (cellules reproductrices haploïdes, $n=23$ chromosomes) à partir de cellules souches diploïdes ($2n=46$) par méiose.
\begin{itemize}
    \item \textbf{Spermatogenèse :} Dans les tubules séminifères des testicules, continue dès la puberté, produit $\approx 2 \times 10^8$ spermatozoïdes/jour (total deux testicules). Un spermatozoïde = flagelle (mobilité) + noyau haploïde + acrosome (enzymes pour pénétrer l'ovule).
    \item \textbf{Oogenèse :} Dans les ovaires, commence \emph{in utero} (stock d'ovocytes I bloqués en prophase I), reprend à la puberté par cycles : un ovocyte II par cycle (ovulation), achève la méiose II seulement si fécondation.
\end{itemize}
\end{definition}

\begin{exoresolu}[Durée et rythme de la spermatogenèse]
Énoncé : La spermatogenèse dure environ 74 jours. La production testiculaire est d'environ $2 \times 10^8$ spermatozoïdes par jour (pour les deux testicules). Combien de spermatozoïdes sont produits par seconde en moyenne ?
\tcblower
Solution :
\[
\text{Production journalière} = 2 \times 10^8 \text{ spermatozoïdes/jour}
\]
\[
\text{Nombre de secondes dans une journée} = 24 \times 3600 = 86\,400 \text{ s}
\]
\[
\text{Production par seconde} = \frac{2 \times 10^8}{8,64 \times 10^4} \approx 2\,315 \text{ spermatozoïdes/s}
\]
\textbf{Conclusion :} La production est continue et massive (plus de 2\,000 spermatozoïdes par seconde), contrairement à l'ovulation unique par cycle chez la femme. Cette différence fondamentale explique la fertilité permanente de l'homme après la puberté.
\end{exoresolu}

\courssub{Du cycle menstruel à la fécondation}

\begin{definition}[Cycle menstruel et ovulation]
Le \cle{cycle menstruel} (durée moyenne 28 jours, variable 21-35 jours) prépare l'utérus à une éventuelle grossesse.
\begin{itemize}
    \item \textbf{Phase folliculaire (J1-J14, variable) :} Sous l'effet de la FSH, un follicule mûrit et sécrète des œstrogènes $\rightarrow$ épaississement de l'endomètre (phase de prolifération).
    \item \textbf{Ovulation (vers J14) :} Pic de LH $\rightarrow$ rupture du follicule $\rightarrow$ libération de l'ovocyte II dans le pavillon de la trompe.
    \item \textbf{Phase lutéale (J15-J28, fixe $\approx$ 14 jours) :} Le follicule devient corps jaune $\rightarrow$ sécrète progestérone + œstrogènes $\rightarrow$ endomètre sécréteur (prêt pour nidation). Sans fécondation, corps jaune régresse $\rightarrow$ chute hormones $\rightarrow$ règles (destruction endomètre).
\end{itemize}
\end{definition}

\begin{exemplebox}[Repères chiffrés du cycle]
\begin{itemize}
    \item Durée moyenne : 28 jours (J1 = premier jour des règles).
    \item Ovulation : vers J14 (période fertile $\approx$ J11-J16, durée de vie ovocyte 12-24 h, spermatozoïdes 3-5 jours dans trompes).
    \item Règles : 3 à 7 jours, perte de 30-80 mL de sang.
    \item Ménarche : premières règles (signe que l'ovulation a commencé, mais cycles souvent anovulatoires les 1-2 premières années).
\end{itemize}
\end{exemplebox}

\begin{popfigure}
\begin{tikzpicture}[scale=0.85, every node/.style={font=\small}]
\node[ellipse, draw=poporange, thick, fill=poporange!10, minimum width=2.5cm, minimum height=1.2cm, align=center] (ovaire) at (0,0) {Ovaire\\ (Ovulation J14)};
\draw[poppink, very thick, ->] (1.5,0) .. controls (2.5,0.5) and (3.5,0.5) .. (4.5,0) node[midway, above, sloped] {Trompe de Fallope};
\node[ellipse, draw=poppurple, thick, fill=poppurple!10, minimum width=2.5cm, minimum height=1.2cm, align=center] (fecond) at (5.5,0) {Fécondation\\ dans la trompe};
\draw[poppink, very thick, ->] (7,0) -- (9,0) node[midway, above] {Migration 3-4 jours};
\node[ellipse, draw=popgreen, thick, fill=popgreen!10, minimum width=3cm, minimum height=1.5cm, align=center] (uterus) at (11,0) {Utérus\\ Nidation J6-J7};
\draw[->, poporange] (ovaire.north) -- ++(0,1) node[above] {Libération ovocyte II};
\draw[->, poppurple] (fecond.north) -- ++(0,1) node[above] {Rencontre spermatozoïde};
\draw[->, popgreen] (uterus.north) -- ++(0,1) node[above] {Implantation blastocyste};
\node[align=center] at (0,-1.8) {\textbf{1. Ovulation}};
\node[align=center] at (5.5,-1.8) {\textbf{2. Fécondation}};
\node[align=center] at (11,-1.8) {\textbf{3. Nidation}};
\end{tikzpicture}
\legende{Trajet de l'ovule : ovulation dans le pavillon de la trompe $\rightarrow$ fécondation dans l'ampoule de la trompe (rencontre spermatozoïde/ovocyte) $\rightarrow$ migration du zygote (puis blastocyste) vers l'utérus $\rightarrow$ nidation dans l'endomètre vers J6-J7 après fécondation.}
\end{popfigure}

\begin{definition}[Fécondation et nidation]
La \cle{fécondation} est la fusion d'un spermatozoïde et d'un ovocyte II dans la trompe (ampoule), restaurant la diploidie ($2n=46$) : formation du zygote. La \cle{nidation} (implantation) est la fixation du blastocyste (stade 100-200 cellules, J5-J6) dans l'endomètre utérin, marquant le début de la grossesse.
\end{definition}

\begin{attention}[Erreur fréquente — « Première fois = pas de risque »]
\textbf{FAUX.} Dès la première ovulation (qui précède souvent la première règles de quelques semaines) et dès le premier rapport sexuel non protégé, une grossesse est possible. De même, les IST se transmettent dès le premier rapport. La « période de sécurité » n'existe pas : la variabilité de l'ovulation (stress, fatigue, maladie) rend toute méthode de calendrier (Ogino) peu fiable chez l'adolescente.
\end{attention}

\begin{propriete}[Conséquence pratique immédiate]
Dès la puberté (ménarche chez la fille, spermarche chez le garçon), \textbf{tout rapport sexuel vaginal non protégé peut entraîner une grossesse ou une IST}. La fertilité est maximale chez l'adolescent(e) en bonne santé. La responsabilité reproductive commence avec la connaissance de ces mécanismes.
\end{propriete}

\begin{exercice}[Bilan de la section]
\begin{enumerate}
    \item Citez deux caractères sexuels secondaires masculins et deux féminins.
    \item Où et quand a lieu la fécondation humaine ? Quelle est la formule chromosomique du zygote ?
    \item Expliquez pourquoi une fille peut tomber enceinte avant ses premières règles.
    \item Complétez : La testostérone est sécrétée par les cellules de \_\_\_\_\_\_ ; les œstrogènes par les \_\_\_\_\_\_ ovariens ; la progestérone par le \_\_\_\_\_\_.
\end{enumerate}
\end{exercice}

\begin{corrige}[Exercice n°1]
\begin{enumerate}
    \item Masculins : pilosité faciale, mue de la voix, développement musculaire, augmentation testicules/pénis. Féminins : développement des seins, élargissement du bassin, pilosité pubienne, ménarche.
    \item Dans la trompe de Fallope (ampoule), vers le 14e jour du cycle (ovulation). Zygote : $2n = 46$ chromosomes (23 paternels + 23 maternels).
    \item La première ovulation survient \emph{avant} les premières règles (ménarche). Si un rapport a lieu lors de ce premier cycle ovulatoire, la fécondation est possible, mais il n'y a pas encore eu de règles pour signaler la fertilité.
    \item Cellules de Leydig ; follicules ; corps jaune.
\end{enumerate}
\end{corrige}

\coursec{Grossesses précoces et non désirées : causes}

Dans la section précédente, tu as découvert comment la puberté rend l'organisme apte à la reproduction : chez la jeune fille, les premières règles (\cle{ménarche}) annoncent le début de l'activité ovarienne cyclique ; chez le garçon, la \cle{spermatogenèse} devient fonctionnelle. Tu as aussi compris qu'un rapport sexuel non protégé, même unique, peut aboutir à une \cle{fécondation} si celui-ci a lieu durant la \cle{période fertile} du cycle. Cette capacité biologique nouvelle, acquise parfois dès 12 ou 13 ans, n'est malheureusement pas toujours accompagnée de la maturité psychologique, de l'information et de l'autonomie nécessaires pour la gérer. C'est là que naît un problème majeur de santé publique au Cameroun : les \cle{grossesses précoces}. Chaque année, des milliers d'élèves de collège et de lycée interrompent leur scolarité à cause d'une maternité non prévue. Comprendre \emph{pourquoi} ces grossesses surviennent est la première étape pour les prévenir.

\courssub{Définitions : grossesse précoce et grossesse non désirée}

\begin{definition}[Grossesse précoce et grossesse non désirée]
Une \cle{grossesse précoce} (ou grossesse chez l'adolescente) est une grossesse survenant chez une jeune fille âgée de moins de 19 ans, c'est-à-dire avant la fin de l'adolescence définie par l'Organisation mondiale de la santé (10 à 19 ans). Elle est le plus souvent \emph{non désirée}.\par\medskip
Une \cle{grossesse non désirée} est une grossesse qui survient alors que la femme ou le couple ne l'avait pas planifiée ni souhaitée à ce moment de sa vie, quel que soit son âge.
\end{definition}

\begin{attention}[Deux notions qui se recoupent sans être synonymes]
Ne confonds pas « grossesse précoce » et « grossesse non désirée » :
\begin{itemize}
\item une grossesse \emph{précoce} est définie par l'\textbf{âge} (avant 19 ans) ;
\item une grossesse \emph{non désirée} est définie par l'\textbf{intention} (absence de projet de maternité).
\end{itemize}
Une femme de 32 ans peut connaître une grossesse non désirée (elle n'est pas précoce) ; une adolescente mariée de 18 ans peut désirer sa grossesse (elle reste précoce, avec tous ses risques médicaux). En pratique, les deux situations se recoupent très souvent : la grande majorité des grossesses d'adolescentes sont non désirées.
\end{attention}

\courssub{Démarche d'investigation : d'où viennent les grossesses précoces ?}

\begin{experience}[Enquête de compréhension dans un établissement scolaire]
\textbf{Observation.} Dans un lycée de la région du Centre, l'infirmier scolaire constate que 12 élèves ont abandonné leurs études en trois ans pour cause de grossesse.\par
\textbf{Hypothèse.} « Les grossesses précoces sont dues uniquement à l'imprudence des adolescentes. »\par
\textbf{Données recueillies.} Un questionnaire anonyme et des entretiens (menés avec l'accord des autorités et dans le respect de la confidentialité) révèlent que : certaines élèves ignoraient le fonctionnement du cycle menstruel et la période fertile ; d'autres subissaient la pression d'un petit ami plus âgé ; d'autres encore venaient de familles très pauvres où un « amant » finançait les fournitures scolaires ; une élève a été victime d'abus.\par
\textbf{Interprétation.} L'hypothèse initiale est fausse : aucune cause unique n'explique les 12 cas. Chaque grossesse résulte d'une \emph{combinaison} de facteurs individuels, familiaux, sociaux et économiques.\par
\textbf{Conclusion.} Les grossesses précoces sont un phénomène \cle{multifactoriel}. Toute prévention efficace doit donc agir sur plusieurs causes à la fois.
\end{experience}

\courssub{Les grandes causes des grossesses précoces}

\textbf{Cause 1 : les rapports sexuels précoces et non protégés.}\par
L'âge des premiers rapports sexuels s'abaisse : de nombreuses enquêtes montrent qu'une partie des adolescents camerounais ont leur premier rapport avant 16 ans. Or, lors de ces premiers rapports, aucune méthode de prévention (préservatif, contraception) n'est utilisée, par ignorance, par gêne, ou parce que le rapport n'était pas prévu. Un seul rapport non protégé pendant la période fertile du cycle (autour de l'ovulation, vers le $14^{\text{ème}}$ jour d'un cycle de 28 jours) suffit à déclencher une grossesse, d'autant que les spermatozoïdes peuvent survivre jusqu'à 5 jours dans les voies génitales féminines. C'est la cause \emph{directe} et immédiate de la quasi-totalité des grossesses précoces.

\textbf{Cause 2 : la méconnaissance de la reproduction et des méthodes de prévention.}\par
Beaucoup d'adolescents ignorent les faits biologiques élémentaires : la date approximative de l'ovulation, la durée de vie des gamètes, l'existence et le mode d'emploi des contraceptifs. Des idées fausses circulent : « on ne peut pas tomber enceinte au premier rapport », « se laver après le rapport protège », « le retrait (\emph{coitus interruptus}) est sûr ». Toutes sont \emph{fausses}. Ce manque d'éducation sexuelle — à la maison comme parfois à l'école — laisse les adolescents sans défense face à leur propre fertilité.

\textbf{Cause 3 : les pressions sociales et économiques.}\par
La pauvreté joue un rôle majeur. Une adolescente dont la famille ne peut pas payer les frais de scolarité, les fournitures ou même la nourriture peut accepter une relation avec un partenaire plus âgé et plus fortuné (phénomène parfois appelé « \cle{sponsor} ») en échange d'un soutien matériel. Le rapport de force est alors inégal : elle peut difficilement exiger l'usage du préservatif. S'y ajoutent les \cle{mariages précoces} : dans certaines régions, des jeunes filles sont mariées dès 13 ou 14 ans, par tradition ou pour soulager la famille d'une charge financière ; la grossesse survient alors presque inévitablement, dans un corps qui n'a pas fini sa croissance (bassin étroit, utérus immature).

\textbf{Cause 4 : l'influence des pairs, des médias et des substances.}\par
La pression du groupe d'amis (« tout le monde l'a fait, sauf toi »), certains contenus médiatiques qui banalisent ou valorisent la sexualité précoce sans en montrer les conséquences, poussent des adolescents à des comportements qu'ils n'auraient pas choisis seuls. La consommation d'\cle{alcool} ou de drogues aggrave le risque : ces substances diminuent la vigilance, le jugement et la capacité à refuser un rapport ou à exiger une protection. Un rapport « sous l'emprise de l'alcool » est très souvent un rapport non protégé.

\textbf{Cause 5 : les violences sexuelles et les abus.}\par
Une partie des grossesses précoces résulte de rapports imposés : viol, abus sexuels commis par un proche, un voisin, un enseignant ou un employeur. La victime, souvent paralysée par la peur, la honte ou les menaces, ne dénonce pas les faits et ne reçoit aucune aide. Il faut le dire clairement : dans ce cas, la responsabilité incombe \emph{entièrement à l'agresseur}, jamais à la victime. Toute violence sexuelle doit être dénoncée auprès d'un adulte de confiance, d'un centre de santé ou des services de protection (au Cameroun, le numéro d'urgence \textbf{1517} permet de signaler les violences faites aux enfants).

\courssub{Une vision d'ensemble : l'arbre des causes}

\begin{popfigure}
\begin{tikzpicture}[
  grow=right,
  level 1/.style={sibling distance=32mm, level distance=42mm},
  level 2/.style={sibling distance=11mm, level distance=52mm},
  every node/.style={font=\small},
  racine/.style={draw=popdark, fill=popdark, text=white, rounded corners=2pt, align=center, inner sep=3pt},
  cat/.style={draw=popblue, fill=popblueL, rounded corners=2pt, align=center, inner sep=2pt},
  feuille/.style={draw=popmuted, fill=popcream, rounded corners=2pt, align=center, inner sep=2pt, font=\footnotesize},
  edge from parent/.style={draw=popline, thick}
]
\node[racine, align=center]{Grossesses\\précoces}
  child { node[cat, align=center]{Causes\\économiques}
    child { node[feuille]{Mariages précoces} }
    child { node[feuille, align=center]{Partenaires âgés\\(« sponsors »)} }
    child { node[feuille]{Pauvreté} }
  }
  child { node[cat, align=center]{Causes\\sociales}
    child { node[feuille]{Violences sexuelles} }
    child { node[feuille, align=center]{Pression des pairs\\et des médias} }
    child { node[feuille]{Alcool et drogues} }
  }
  child { node[cat, align=center]{Causes\\familiales}
    child { node[feuille, align=center]{Absence de dialogue\\parents--enfants} }
    child { node[feuille, align=center]{Traditions\\(mariage forcé)} }
  }
  child { node[cat, align=center]{Causes\\individuelles}
    child { node[feuille]{Rapports non protégés} }
    child { node[feuille, align=center]{Méconnaissance\\de la reproduction} }
  };
\end{tikzpicture}
\legende{Arbre des causes des grossesses précoces. Le phénomène (tronc, à gauche) résulte de quatre familles de causes — individuelles, familiales, sociales et économiques — qui se combinent le plus souvent.}
\end{popfigure}

\courssub{Que déclarent les adolescentes elles-mêmes ?}

Pour se représenter le poids relatif des différentes causes, on peut interroger directement les jeunes mères. Le diagramme circulaire ci-dessous présente une répartition \emph{indicative et fictive}, construite à but pédagogique, des principales causes citées par des adolescentes enceintes lors d'une enquête imaginaire.

\begin{popfigure}
\begin{tikzpicture}
\def\r{2.6}
% Rapports non protégés : 35% -> 126°
\fill[popblue] (0,0) -- (90:\r) arc (90:-36:\r) -- cycle;
% Méconnaissance : 25% -> 90°
\fill[poporange] (0,0) -- (-36:\r) arc (-36:-126:\r) -- cycle;
% Pauvreté/pressions éco : 20% -> 72°
\fill[popgreen] (0,0) -- (-126:\r) arc (-126:-198:\r) -- cycle;
% Pairs/alcool/médias : 12% -> 43.2°
\fill[poppurple] (0,0) -- (-198:\r) arc (-198:-241.2:\r) -- cycle;
% Violences : 8% -> 28.8°
\fill[poppink] (0,0) -- (-241.2:\r) arc (-241.2:-270:\r) -- cycle;
\draw[popdark, thick] (0,0) circle (\r);
% Étiquettes
\node[font=\footnotesize, text=white, align=center] at (27:1.5) {Rapports non\\protégés 35\,\%};
\node[font=\footnotesize, text=white, align=center] at (-81:1.5) {Méconnaissance\\25\,\%};
\node[font=\footnotesize, text=white, align=center] at (-162:1.5) {Pauvreté et\\pressions 20\,\%};
\node[font=\footnotesize, text=white, align=center] at (-220:1.6) {Pairs, alcool\\12\,\%};
\node[font=\footnotesize, text=white, align=center] at (-256:1.7) {Violences\\8\,\%};
\end{tikzpicture}
\legende{Répartition indicative (données fictives à but pédagogique) des causes citées par des adolescentes enceintes. Aucune cause ne domine totalement : les rapports non protégés et la méconnaissance totalisent à eux deux 60\,\% des cas, ce qui montre l'importance de l'information et de la prévention.}
\end{popfigure}

\begin{attention}[Erreur fréquente : la cause unique]
Au Probatoire comme dans la vie, évite d'attribuer les grossesses précoces à une seule cause (« c'est la faute des filles », « c'est seulement la pauvreté »). Une copie qui ne cite qu'un facteur est incomplète. La démarche correcte consiste à montrer que plusieurs facteurs \emph{se combinent} : par exemple, la pauvreté (cause économique) pousse vers un partenaire plus âgé (cause sociale), dans un contexte d'ignorance des méthodes de prévention (cause individuelle) — et la grossesse survient.
\end{attention}

\courssub{Méthode : analyser une situation concrète}

\begin{methode}[Analyser une situation de grossesse précoce en quatre étapes]
Face à un cas concret (exercice, exposé, situation de la vie réelle), procède ainsi :
\begin{enumerate}
\item \textbf{Décris les faits} : âge de la jeune fille, contexte familial, scolaire et économique, circonstances du rapport.
\item \textbf{Identifie les causes} en les classant : individuelle (ignorance, rapport non protégé), familiale (absence de dialogue, mariage forcé), sociale (pression des pairs, violence) ou économique (pauvreté, dépendance matérielle).
\item \textbf{Repère les combinaisons} : quelles causes se renforcent mutuellement ?
\item \textbf{Propose des solutions ciblées} : chaque cause identifiée appelle une réponse adaptée (information et éducation sexuelle contre la méconnaissance ; soutien scolaire et aide matérielle contre la pauvreté ; dénonciation et protection contre les violences).
\end{enumerate}
\end{methode}

\begin{exoresolu}[Analyse d'un cas concret]
\textbf{Énoncé.} Amina, 15 ans, élève en classe de Seconde à Maroua, est enceinte. Son père, très pauvre, l'avait promise en mariage à un commerçant de 35 ans qui aidait la famille financièrement. Amina ne connaissait ni le fonctionnement de son cycle menstruel ni l'existence de moyens de prévention. Identifie et classe les causes de cette grossesse précoce, puis propose une mesure de prévention pour chacune.
\tcblower
\textbf{Solution détaillée.}
\begin{itemize}
\item \textbf{Cause économique} : la pauvreté de la famille, qui accepte l'aide financière du prétendant. \emph{Prévention} : bourses scolaires, programmes de maintien des filles à l'école, autonomisation économique des familles.
\item \textbf{Cause sociale et familiale} : le mariage précoce imposé par la tradition, sans le consentement éclairé d'Amina. \emph{Prévention} : application de la loi (au Cameroun, le Code pénal réprime le mariage d'une fille mineure de 18 ans), sensibilisation des communautés et des chefs traditionnels.
\item \textbf{Cause individuelle} : la méconnaissance de la reproduction et des méthodes de prévention. \emph{Prévention} : éducation sexuelle à l'école et dans les centres de santé, information sur le cycle menstruel et la contraception.
\end{itemize}
\textbf{Conclusion} : la grossesse d'Amina résulte de la \emph{combinaison} de trois causes ; une prévention efficace devrait agir sur les trois à la fois.
\end{exoresolu}

\begin{savaistu}[La situation au Cameroun]
Selon les enquêtes démographiques et de santé (EDS-MICS 2018), environ \textbf{1 adolescente camerounaise sur 4} a déjà commencé une vie de mère avant l'âge de 20 ans : elle est soit déjà mère, soit enceinte de son premier enfant. Ce taux est plus élevé dans les régions de l'Extrême-Nord, du Nord et de l'Adamaoua, où les mariages précoces sont fréquents, et chez les jeunes filles non scolarisées ou issues des milieux les plus pauvres. À l'inverse, plus une jeune fille reste longtemps à l'école, plus le risque de grossesse précoce diminue : l'éducation est le meilleur « contraceptif social ».
\end{savaistu}

\begin{aretenir}[L'essentiel de la section]
\begin{itemize}
\item Une grossesse précoce survient avant 19 ans ; elle est le plus souvent non désirée, mais les deux notions ne sont pas synonymes.
\item Les causes sont multiples et combinées : rapports sexuels précoces et non protégés, méconnaissance de la reproduction, pressions sociales et économiques (pauvreté, partenaires plus âgés, mariages précoces), influence des pairs et des substances, violences sexuelles.
\item La responsabilité d'une grossesse due à une violence incombe toujours à l'agresseur, jamais à la victime.
\item Analyser une situation, c'est classer les causes (individuelles, familiales, sociales, économiques) avant de proposer des solutions adaptées à chacune.
\end{itemize}
\end{aretenir}

Maintenant que tu sais \emph{pourquoi} ces grossesses surviennent, la section suivante examinera leurs \emph{conséquences} — médicales, scolaires, sociales et économiques — pour l'adolescente, son enfant et la société tout entière.

\coursec{Grossesses précoces : conséquences}

Dans la section précédente, nous avons identifié les causes des grossesses précoces et non désirées : rapports sexuels non protégés, méconnaissance du fonctionnement du corps, pressions sociales, pauvreté, violences. Mais pourquoi faut-il à tout prix les éviter ? Parce qu'une grossesse survenant chez une adolescente dont le corps n'a pas fini de grandir engage deux vies — celle de la mère et celle de l'enfant — et bouleverse durablement un parcours scolaire, une famille, une communauté. C'est ce que nous allons montrer ici, de manière rigoureuse et documentée.

\courssub{Conséquences sur la santé de la mère}

\textbf{Une observation pour commencer.}\par

Dans les services de maternité des hôpitaux camerounais (Laquintinie à Douala, Hôpital Central à Yaoundé, hôpitaux régionaux), les sages-femmes et obstétriciens constatent un fait récurrent : les parturientes de 13 à 17 ans présentent des accouchements plus longs, plus difficiles et plus souvent compliqués que les femmes de 20 à 30 ans. Comment l'expliquer ?

\textbf{Le mécanisme : un bassin encore immature.}\par

À la puberté, la croissance staturale est achevée, mais la maturation du squelette, notamment du \cle{bassin}, se poursuit. Le \cle{vrai bassin} (ceinture osseuse par laquelle passe le nouveau-né lors de l'accouchement) n'atteint ses dimensions définitives, en particulier le diamètre transverse du détroit supérieur et le diamètre antéro-postérieur du détroit inférieur, que vers 18--20 ans. Chez une adolescente, le détroit du bassin peut être trop étroit pour laisser passer la tête du fœtus : on parle de \cle{disproportion fœto-pelvienne} (DFP). Le travail se prolonge alors pendant des heures, parfois des jours, avec trois issues possibles :

\begin{itemize}
\item une \textbf{césarienne} d'urgence, si l'on est près d'un hôpital équipé en bloc opératoire et banque de sang ;
\item une \textbf{fistule obstétricale}, si le travail prolongé comprime les tissus mères contre les os du bassin ;
\item le \textbf{décès} de la mère (hémorragie, rupture utérine, infection) ou de l'enfant (asphyxie périnatale), en l'absence de soins obstétricaux d'urgence (SONU).
\end{itemize}

\begin{definition}[Fistule obstétricale]
La \cle{fistule obstétricale} est une communication anormale, acquise par nécrose ischémique, entre le \textbf{vagin} et la \textbf{vessie} (fistule vésico-vaginale, FVV) ou entre le vagin et le \textbf{rectum} (fistule recto-vaginale, FRV). Elle résulte d'un travail d'accouchement prolongé et obstructif : la tête du fœtus comprime durablement les tissus mous (paroi vaginale, vessie, urètre, rectum) contre les os du bassin, entraînant leur nécrose par ischémie. Conséquence : des fuites permanentes et incontrôlées d'urine (FVV) ou de matières fécales (FRV), souvent associées.
\end{definition}

\begin{savaistu}[La fistule obstétricale, drame silencieux et évitable]
La fistule obstétricale touche surtout les très jeunes mères (souvent \textless 18 ans) des zones rurales pauvres, éloignées des centres de santé capables de pratiquer une césarienne. Au-delà de la souffrance physique (dermatites, infections urinaires récidivantes, infertilité secondaire), les fuites permanentes d'urine provoquent odeurs et stigmatisation : la jeune femme est souvent rejetée par son époux, abandonnée par sa famille, exclue de la vie communautaire et scolaire. Pourtant, cette lésion est \textbf{réparable chirurgicalement} (fistulectomie, plastie vésicale ou rectale) avec un taux de succès élevé (> 90 \%) si elle est prise en charge par des chirurgiens formés. Au Cameroun, le Ministère de la Santé Publique (MINSANTÉ), avec l'appui de l'UNFPA et de partenaires (ONG comme \emph{Fistula Foundation}, \emph{MSF}), organise régulièrement des \textbf{campagnes de prise en charge gratuite} (dépistage, chirurgie, réinsertion socio-économique) dans les hôpitaux de référence (ex. Hôpital Gynéco-Obstétrique et Pédiatrique de Yaoundé, Hôpital de District de Bafia, etc.). La meilleure solution reste la \textbf{prévention primaire} : éviter les grossesses avant la maturité pelvienne complète (18--20 ans), assurer un suivi prénatal de qualité et garantir l'accès à la césarienne en cas de DFP (SONU).
\end{savaistu}

\textbf{Les autres risques médicaux majeurs.}\par

\begin{itemize}
\item Les \textbf{hémorragies de la délivrance (HDD)} : l'utérus immature (hypoplasie utérine relative) se contracte moins efficacement après l'expulsion du fœtus (atonie utérine), ce qui favorise les hémorragies du post-partum immédiat, \textbf{première cause de mortalité maternelle} au Cameroun et dans le monde.
\item L'\textbf{anémie gravidique} (besoins en fer accrus pour la croissance maternelle + fœtus) et la \textbf{pré-éclampsie} (hypertension artérielle gravidique avec protéinurie, survenant avant 37 SA), dont l'incidence est significativement plus élevée chez les adolescentes (\textless 18 ans) que chez les femmes de 20--29 ans.
\item Les \textbf{avortements non sécurisés (clandestins)} : face à une grossesse non désirée, certaines adolescentes, par peur, honte, ou contrainte légale/financière, ont recours à des interruptions volontaires de grossesse (IVG) pratiquées dans de mauvaises conditions (personnel non qualifié, instruments non stériles, produits toxiques -- décoctions, cytotéc mal dosé, corps étrangers). Les conséquences sont les \textbf{infections} graves (endométrite, salpingite, péritonite, septicémie), les \textbf{hémorragies} massives, les \textbf{perforations utérines} (nécessitant souvent hystérectomie), la \textbf{stérilité} définitive (syndrome d'Asherman, obstruction tubaire), et trop souvent le \textbf{décès maternel}. Au Cameroun, les complications des avortements non sécurisés représentent une part majeure des admissions en urgence gynécologique.
\end{itemize}

\begin{exemplebox}[Un chiffre qui doit alerter]
Selon l'Organisation mondiale de la santé (OMS, 2023), les complications liées à la grossesse et à l'accouchement (hémorragies, infections, pré-éclampsie, avortements non sécurisés, travail obstructif) constituent la \textbf{deuxième cause de mortalité} chez les filles de 15 à 19 ans dans le monde (après les blessures non intentionnelles), et la \textbf{première cause} dans les pays à revenu faible ou intermédiaire. Autrement dit, pour une adolescente, tomber enceinte n'est pas seulement un problème social : c'est un \textbf{risque vital} immédiat.
\end{exemplebox}

\courssub{Conséquences sur l'enfant}

L'enfant né d'une mère très jeune (\textless 18 ans) paie lui aussi un lourd tribut, lié à l'immaturité biologique de la mère et à la compétition nutritionnelle mère-fœtus :

\begin{itemize}
\item la \textbf{prématurité} (naissance avant 37 semaines d'aménorrhée -- SA) est 1,5 à 2 fois plus fréquente ;
\item le \textbf{faible poids à la naissance} (PN \textless \SI{2500}{g}, souvent \textless \SI{2000}{g}) expose le nouveau-né à l'hypothermie, l'hypoglycémie, les infections néonatales, les difficultés d'alimentation (réflexes de succion/déglutition immatures) et un développement neuro-moteur retardé ;
\item la \textbf{mortalité néonatale} (premiers 28 jours) et \textbf{infantile} (première année) est nettement plus élevée : les nouveau-nés de mères adolescentes présentent un risque de décès dans la première année supérieur de 50 à 100 \% à celui des enfants de mères de 20 à 29 ans ;
\item à plus long terme, l'enfant grandit souvent dans la \textbf{précarité}, avec une alimentation insuffisante (sevrage précoce, diversification tardive), un suivi médical irrégulier (vaccination incomplète), et une stimulation cognitive faible, ce qui compromet sa croissance staturo-pondérale, son développement cognitif et sa scolarité future.
\end{itemize}

\courssub{Conséquences scolaires}

Une grossesse précoce interrompt presque toujours le parcours scolaire de l'adolescente :

\begin{itemize}
\item \textbf{Abandon des études} : la jeune mère quitte l'école (souvent définitivement) pour s'occuper de l'enfant ou subvenir à ses besoins par des activités informelles ;
\item même en cas de retour possible (réintégration autorisée par la circulaire MINEDUB/MINESEC), les \textbf{absences répétées} (consultations prénatales, accouchement, congé de maternité, maladies de l'enfant, manque de garde) provoquent des lacunes, des échecs et des redoublements ;
\item la \textbf{perte des chances de formation} est durable : sans diplôme (BEPC, Probatoire, Baccalauréat), l'accès à un emploi qualifié ou à l'enseignement supérieur devient très difficile, ce qui enferme la jeune femme dans la pauvreté et la dépendance.
\end{itemize}

Notons que le garçon responsable de la grossesse peut lui aussi être déscolarisé (exclusion temporaire, pression familiale), mais dans les faits, c'est presque toujours la fille qui porte seule, et de manière disproportionnée, les conséquences scolaires, sanitaires et sociales.

\courssub{Conséquences économiques}

Sans diplôme ni formation professionnelle, la jeune mère ne peut accéder qu'à des activités précaires et mal rémunérées (commerce ambulant, travail domestique, agriculture de subsistance). Il en résulte :

\begin{itemize}
\item une \textbf{dépendance économique} vis-à-vis de la famille d'origine (parents, grands-parents) ou du partenaire (souvent lui-même précaire) ;
\item une grande \textbf{difficulté à subvenir aux besoins de l'enfant} : alimentation diversifiée, soins de santé (mutuelle, médicaments), vêtements, frais de scolarité (même publics : PTA, fournitures, cantine) ;
\item un \textbf{cercle vicieux de la pauvreté intergénérationnelle} : la pauvreté favorise les grossesses précoces (manque d'information, échange sexe/argent, mariage d'enfants), qui elles-mêmes aggravent la pauvreté (perte de capital humain), et l'enfant né dans ces conditions a lui-même plus de risques de reproduire ce parcours (faible capital santé/éducation, reproduction précoce).
\end{itemize}

\courssub{Conséquences sociales et psychologiques}

\begin{itemize}
\item la \textbf{stigmatisation} : la jeune fille enceinte est souvent montrée du doigt, insultée (\"fille facile\", \"mal élevée\"), exclue des activités de la communauté, de l'église, des groupes de pairs ;
\item le \textbf{rejet familial} : certains parents chassent leur fille de la maison (\"honte sur la famille\"), la laissant sans ressource ni hébergement, ce qui la pousse vers la rue ou des unions précaires ;
\item la \textbf{honte}, la culpabilité, l'isolement social, la perte d'estime de soi, qui peuvent conduire à la \textbf{dépression post-partum}, à l'anxiété, voire à des conduites dangereuses (récidive d'avortement non sécurisé, fugue, consommation de substances, tentative de suicide) ;
\item les \textbf{mariages forcés ou précoces} précipités pour « réparer » la situation (dot versée, union coutumière/religieuse), qui compromettent définitivement l'autonomie et l'avenir de l'adolescente, la plaçant sous l'autorité d'un époux souvent plus âgé.
\end{itemize}

\courssub{Conséquences pour la communauté et la nation}

Les grossesses précoces ne concernent pas que l'individu : elles pèsent sur toute la société et freinent le développement national (Objectifs de Développement Durable -- ODD 3, 4, 5, 8, 10).

\begin{itemize}
\item \textbf{Paupérisation du capital humain} : chaque adolescente déscolarisée est une force de travail qualifiée et un potentiel d'innovation perdus pour le développement du pays (dividende démographique non capté).
\item \textbf{Charge pour les familles élargies} : ce sont souvent les grands-mères, déjà vulnérables, qui élèvent les enfants des adolescentes, au prix de leurs maigres ressources (pensions, petit commerce), sacrifiant parfois la scolarité de leurs propres enfants cadets.
\item \textbf{Charge pour le système de santé} : les complications (césariennes d'urgence, réparations de fistules, soins aux grands prématurés en néonatologie, prise en charge des suites d'avortements non sécurisés) mobilisent des lits, du sang, des médicaments, du personnel qualifié et des moyens financiers coûteux qui manquent pour d'autres besoins de santé publique (paludisme, VIH, maladies non transmissibles).
\end{itemize}

\begin{popfigure}
\begin{center}
\renewcommand{\arraystretch}{1.4}
\begin{tabularx}{\linewidth}{|l|X|X|X|}
\hline
\rowcolor{popblueL} \textbf{Type de conséquence} & \textbf{Pour la mère} & \textbf{Pour l'enfant} & \textbf{Pour la société} \\
\hline
\textbf{Sanitaires} & Disproportion fœto-pelvienne, travail prolongé, \textbf{fistule obstétricale}, \textbf{hémorragie de la délivrance}, pré-éclampsie, séquelles d'avortement non sécurisé, \textbf{décès maternel} & \textbf{Prématurité}, \textbf{faible poids de naissance} (\textless \SI{2500}{g}), \textbf{mortalité néonatale/infantile} accrue, retard de croissance/développement & Services de santé surchargés (urgences obstétricales, néonatologie), coûts élevés (SONU, chirurgie réparatrice), consommation de produits sanguins \\
\hline
\textbf{Scolaires} & \textbf{Abandon des études}, perte de formation, illettrisme fonctionnel & Scolarité future compromise (retard cognitif, absentéisme, manque de soutien) & Perte de main-d'œuvre qualifiée, baisse du taux d'achèvement primaire/secondaire (filles), gaspillage investissement éducatif \\
\hline
\textbf{Économiques} & \textbf{Précarité}, dépendance, emploi informel, bas revenu & Besoins non couverts (nutrition, santé, éducation), travail des enfants & Paupérisation des ménages, cercle vicieux pauvreté intergénérationnelle, coût social (assistance) \\
\hline
\textbf{Sociales et psychologiques} & Stigmatisation, rejet familial, dépression, anxiété, mariage forcé/précoce & Abandon, absence de repères parentaux, risque de maltraitance & Exclusion sociale, tensions familiales/communautaires, instabilité \\
\hline
\end{tabularx}
\end{center}
\legende{Tableau récapitulatif des conséquences des grossesses précoces, selon le domaine et la personne concernée. À l'examen, sache mobiliser au moins une conséquence précise par case (mécanisme + terme médical exact).}
\end{popfigure}

\begin{popfigure}
\begin{center}
\begin{tikzpicture}[
  node distance=0.6cm,
  box/.style={draw=popblue, fill=popblueL, rounded corners=3pt, align=center, text width=3.2cm, minimum height=1.1cm, font=\small},
  badbox/.style={draw=poporange, fill=poporangeL, rounded corners=3pt, align=center, text width=3.2cm, minimum height=1.1cm, font=\small},
  arr/.style={-{Stealth[length=3mm]}, thick, popdark}
]
\node[badbox] (g) {\textbf{Grossesse précoce} non désirée (\textless 18 ans)};
\node[box, below=of g] (d) {\textbf{Déscolarisation} : abandon, perte de diplôme, illettrisme};
\node[box, below=of d] (p) {\textbf{Précarité économique} : pas d'emploi qualifié, dépendance, bas revenu};
\node[box, below=of p] (e) {\textbf{Difficultés pour l'enfant} : malnutrition, soins insuffisants, échec scolaire, santé fragile};
\node[badbox, below=of e] (r) {Risque de \textbf{reproduction du cycle} à la génération suivante (pauvreté + grossesse précoce)};
\draw[arr] (g) -- (d);
\draw[arr] (d) -- (p);
\draw[arr] (p) -- (e);
\draw[arr] (e) -- (r);
\draw[arr, poporange, dashed, bend left=45] (r.west) to (g.west);
\node[poporange, font=\small\itshape, align=center] at (-4.5, -5.5) {\textbf{Cercle vicieux}\\\textbf{intergénérationnel}};
\end{tikzpicture}
\end{center}
\legende{Chaîne des conséquences d'une grossesse précoce : chaque maillon aggrave le suivant, formant un cercle vicieux intergénérationnel (flèche en pointillés) que seule la prévention (information, contraception, maintien à l'école, autonomisation) permet de briser durablement.}
\end{popfigure}

\begin{attention}[Erreur fréquente à l'examen / Piège de rédaction]
Beaucoup d'élèves ne citent que les conséquences \textbf{scolaires} (« elle arrête l'école ») ou sociales (« on la montre du doigt ») en oubliant les risques \textbf{médicaux graves et spécifiques à l'adolescente}. Or les plus lourdes conséquences d'une grossesse précoce sont \textbf{sanitaires} : \textbf{disproportion fœto-pelvienne} (bassin immature), \textbf{fistules obstétricales} (lésion invalidante), \textbf{hémorragies de la délivrance} (première cause de mortalité maternelle), \textbf{pré-éclampsie}, \textbf{avortements non sécurisés mortels}. Une réponse complète et probatoire doit \textbf{toujours commencer par la santé de la mère et de l'enfant} (mécanismes physiopathologiques), avant d'aborder les dimensions scolaires, économiques et sociales. N'oubliez pas le vocabulaire précis : \"fistule vésico-vaginale\", \"travail obstructif\", \"faible poids de naissance\", \"mortalité néonatale\".
\end{attention}

\courssub{Savoir expliquer ces risques : un acte de prévention (Savoir-faire)}

Comprendre ces conséquences ne suffit pas : le programme attend de toi que tu saches les \textbf{expliquer à un pair} (camarade, frère/sœur cadet, ami) pour l'aider à prendre une décision éclairée. Informer sans juger, c'est déjà prévenir.

\begin{methode}[Construire un argumentaire de prévention en trois étapes (Probatoire)]
\begin{enumerate}
\item \textbf{Énoncer le risque biologique} clairement, sans exagération ni catastrophisme, en nommant le mécanisme : « Chez une fille de 15--16 ans, le vrai bassin (détroit) n'a pas atteint sa taille adulte. Si elle tombe enceinte, la tête du bébé peut ne pas passer : c'est la disproportion fœto-pelvienne. Le travail se bloque, ce qui cause des fistules ou des hémorragies mortelles si elle n'est pas opérée à temps. »
\item \textbf{Donner une preuve concrète} (chiffre officiel, exemple local, témoignage) qui ancre le risque dans la réalité : « Au Cameroun, les complications de la grossesse sont une cause majeure de décès chez les 15--19 ans. À l'hôpital Laquintinie, on répare encore trop de fistules sur des filles de 14--17 ans qui ont accouché chez elles ou trop tard. Même sans drame médical, 9 filles sur 10 qui tombent enceintes au lycée n'ont pas le Bac. »
\item \textbf{Proposer une attitude responsable et une ressource} : « Le plus sûr pour ton avenir et ta santé, c'est de terminer tes études d'abord. Si tu as des rapports, protège-toi (préservatif) et va te faire conseiller \textbf{gratuitement et confidentiellement} dans un Centre de Santé (CSI, CMA, Hôpital) ou un Centre d'Information et d'Orientation Jeunes (CIOJ) : on t'y expliquera la contraception sans te juger. »
\end{enumerate}
Un bon argumentaire reste \textbf{respectueux, factuel, empathique et orienté solution} : il informe sans humilier, sans menacer, et ouvre une porte d'action concrète.
\end{methode}

\begin{exoresolu}[Expliquer les risques à un pair -- Mise en situation probatoire]
\textbf{Énoncé.} Ta camarade de classe, âgée de 16 ans, en classe de Première, confie qu'elle envisage d'avoir un enfant avec son petit ami car « beaucoup de filles le font dans le quartier et s'en sortent bien ». Rédige un court argumentaire structuré (trois étapes) pour l'éclairer sur les risques réels, sans la juger ni la moraliser.
\tcblower
\textbf{Solution détaillée.}
\begin{enumerate}
\item \emph{Énoncer le risque biologique} : « Écoute, à 16 ans, ton corps est encore en croissance. Ton bassin (le passage pour le bébé) n'est pas encore assez large. Si tu tombes enceinte maintenant, tu risques une \textbf{disproportion fœto-pelvienne} : le bébé ne passe pas, le travail dure trop longtemps, et ça peut te causer une \textbf{fistule} (fuites d'urine permanentes) ou une \textbf{hémorragie grave}. Pour ton bébé, le risque, c'est la \textbf{prématurité} et un \textbf{poids trop faible} à la naissance. »
\item \emph{Donner une preuve concrète} : « Ce n'est pas pour te faire peur, c'est la réalité des urgences ici. L'OMS dit que les complications de la grossesse sont la \textbf{deuxième cause de mortalité} mondiale chez les 15--19 ans. Au Cameroun, nos maternités sont pleines de jeunes filles qu'on opère en urgence (césarienne) ou qu'on soigne pour des suites d'avortements mal faits. Et même si l'accouchement se passe bien, regarde autour de toi : celles qui ont eu un bébé en Première, combien ont eu le Bac ? Presque aucune. Elles galèrent pour nourrir l'enfant, sans diplôme, sans argent. »
\item \emph{Proposer une attitude responsable et une ressource} : « Ton avenir, c'est ton Bac, puis une formation, un métier. C'est comme ça que tu pourras offrir une vraie vie à un enfant, plus tard, quand \textbf{toi} tu seras prête (corps mature, diplôme, travail). En attendant, si vous avez des rapports avec ton copain, \textbf{protégez-vous} (préservatif = double protection grossesse + IST). Et vas voir l'infirmière du lycée ou le CSI du quartier : c'est gratuit, confidentiel, et ils te donneront la pilule ou l'implant si tu veux. Réfléchis-y, ton avenir en vaut la peine. »
\end{enumerate}
\emph{Remarque pédagogique} : Cet argumentaire est efficace car il suit la méthode (Mécanisme $\rightarrow$ Preuve locale/Chiffre $\rightarrow$ Solution/Action), utilise le vocabulaire scientifique exact (disproportion fœto-pelvienne, fistule, prématurité), cite une source d'autorité (OMS, hôpital local), et se termine par une orientation concrète vers le service de santé (comportement responsable).
\end{exoresolu}

\begin{aretenir}[L'essentiel de la section -- Fiche de révision Probatoire]
\begin{itemize}
\item La grossesse précoce (\textless 18 ans) met en danger vital la \textbf{santé de la mère} : \textbf{bassin immature} $\Rightarrow$ \textbf{disproportion fœto-pelvienne}, travail obstructif $\Rightarrow$ \textbf{fistules obstétricales} (FVV, FRV), \textbf{hémorragies de la délivrance} (atonie utérine), \textbf{pré-éclampsie}, séquelles d'\textbf{avortements non sécurisés} (septicémie, perforation, stérilité, décès).
\item Elle menace gravement l'\textbf{enfant} : \textbf{prématurité}, \textbf{faible poids de naissance} (\textless \SI{2500}{g}), \textbf{mortalité néonatale/infantile} majorée, retard de croissance/développement.
\item Elle provoque la \textbf{déscolarisation} quasi systématique, la \textbf{précarité économique}, la \textbf{stigmatisation sociale}, des \textbf{troubles psychologiques} (dépression), et des \textbf{mariages précoces/forcés}.
\item Ces conséquences s'enchaînent en \textbf{cercle vicieux intergénérationnel} : grossesse précoce $\rightarrow$ déscolarisation $\rightarrow$ précarité $\rightarrow$ difficultés pour l'enfant $\rightarrow$ reproduction du cycle.
\item \textbf{Savoir-faire probatoire} : Construire un argumentaire de prévention en 3 étapes (Risque biologique nommé $\rightarrow$ Preuve chiffrée/locale $\rightarrow$ Attitude responsable + Ressource santé) pour informer un pair sans le juger.
\end{itemize}
\end{aretenir}

\coursec{Infections sexuellement transmissibles (IST)}

Dans la section précédente, nous avons vu que les grossesses précoces et non désirées constituent l'un des risques majeurs des rapports sexuels non protégés chez les adolescents. Mais la grossesse n'est pas le seul danger : le même rapport non protégé peut transmettre des maladies parfois graves, parfois silencieuses, parfois incurables. Ce sont les \cle{infections sexuellement transmissibles} (IST). Comprendre ce que sont ces infections, comment elles se transmettent et comment elles évoluent est indispensable pour adopter un comportement responsable en matière de santé reproductive.

\courssub{Définition et situation problème}

\textbf{Une observation de la vie courante.}\par
Dans un centre de santé de Yaoundé, une jeune fille de 19 ans consulte pour des douleurs persistantes du bas-ventre. Après examen et analyses, le médecin diagnostique une infection des trompes de Fallope (salpingite) causée par une \emph{chlamydiose} jamais dépistée ni traitée. Pourtant, la jeune fille n'a jamais ressenti le moindre symptôme auparavant : elle se croyait en parfaite santé. Comment une maladie peut-elle détruire un organe sans jamais se manifester ? C'est toute la particularité des IST.

\begin{definition}[Infection sexuellement transmissible (IST)]
Une \cle{infection sexuellement transmissible} (IST) est une infection causée par un micro-organisme (virus, bactérie ou parasite) et transmise \textbf{principalement lors de rapports sexuels non protégés} (sans préservatif), qu'ils soient vaginaux, anaux ou bucco-génitaux. Certaines IST se transmettent aussi par le \textbf{sang contaminé} ou \textbf{de la mère à l'enfant} pendant la grossesse, l'accouchement ou l'allaitement.
\end{definition}

On distinguait autrefois les « maladies sexuellement transmissibles » (MST) ; le terme IST est aujourd'hui préféré, car il insiste sur un point essentiel : on peut être \emph{infecté} (porter le microbe et le transmettre) sans être \emph{malade} (sans aucun signe visible).

\begin{propriete}[Deux grandes familles d'IST]
Selon la nature de l'agent responsable, on distingue :
\begin{itemize}
\item les \cle{IST bactériennes} : gonococcie (blennorragie), syphilis, chlamydioses. Elles se \textbf{guérissent} par un traitement antibiotique bien suivi, à condition d'être dépistées tôt ;
\item les \cle{IST virales} : VIH/SIDA, hépatite B, herpès génital, papillomavirus (HPV). Elles ne se guérissent généralement pas : certaines se \textbf{contrôlent} par des traitements au long cours (VIH, hépatite B chronique), d'autres \textbf{persistent} toute la vie dans l'organisme avec des réactivations possibles (herpès).
\end{itemize}
La trichomonose, due à un parasite unicellulaire (\emph{Trichomonas vaginalis}), se guérit également par traitement antibiotique (métronidazole).
\end{propriete}

\begin{attention}[Erreur fréquente : « on voit les malades à leur apparence »]
Il est \textbf{impossible} de reconnaître une personne infectée à son apparence. Une personne séropositive au VIH ou porteuse de chlamydiae peut avoir un aspect parfaitement sain pendant des années. De nombreuses IST sont \cle{asymptomatiques} (silencieuses), surtout chez la femme. Seul un \textbf{test de dépistage} en laboratoire permet de savoir si l'on est infecté. Dire « il/elle a l'air en bonne santé, donc pas de risque » est un raisonnement dangereux et scientifiquement faux.
\end{attention}

\courssub{Les principales IST : agents, signes et complications}

\textbf{Démarche d'investigation : l'identification des agents.}\par
Comment les scientifiques identifient-ils l'agent responsable d'une IST ? Prenons l'exemple historique de la blennorragie.
\begin{enumerate}
\item \textbf{Observation} : des malades présentent un écoulement purulent par l'urètre et des brûlures en urinant.
\item \textbf{Hypothèse} : un microbe présent dans le pus est responsable de la maladie.
\item \textbf{Expérience/données} : en 1879, Albert Neisser observe au microscope, dans le pus des malades, une bactérie en forme de grains de café accolés (diplocoque) : le \emph{gonocoque} (\emph{Neisseria gonorrhoeae}). La bactérie est ensuite cultivée et retrouvée chez tous les malades, jamais chez les personnes saines.
\item \textbf{Interprétation} : le gonocoque, déposé sur les muqueuses génitales lors d'un rapport non protégé, s'y multiplie et provoque l'inflammation.
\item \textbf{Conclusion} : la blennorragie est une infection bactérienne sexuellement transmissible ; un antibiotique adapté la guérit.
\end{enumerate}
C'est cette même démarche (observation microscopique, culture, tests sérologiques ou moléculaires) qui a permis d'identifier l'agent de chaque IST : le tréponème pâle pour la syphilis (Schaudinn et Hoffmann, 1905), le VIH pour le SIDA (Montagnier, Barré-Sinoussi et al., 1983), etc.

\textbf{Les signes évocateurs.}\par
Lorsqu'ils existent, les signes les plus fréquents d'une IST sont :
\begin{itemize}
\item un \cle{écoulement} anormal par le pénis ou le vagin (purulent, jaunâtre ou verdâtre, parfois malodorant) ;
\item des \cle{brûlures urinaires} (dysurie : douleur en urinant) ;
\item des \cle{ulcérations génitales} : chancre indolore de la syphilis (stade primaire), vésicules douloureuses regroupées en bouquets de l'herpès ;
\item des douleurs du bas-ventre, des démangeaisons, des verrues génitales (condylomes, dus aux HPV de bas risque oncogénique).
\end{itemize}
Mais rappelons-le : \textbf{l'absence de signes ne signifie pas l'absence d'infection}. La chlamydiose est asymptomatique dans environ 70 \% des cas chez la femme et 50 \% chez l'homme ; le VIH peut rester cliniquement silencieux 8 à 10 ans (phase de latence clinique).

\textbf{Les complications en l'absence de traitement.}\par
Une IST non traitée peut entraîner des séquelles graves, souvent irréversibles :
\begin{itemize}
\item la \cle{stérilité} (obstruction des trompes de Fallope par fibrose séquellaire de salpingites à répétition chez la femme ; atteinte des épididymes ou canaux déférents chez l'homme), surtout avec la chlamydiose et la gonococcie ;
\item les \cle{grossesses extra-utérines} (l'œuf se développe dans une trompe abîmée : urgence vitale par rupture hémorragique) ;
\item des \cle{cancers} : certains papillomavirus à haut risque (HPV 16 et 18 principalement) sont la cause nécessaire du cancer du col de l'utérus, l'un des cancers les plus fréquents et meurtriers chez la femme camerounaise ;
\item la \cle{transmission mère-enfant} : le VIH, la syphilis, l'hépatite B ou le gonocoque peuvent infecter le fœtus ou le nouveau-né (la syphilis non traitée provoque fausses couches, mortinaissances et malformations congénitales ; le gonocoque peut provoquer une blennorragie néonatale avec risque de cécité) ;
\item l'\cle{affaiblissement du système immunitaire} : le VIH détruit progressivement les lymphocytes T4 (CD4+), cellules chefs d'orchestre de la réponse immune adaptative ; sans traitement, survient le \cle{SIDA} (stade d'immunodépression sévère), où le corps ne se défend plus contre les infections dites « opportunistes » (tuberculose, méningites cryptococciques, pneumonies à \emph{Pneumocystis jirovecii}, etc.).
\end{itemize}

\begin{center}
\rowcolors{1}{popcream}{white}
\begin{tabularx}{\linewidth}{|l|l|X|X|X|}
\hline
\rowcolor{popblueL} \textbf{IST} & \textbf{Agent} & \textbf{Signes possibles} & \textbf{Complications majeures} & \textbf{Évolution / Prise en charge} \\
\hline
VIH/SIDA & Virus (VIH-1/2) & Longue phase silencieuse, puis infections répétées, amaigrissement & Destruction des défenses immunitaires (SIDA), cancers & Non guérison ; contrôlée à vie par antirétroviraux (ARV) ; charge virale indétectable = intransmissibilité (I=I) \\
\hline
Gonococcie (blennorragie) & Bactérie (\emph{Neisseria gonorrhoeae}) & Écoulement purulent abondant, brûlures urinaires vives & Stérilité, blennorragie néonatale (cécité), septicémie & Guérison par antibiotiques (résistances croissantes : surveillance) \\
\hline
Syphilis & Bactérie (\emph{Treponema pallidum}) & Chancre indolore (stade 1), éruptions cutanées (stade 2), atteintes viscérales/neurologiques (stade 3) & Atteintes cardiaques, neurologiques (neurosyphilis), mort fœtale & Guérison par pénicilline (efficace à tous stades, mais séquelles irréversibles si tardif) \\
\hline
Chlamydiose & Bactérie (\emph{Chlamydia trachomatis}) & Souvent aucun signe ; écoulement clair, brûlures modérées & Stérilité tubaire, grossesse extra-utérine, pneumonie/ conjonctivite néonatale & Guérison par antibiotiques (azithromycine, doxycycline) \\
\hline
Trichomonose & Parasite (\emph{Trichomonas vaginalis}) & Écoulement mousseux jaune-vert, démangeaisons vulvaires & Inflammation facilitant transmission VIH, accouchement prématuré & Guérison par métronidazole (traitement partenaire obligatoire) \\
\hline
Hépatite B & Virus (VHB) & Fatigue, jaunisse (ictère), urines foncées ; souvent silencieuse & Cirrhose, cancer du foie (hépatocarcinome) & Guérison spontanée fréquente (adulte) ; chronicité possible contrôlée par antiviraux ; \textbf{vaccin préventif} (EPI) \\
\hline
Herpès génital & Virus (HSV-1/2) & Vésicules douloureuses récidivantes, ulcérations & Poussées répétées, transmission néonatale grave (encéphalite) & Non éradication (latence ganglionnaire) ; antiviraux pour réduire poussées \\
\hline
Papillomavirus (HPV) & Virus (HPV) & Verrues génitales (condylomes, bas risque), lésions précancéreuses invisibles (haut risque) & Cancer du col utérus, anus, oropharynx & Non traitement viral ; traitement des lésions ; \textbf{vaccin préventif} (filles/garçons) \\
\hline
\end{tabularx}
\end{center}

\begin{savaistu}[Une infection silencieuse peut rendre stérile]
La chlamydiose est l'une des IST les plus fréquentes chez les jeunes de 15 à 24 ans au Cameroun comme ailleurs. Chez la femme, elle évolue le plus souvent \textbf{sans aucun symptôme}. Silencieusement, la bactérie remonte du col utérin vers l'utérus puis les trompes de Fallope, qu'elle enflamme (salpingite) et obstrue par fibrose. Des années plus tard, la jeune femme découvre qu'elle ne peut plus avoir d'enfant, sans jamais avoir su qu'elle était malade. D'où l'importance capitale du \textbf{dépistage régulier} (urines ou prélèvement vaginal), même en l'absence de tout signe, surtout après un rapport à risque ou un changement de partenaire.
\end{savaistu}

\courssub{Les modes de transmission : focus sur le VIH}

\textbf{Observation.}\par
Lors des campagnes de lutte contre le SIDA au Cameroun, une question revient sans cesse : « Peut-on attraper le VIH en serrant la main d'un malade, en partageant son repas ou par une piqûre de moustique ? ». Les données épidémiologiques mondiales, accumulées depuis plus de quarante ans, sont formelles : \textbf{non}. Le VIH est un virus enveloppé, fragile hors du corps, qui ne survit pas longtemps dans l'environnement et qui n'est présent en quantité infectante suffisante que dans certains liquides biologiques : le sang, le sperme, les sécrétions vaginales, le liquide séminal et le lait maternel. La salive, les larmes, la sueur, l'urine ne transmettent pas le VIH (charge virale trop faible, présence d'inhibiteurs).

\begin{propriete}[Les trois modes de transmission du VIH]
Le VIH se transmet uniquement par trois voies :
\begin{enumerate}
\item \textbf{Voie sexuelle} : rapports sexuels (vaginaux, anaux, bucco-génitaux) non protégés avec une personne infectée (voie la plus fréquente au Cameroun, > 90 \% des cas) ;
\item \textbf{Voie sanguine} : transfusion de sang contaminé (très rare aujourd'hui grâce au dépistage systématique), partage de seringues/aiguilles (usagers de drogues injectables), utilisation de lames ou d'objets tranchants souillés non stérilisés (scarifications traditionnelles, circoncision hors structure médicale, tatouages, perçages) ;
\item \textbf{Transmission mère-enfant (verticale)} : pendant la grossesse (transplacentaire), l'accouchement (contact sang/sécrétions) ou l'allaitement, si la mère infectée n'est pas sous traitement antirétroviral efficace.
\end{enumerate}
L'hépatite B (VHB) se transmet par ces mêmes trois voies (le VHB est 50 à 100 fois plus infectieux que le VIH par voie sanguine). En revanche, la gonococcie, la syphilis, les chlamydioses, la trichomonose, l'herpès et les HPV se transmettent essentiellement par contact sexuel direct avec les muqueuses ou lésions (et de la mère à l'enfant lors de l'accouchement pour certaines).
\end{propriete}

\begin{popfigure}
\begin{tikzpicture}[font=\small, >=Stealth, node distance=1.2cm and 1.5cm]
% Central node
\node[draw=popdark, fill=popblueL, rounded corners, thick, align=center, minimum width=3.5cm, minimum height=1.2cm] (vih) {\textbf{Transmission du VIH}};
% Three modes
\node[draw=popdark, fill=poporangeL, rounded corners, align=center, minimum width=3.8cm] (sex) at (-5.5,2.8) {\textbf{Voie sexuelle}\\ Rapports non protégés\\ (sperme, sécrétions vaginales,\\ liquide séminal)};
\node[draw=popdark, fill=poppinkL, rounded corners, align=center, minimum width=3.8cm] (sang) at (0,3.6) {\textbf{Voie sanguine}\\ Seringues partagées, sang\\ contaminé, lames/objets\\ tranchants non stérilisés};
\node[draw=popdark, fill=popgreenL, rounded corners, align=center, minimum width=3.8cm] (mere) at (5.5,2.8) {\textbf{Mère -- enfant}\\ Grossesse, accouchement,\\ allaitement (si mère non traitée)};
\draw[->, thick, popdark] (sex.south) -- (vih.north west);
\draw[->, thick, popdark] (sang.south) -- (vih.north);
\draw[->, thick, popdark] (mere.south) -- (vih.north east);
% Non-transmission box
\node[draw=popgreen, fill=popgreenL, rounded corners, thick, align=center, text width=12cm] (non) at (0,-2.8) {\textbf{Le VIH ne se transmet PAS par la vie quotidienne :} poignée de main, embrassades, baisers sociaux, partage des repas/verre/couverts, toilettes, piscine, sueur, larmes, toux, éternuements, piqûres de moustiques (le moustique n'injecte pas le sang du repas précédent, le virus ne se réplique pas dans le moustique).};
\draw[->, thick, popgreen] (vih.south) -- (non.north);
\end{tikzpicture}
\legende{Les trois modes de transmission du VIH (en haut) et les contacts de la vie quotidienne qui ne transmettent jamais le virus (en bas). Le VIH n'est présent en quantité infectante que dans le sang, le sperme, le liquide séminal, les sécrétions vaginales et le lait maternel.}
\end{popfigure}

\begin{attention}[Ni poignée de main, ni moustique : lutter contre la stigmatisation]
Le VIH ne se transmet \textbf{jamais} par la poignée de main, les embrassades, le partage d'un repas ou d'un verre, les toilettes, la piscine, la sueur, ni par les piqûres de moustiques. Croire le contraire entretient la \textbf{stigmatisation} et la discrimination envers les personnes vivant avec le VIH (PVVIH), les éloignant du dépistage, des soins et de leur entourage. Vivre, étudier, manger, travailler ou jouer avec une personne séropositive sous traitement (charge virale indétectable) ne présente \textbf{aucun danger}. La solidarité, c'est aussi la prévention.
\end{attention}

\courssub{Dépistage, traitement et prévention : agir tôt et se protéger}

\textbf{Pourquoi et comment dépister ?}\par
Puisque beaucoup d'IST sont silencieuses, la seule stratégie efficace est le \cle{dépistage} : une prise de sang (sérologie VIH, syphilis, hépatites B/C) et/ou un prélèvement local (urines, frottis urétral/vaginal/cervical pour gonocoque, chlamydiae, trichomonas, PCR HPV) analysé en laboratoire. Au Cameroun, le test de dépistage du VIH (TROD : Test Rapide d'Orientation Diagnostique ou ELISA) est \textbf{gratuit, confidentiel et accessible} dans tous les hôpitaux, centres de santé agréés et lors des campagnes mobiles (Journée mondiale de lutte contre le SIDA, vacances « Sans SIDA », consultations prénatales). Le dépistage est recommandé : après tout rapport à risque, avant d'arrêter le préservatif dans un couple stable, au début d'une grossesse, et de façon régulière pour les personnes multipartenaires.

\begin{propriete}[Guérir, contrôler ou prévenir : les trois piliers]
\begin{itemize}
\item \textbf{Guérir} : Les IST \textbf{bactériennes} (syphilis, gonococcie, chlamydioses) et la trichomonose (parasite) se \textbf{guérissent complètement} par des antibiotiques (ou antiparasitaire), à condition que le traitement soit pris \textbf{intégralement} (respect de la durée) et que \textbf{le(s) partenaire(s) soit(ent) traité(s) simultanément} pour éviter le « ping-pong » (réinfection croisée).
\item \textbf{Contrôler} : Le VIH ne se guérit pas encore, mais les \cle{antirétroviraux} (ARV), pris \textbf{à vie} avec une observance stricte, bloquent la multiplication du virus. Une personne traitée correctement vit longtemps en bonne santé, et sa charge virale plasmatique devient \textbf{indétectable} (< 50 copies/mL), ce qui supprime pratiquement le risque de transmission sexuelle (concept \textbf{I = I : Indétectable = Intransmissible}). Le traitement d'une femme enceinte séropositive (débuté au plus tôt) réduit à moins de 2 \% le risque de transmission à son bébé.
\item \textbf{Prévenir par la vaccination} : L'hépatite B (vaccin inclus dans le Programme Élargi de Vaccination - PEV du nourrisson au Cameroun, 3 doses) et les infections à HPV (vaccin anti-HPV, 2 doses pour les filles de 9 à 14 ans, idéalement avant les premiers rapports) se \textbf{préviennent efficacement par la vaccination}. Le vaccin HPV protège contre les génotypes 16 et 18 (70 \% des cancers du col) et souvent 6 et 11 (condylomes).
\end{itemize}
\end{propriete}

\begin{exemplebox}[Le préservatif : une barrière efficace, la « double protection »]
Le \cle{préservatif} (masculin en latex ou polyuréthane, ou féminin en nitrile), utilisé \textbf{correctement et à chaque rapport}, du début à la fin, constitue une barrière physique étanche qui empêche le contact entre les muqueuses et les liquides biologiques infectants. Les études épidémiologiques (méta-analyses) montrent que son usage correct et systématique réduit d'environ \textbf{80 à 95 \%} le risque de transmission du VIH, et réduit fortement celui de la plupart des autres IST bactériennes (gonococcie, chlamydioses, trichomonose) et de l'hépatite B. Son efficacité est moindre pour les IST à transmission cutanée-muqueuse étendue (herpès, HPV, syphilis si lésion hors zone couverte) mais reste significative. Il protège \textbf{en même temps contre les grossesses non désirées} : c'est le seul moyen de « double protection ».
\textbf{Règles d'or :} vérifier la date de péremption et l'intégrité de l'emballage (marque CE/NF), l'ouvrir avec précaution (pas de dents, pas de ciseaux), le mettre avant tout contact génital, le retirer après l'éjaculation en tenant la base, ne \textbf{jamais} le réutiliser, ne pas utiliser deux préservatifs superposés (risque de rupture par frottement).
\end{exemplebox}

\begin{exoresolu}[Analyser une situation et conseiller un comportement responsable]
\textbf{Énoncé.} Arnaud, 18 ans, déclare à son ami : « Ma partenaire est en pleine forme, elle n'a aucun bouton ni écoulement, donc elle ne peut pas avoir d'IST. Inutile d'utiliser un préservatif. »
\begin{enumerate}
\item Relever les deux erreurs de raisonnement commises par Arnaud.
\item Expliquer pourquoi un dépistage peut être utile même en l'absence de symptômes.
\item Citer trois moyens concrets de se protéger des IST.
\end{enumerate}
\tcblower
\textbf{Solution détaillée.}
\begin{enumerate}
\item \textbf{Première erreur} : Arnaud croit qu'on reconnaît une personne infectée à son apparence clinique. Or de nombreuses IST sont \textbf{asymptomatiques} : la chlamydiose est silencieuse dans environ 70 \% des cas chez la femme, et le VIH peut rester cliniquement invisible pendant 8 à 10 ans (phase de latence). \textbf{Deuxième erreur} : il conclut de cette apparence saine qu'aucune protection n'est nécessaire, alors que c'est précisément quand on ne connaît pas le statut sérologique/bactériologique de son partenaire (ou de soi-même) que le préservatif est indispensable.
\item Le dépistage (sérologie, PCR, culture) est le \textbf{seul moyen fiable} de connaître son statut, car il détecte l'agent infectieux (antigène, ADN/ARN viral, bactérie) ou les anticorps spécifiques dans le sang ou les sécrétions, même sans aucun symptôme. Dépister tôt permet de traiter tôt : les IST bactériennes se guérissent alors avant d'entraîner des complications irréversibles (stérilité, grossesse extra-utérine), et une infection à VIH prise en charge précocement se contrôle efficacement grâce aux antirétroviraux, préservant le système immunitaire et évitant la transmission.
\item Trois moyens de protection concrets : (1) \textbf{Utiliser un préservatif (masculin ou féminin) correctement à chaque rapport sexuel} (double protection IST + grossesse). (2) \textbf{Pratiquer le dépistage régulier}, idéalement en couple avant d'envisager l'abandon du préservatif. (3) \textbf{Se faire vacciner} contre l'hépatite B (si non fait dans l'enfance) et contre les HPV (filles/garçons selon recommandations). On peut ajouter : la fidélité mutuelle entre partenaires dépistés négatifs, l'abstinence, la non-utilisation de drogues injectables / matériel stérile.
\end{enumerate}
\end{exoresolu}

\begin{aretenir}[L'essentiel de la section : IST et prévention]
\begin{itemize}
\item Une IST est une infection transmise principalement lors de rapports sexuels non protégés ; certaines passent aussi par le sang (VIH, VHB) ou de la mère à l'enfant (VIH, syphilis, VHB, gonocoque, chlamydiae, herpès).
\item Les IST \textbf{bactériennes} (gonococcie, syphilis, chlamydioses) et la trichomonose se \textbf{guérissent} par antibiotiques/antiparasitaires (traitement complet + partenaire) ; les IST \textbf{virales} (VIH, VHB, herpès, HPV) ne se guérissent pas (sauf VHB aigu) mais se \textbf{contrôlent} (ARV, antiviraux) ou se \textbf{préviennent par vaccin} (VHB, HPV).
\item \textbf{Beaucoup d'IST sont asymptomatiques} : l'apparence physique ne renseigne \textbf{jamais} sur le statut d'une personne ; seul le dépistage (gratuit, confidentiel au Cameroun) est fiable.
\item Non traitées, elles provoquent : stérilité, grossesses extra-utérines, cancers (col utérus \textleftarrow HPV), transmission mère-enfant, et pour le VIH : immunodépression sévère (SIDA).
\item Le VIH ne se transmet \textbf{ni} par les contacts quotidiens (main, repas, baiser, piscine), \textbf{ni} par les moustiques : il faut combattre la stigmatisation des PVVIH.
\item \textbf{Piliers de la protection} : Dépistage précoce et régulier $\rightarrow$ Traitement complet et précoce $\rightarrow$ Vaccination (VHB, HPV) $\rightarrow$ Préservatif correct et systématique (double protection) $\rightarrow$ Comportement responsable (fidélité mutuelle dépistée, réduction des partenaires, refus du partage d'objets tranchants).
\end{itemize}
\end{aretenir}

Maintenant que nous connaissons les IST, leurs dangers et les moyens de les dépister, de les traiter et de s'en protéger, la section suivante précisera les méthodes de \textbf{planification familiale} (notions générales) pour permettre à tout adolescent ou jeune de choisir, en toute connaissance de cause, s'il veut avoir des enfants, quand et combien, dans le respect de sa santé et de ses droits.

\coursec{Prévention des IST et comportement responsable}

Dans la section précédente, nous avons dressé l'inventaire des principales infections sexuellement transmissibles (IST) : le VIH/Sida, la syphilis, la gonococcie, la chlamydiose, l'hépatite B, les infections à papillomavirus (HPV), la trichomonose et l'herpès génital. Nous avons vu que certaines se soignent facilement quand elles sont détectées tôt, que d'autres (comme le VIH) ne se guérissent pas mais se contrôlent, et que toutes peuvent avoir des conséquences graves : stérilité, cancers, transmission à l'enfant, décès prématuré. Face à ce constat, une question s'impose naturellement : \emph{comment un adolescent ou une adolescente peut-il(elle) se protéger concrètement ?} C'est tout l'objet de cette section : transformer la connaissance des risques en \cle{comportements responsables} qui protègent la santé.

\courssub{Les trois piliers de la prévention : la stratégie A-F-P}

\textbf{Situation concrète.} Lors d'une causerie éducative organisée dans un lycée de Yaoundé, une infirmière pose la question aux élèves : « Que peut faire un jeune pour éviter les IST ? » Les réponses fusent : « ne pas avoir de rapports », « rester avec un seul partenaire », « utiliser un préservatif ». Sans le savoir, les élèves viennent d'énoncer les trois piliers reconnus mondialement de la prévention des IST, résumés par le sigle \cle{A-F-P}.

\begin{definition}[La stratégie A-F-P]
La stratégie \cle{A-F-P} est la démarche de prévention des IST fondée sur trois comportements complémentaires, classés par ordre d'efficacité :
\begin{itemize}
\item \textbf{A — Abstinence} : ne pas avoir de rapports sexuels. C'est la seule méthode efficace à 100 \% contre les IST transmises par voie sexuelle et contre les grossesses.
\item \textbf{F — Fidélité} : avoir un \cle{partenaire unique}, mutuellement fidèle, et dont on sait (par un dépistage) qu'il n'est pas infecté.
\item \textbf{P — Préservatif} : utiliser correctement et systématiquement un préservatif (masculin ou féminin) lors de tout rapport sexuel avec un partenaire dont le statut infectieux est inconnu ou incertain.
\end{itemize}
\end{definition}

\begin{attention}[La fidélité ne protège que si elle est partagée et vérifiée]
Être fidèle soi-même ne suffit pas : la fidélité ne protège que si \emph{les deux} partenaires sont fidèles \emph{et} non infectés. Or certaines IST (VIH, chlamydiose, HPV) peuvent être présentes sans aucun signe visible pendant des années. D'où l'importance du \cle{dépistage en couple} avant de renoncer au préservatif. La confiance ne remplace pas un test de laboratoire.
\end{attention}

\begin{popfigure}
\begin{center}
\begin{tikzpicture}
% Pyramide A-F-P centrée sur x=0 pour un meilleur centrage dans la figure
\fill[popgreenL] (-3,0) -- (3,0) -- (2.5,1.8) -- (-2.5,1.8) -- cycle;
\fill[poporangeL] (-2.5,1.8) -- (2.5,1.8) -- (2,3.6) -- (-2,3.6) -- cycle;
\fill[popblueL] (-2,3.6) -- (2,3.6) -- (0,5.2) -- cycle;
\draw[popdark,thick] (-3,0) -- (3,0) -- (0,5.2) -- cycle;
\draw[popdark] (-2.5,1.8) -- (2.5,1.8);
\draw[popdark] (-2,3.6) -- (2,3.6);
\node at (0,0.9) {\textbf{P — Préservatif}};
\node at (0,2.7) {\textbf{F — Fidélité}};
\node at (0,4.15) {\textbf{A — Abstinence}};
% Annotations latérales (à droite)
\draw[popmuted,thin] (3,0.9) -- (4.3,0.9) node[right,align=left,text width=5cm] {\small Protection efficace si usage \textbf{correct et systématique} ; protège aussi des grossesses.};
\draw[popmuted,thin] (2.5,2.7) -- (4.3,2.7) node[right,align=left,text width=5cm] {\small Partenaire unique, mutuellement fidèle, \textbf{dépisté non infecté}.};
\draw[popmuted,thin] (2,4.4) -- (4.3,4.4) node[right,align=left,text width=5cm] {\small Seule méthode efficace à \textbf{100 \%} contre IST et grossesses.};
% Flèche efficacité (à gauche)
\draw[-{Stealth},very thick,poppurple] (-4.2,0.3) -- (-4.2,4.9) node[midway,left,rotate=90,align=center] {\small efficacité croissante};
\end{tikzpicture}
\end{center}
\legende{Schéma pyramidal de la stratégie A-F-P. L'abstinence (sommet) est la méthode la plus efficace ; la fidélité vérifiée par un dépistage vient ensuite ; le préservatif (base) est le moyen de protection indispensable lorsque les deux premières conditions ne sont pas réunies.}
\end{popfigure}

\courssub{Le préservatif : une barrière mécanique à double protection}

\textbf{Observation.} Au laboratoire, on peut réaliser une expérience simple : on remplit un préservatif intact avec de l'eau colorée contenant des particules fines, on le ferme et on le presse. Aucune particule ne traverse la membrane. Cette observation illustre le principe même de la protection.

\begin{propriete}[Efficacité du préservatif (admise)]
L'usage \cle{correct} et \cle{systématique} du préservatif réduit très fortement le risque de transmission des IST (y compris le VIH) et le risque de grossesse. Cette propriété, admise au niveau de la Première, est fondée sur les données épidémiologiques : les études de suivi de couples dont un seul partenaire est porteur du VIH montrent que l'utilisation systématique du préservatif réduit le risque de transmission de plus de 90 \%.
\end{propriete}

\textbf{Justification scientifique (mécanisme).} Le préservatif est une gaine fine et étanche (en latex ou en polyuréthane) qui recouvre le pénis en érection (préservatif masculin) ou qui tapisse le vagin (préservatif féminin). Il constitue une \cle{barrière mécanique} qui empêche tout contact direct entre :
\begin{itemize}
\item les sécrétions génitales (sperme, sécrétions vaginales) des deux partenaires ;
\item le sang et les lésions éventuelles des muqueuses génitales.
\end{itemize}
Or ce sont précisément ces liquides biologiques qui transportent les agents infectieux (virus du VIH, bactéries de la gonococcie ou de la chlamydiose, tréponème de la syphilis, HPV\ldots). En bloquant leur passage, la barrière mécanique interrompt la \cle{chaîne de transmission}. Le préservatif est ainsi le \textbf{seul moyen qui protège à la fois des IST et des grossesses}, puisqu'il retient aussi les spermatozoïdes.

\begin{methode}[Conditions d'efficacité du préservatif]
Pour que la protection soit réelle, il faut un usage \textbf{correct} et \textbf{systématique} :
\begin{enumerate}
\item Vérifier la date de péremption et l'intégrité de l'emballage (le sachet doit contenir de l'air).
\item Ouvrir le sachet avec les doigts, jamais avec les dents ou un objet coupant.
\item Dérouler le préservatif sur le pénis en érection \emph{avant tout contact} génital, en pinçant le réservoir pour chasser l'air.
\item Utiliser un préservatif \emph{neuf à chaque rapport} ; ne jamais en utiliser deux à la fois (le frottement les déchire).
\item Le retirer immédiatement après l'éjaculation en le tenant à la base, le nouer et le jeter dans une poubelle (jamais dans les toilettes).
\item Conserver les préservatifs à l'abri de la chaleur et de la lumière (pas dans le portefeuille ni dans la boîte à gants d'une voiture).
\end{enumerate}
Le préservatif féminin suit des principes analogues : il s'insère dans le vagin avant le rapport et offre à la femme un moyen de protection qu'elle contrôle elle-même.
\end{methode}

\begin{attention}[Erreur fréquente : la pilule ne protège pas des IST]
Beaucoup de jeunes croient qu'une fille qui prend la \cle{pilule contraceptive} est « protégée ». C'est faux et dangereux : la pilule agit sur les hormones pour bloquer l'ovulation ; elle protège \textbf{uniquement de la grossesse} et n'offre \emph{aucune} barrière contre les agents infectieux. De même, ni le stérilet, ni l'implant, ni les injectables ne protègent des IST. Seul le préservatif assure la double protection.
\end{attention}

\courssub{Le dépistage volontaire : connaître son statut}

\textbf{Situation concrète.} Chaque année, lors de la Journée mondiale de lutte contre le Sida (1\textsuperscript{er} décembre), des stands de dépistage gratuit sont installés dans les quartiers, les marchés et les campus universitaires du Cameroun. En quelques minutes, une goutte de sang prélevée au bout du doigt permet de connaître son statut vis-à-vis du VIH.

\begin{definition}[Dépistage volontaire]
Le \cle{dépistage volontaire} est un examen médical (test sanguin ou analyse de sécrétions) réalisé à la demande de la personne, pour rechercher la présence d'un agent infectieux. Il est \textbf{gratuit} (pour le VIH) et \textbf{confidentiel} dans les formations sanitaires publiques et de nombreux centres associatifs au Cameroun.
\end{definition}

L'intérêt du dépistage est triple :
\begin{itemize}
\item \textbf{Pour soi} : une IST détectée tôt se traite tôt, avant les complications (stérilité, cancer, Sida déclaré). Une personne séropositive mise sous traitement antirétroviral peut vivre longtemps et, si sa charge virale devient indétectable, ne transmet plus le virus.
\item \textbf{Pour le couple} : le \cle{dépistage en couple}, avant le mariage ou avant d'arrêter le préservatif, permet de fonder la fidélité sur une certitude médicale et non sur une simple confiance.
\item \textbf{Pour la communauté} : connaître son statut brise la chaîne de transmission.
\end{itemize}

\begin{savaistu}[Le dépistage du VIH est gratuit au Cameroun]
Au Cameroun, le test de dépistage du VIH est \textbf{gratuit} dans les hôpitaux publics, les centres de santé intégrés et de nombreuses associations de jeunes (clubs santé, ONG de lutte contre le Sida). Des campagnes régulières proposent même le dépistage à domicile ou sur les lieux de rassemblement des jeunes. Le résultat est strictement \textbf{confidentiel} : nul ne peut être dépisté de force, et le secret médical est protégé par la loi.
\end{savaistu}

\courssub{La vaccination : une prévention spécifique}

Certaines IST peuvent être prévenues par la \cle{vaccination}, qui apprend à l'organisme à fabriquer des anticorps avant tout contact avec l'agent infectieux :

\begin{itemize}
\item \textbf{Vaccin contre l'hépatite B} : le virus de l'hépatite B (VHB) se transmet par le sang et les rapports sexuels ; il peut provoquer une cirrhose ou un cancer du foie. La vaccination (schéma en plusieurs doses) confère une protection durable. Elle fait partie des vaccins administrés aux nourrissons au Cameroun, et un rattrapage est possible à tout âge.
\item \textbf{Vaccin contre le papillomavirus (HPV)} : les HPV sont responsables de la quasi-totalité des \cle{cancers du col de l'utérus}, première cause de mortalité par cancer chez la femme en Afrique subsaharienne.
\end{itemize}

\begin{exemplebox}[Exemple chiffré : le vaccin anti-HPV]
Le cancer du col de l'utérus est causé dans plus de 95 \% des cas par une infection persistante à certains types de HPV. Or la vaccination anti-HPV, administrée \emph{avant les premiers rapports sexuels} (idéalement entre 9 et 14 ans), prévient l'infection par les types de HPV les plus cancérigènes et peut ainsi prévenir la \textbf{grande majorité des cancers du col de l'utérus}. Le Cameroun a intégré ce vaccin dans son Programme élargi de vaccination, gratuit pour les filles de la tranche d'âge cible. C'est un exemple remarquable où un simple vaccin évite un cancer.
\end{exemplebox}

\courssub{Hygiène et précautions face au sang}

Certaines IST (VIH, hépatites B et C) se transmettent aussi par le \cle{sang}. Le comportement responsable impose donc de :
\begin{itemize}
\item ne \textbf{jamais partager} d'objets tranchants ou piquants susceptibles d'être souillés de sang : lames de rasoir, aiguilles de couture ou d'injection, outils de scarification ou de tatouage, coupe-ongles, brosses à dents ;
\item exiger, lors de soins, de circoncisions, de tatouages ou de perçages, du \textbf{matériel stérile à usage unique} ;
\item se protéger les mains (gants, sachet plastique) pour secourir une personne blessée qui saigne ;
\item vérifier, en cas de transfusion, que le sang a été testé (c'est une obligation dans les formations sanitaires).
\end{itemize}

\begin{attention}[Précision importante]
Le VIH ne se transmet \textbf{pas} par les poignées de main, les embrassades, les piqûres de moustiques, les repas partagés, les toilettes ou les larmes. Ces fausses croyances alimentent la discrimination envers les personnes séropositives ; les combattre fait aussi partie du comportement responsable.
\end{attention}

\courssub{En cas de signes : consulter tôt, traiter le couple}

Lorsque des signes évocateurs apparaissent (écoulement anormal, brûlures en urinant, boutons ou ulcérations génitales, douleurs du bas-ventre), la conduite responsable est claire :
\begin{enumerate}
\item \textbf{Consulter sans honte et sans délai} dans une formation sanitaire : la plupart des IST bactériennes (gonococcie, chlamydiose, syphilis) se guérissent complètement avec des antibiotiques bien suivis.
\item \textbf{Ne jamais s'automédiquer} ni acheter des médicaments au marché : des doses insuffisantes rendent les bactéries résistantes et masquent la maladie qui continue d'évoluer.
\item \textbf{Traiter aussi le ou les partenaires} : soigner une seule personne du couple conduit à une « réinfection en ping-pong ». Le personnel de santé remet souvent un traitement pour le partenaire.
\item \textbf{S'abstenir de tout rapport} (ou utiliser un préservatif) jusqu'à la guérison confirmée.
\end{enumerate}

\courssub{Adopter un comportement responsable}

Au-delà des moyens techniques, la prévention des IST repose sur des attitudes :

\begin{itemize}
\item \textbf{S'informer} auprès de sources fiables (personnel de santé, enseignants, centres de jeunes) et non auprès de rumeurs ;
\item \textbf{Dialoguer} avec son partenaire : parler de prévention, de dépistage et de fidélité avant tout engagement est une preuve de respect, pas de méfiance ;
\item \textbf{Savoir dire non} : refuser un rapport non protégé, refuser tout rapport forcé. Aucun rapport sexuel ne doit être subi ; le consentement libre et éclairé est un droit ;
\item \textbf{Se respecter et respecter autrui} : ne pas exercer de pression sur un(e) partenaire, ne pas stigmatiser les personnes vivant avec le VIH, assumer la responsabilité de ses actes.
\end{itemize}

\begin{methode}[Justifier l'utilisation d'un moyen de prévention (méthode de rédaction)]
Au Probatoire, pour justifier un moyen de prévention, rédige en trois étapes :
\begin{enumerate}
\item \textbf{Nommer le risque} : identifier l'IST ou la grossesse non désirée contre laquelle on se protège, et le liquide biologique vecteur (sperme, sécrétions vaginales, sang).
\item \textbf{Expliquer le mécanisme de protection} : montrer comment le moyen proposé interrompt la transmission (barrière mécanique pour le préservatif, anticorps pour le vaccin, absence d'exposition pour l'abstinence).
\item \textbf{Conclure sur le bénéfice} : préciser l'avantage concret (pas de transmission, pas de grossesse, double protection, gratuité, accessibilité).
\end{enumerate}
\end{methode}

\begin{exoresolu}[Justifier un conseil de prévention]
\textbf{Énoncé.} Ton ami de 17 ans te dit : « Ma copine prend la pilule, donc nous n'avons pas besoin de préservatif. » En trois arguments rédigés, explique-lui pourquoi ce raisonnement est dangereux et quel comportement tu lui conseilles.
\tcblower
\textbf{Solution détaillée.}
\begin{enumerate}
\item \emph{Nommer le risque} : la pilule contraceptive bloque l'ovulation et protège uniquement de la grossesse ; elle n'offre aucune protection contre les IST (VIH, chlamydiose, gonococcie, HPV\ldots), qui se transmettent par les sécrétions génitales et le sang.
\item \emph{Expliquer le mécanisme} : seul le préservatif constitue une barrière mécanique qui empêche le contact entre les sécrétions génitales des deux partenaires ; il interrompt ainsi la chaîne de transmission des agents infectieux tout en retenant les spermatozoïdes.
\item \emph{Conclure sur le bénéfice} : je lui conseille donc d'utiliser le préservatif correctement et à chaque rapport (double protection IST + grossesse), et, s'ils souhaitent y renoncer un jour, de faire d'abord un dépistage en couple dans un centre de santé, où le test du VIH est gratuit et confidentiel.
\end{enumerate}
\end{exoresolu}

\begin{popfigure}
\begin{center}
\begin{tikzpicture}[
case/.style={draw=popdark,fill=popcream,rounded corners=2pt,minimum width=5.6cm,minimum height=2.6cm},
bulle/.style={draw=popblue,fill=popblueL,rounded corners=6pt,align=center,text width=4.6cm,font=\small}]
% Case 1
\node[case] (c1) at (0,0) {};
\node[font=\small\bfseries,text=popdark] at (0,1.05) {Case 1 : à compléter};
\node[bulle] at (0,-0.15) {« Viens, on le fait sans préservatif, fais-moi confiance\ldots »};
% Case 2
\node[case] (c2) at (6.4,0) {};
\node[font=\small\bfseries,text=popdark] at (6.4,1.05) {Case 2 : à compléter};
\node[bulle,draw=poporange,fill=poporangeL, align=center] at (6.4,-0.15){« \ldots\ldots\ldots\ldots\ldots\ldots \\ \ldots\ldots\ldots\ldots\ldots\ldots »};
% Case 3
\node[case] (c3) at (0,-3.6) {};
\node[font=\small\bfseries,text=popdark] at (0,-2.55) {Case 3 : à compléter};
\node[bulle, align=center] at (0,-3.75){« \ldots\ldots\ldots\ldots\ldots\ldots \\ \ldots\ldots\ldots\ldots\ldots\ldots »};
% Case 4
\node[case] (c4) at (6.4,-3.6) {};
\node[font=\small\bfseries,text=popdark] at (6.4,-2.55) {Case 4 : à compléter};
\node[bulle,draw=popgreen,fill=popgreenL] at (6.4,-3.75) {« D'accord, allons d'abord nous faire dépister ensemble, c'est gratuit au centre de santé. »};
% Flèches de lecture
\draw[-{Stealth},thick,popmuted] (3.1,0) -- (3.3,0);
\draw[-{Stealth},thick,popmuted] (6.4,-1.6) -- (6.4,-2.0);
\draw[-{Stealth},thick,popmuted] (3.3,-3.6) -- (3.1,-3.6);
\end{tikzpicture}
\end{center}
\legende{Bande dessinée muette à compléter : deux adolescents discutent d'un rapport non protégé. Rédige les répliques des cases 2 et 3 en appliquant les trois étapes de la méthode (nommer le risque, expliquer le mécanisme, conclure sur le bénéfice) : la case 2 doit contenir un refus argumenté, la case 3 une proposition de solution (préservatif ou dépistage en couple).}
\end{popfigure}

\begin{exercice}[À toi de jouer]
\begin{enumerate}
\item Complète la bande dessinée ci-dessus en rédigeant les répliques des cases 2 et 3.
\item Cite les trois piliers de la stratégie A-F-P et classe-les par ordre d'efficacité.
\item Explique, en trois phrases, pourquoi le préservatif est dit « à double protection ».
\item Une camarade affirme : « Je suis vaccinée contre l'hépatite B, je ne risque plus aucune IST. » Corrige cette affirmation.
\end{enumerate}
\end{exercice}

\begin{corrige}[Exercice : À toi de jouer]
\begin{enumerate}
\item \emph{Exemple de réponse attendue.} Case 2 : « Non. Sans préservatif, je risque une IST comme le VIH ou une grossesse : nos sécrétions seraient en contact direct, et la confiance ne remplace pas un test. » Case 3 : « Soit on utilise un préservatif à chaque rapport, soit on attend et on va ensemble au centre de santé pour un dépistage gratuit et confidentiel. »
\item A — Abstinence (efficace à 100 \%) ; F — Fidélité à un partenaire unique non infecté (vérifiée par un dépistage) ; P — Préservatif (usage correct et systématique).
\item Le préservatif est une barrière mécanique qui empêche le contact entre les sécrétions génitales : il bloque donc le passage des agents infectieux (protection contre les IST) \emph{et} retient les spermatozoïdes (protection contre la grossesse). C'est le seul moyen offrant ces deux protections simultanément.
\item Le vaccin contre l'hépatite B ne protège que contre le virus de l'hépatite B. Il n'a aucune action contre le VIH, la gonococcie, la syphilis, la chlamydiose ou le HPV. La vaccination est une prévention \emph{spécifique} qui doit être complétée par la stratégie A-F-P.
\end{enumerate}
\end{corrige}

\begin{aretenir}[L'essentiel de la section]
\begin{itemize}
\item La prévention des IST repose sur la stratégie \cle{A-F-P} : Abstinence, Fidélité vérifiée, Préservatif.
\item Le \cle{préservatif} (masculin ou féminin) est une barrière mécanique : c'est le seul moyen qui protège à la fois des IST et des grossesses, à condition d'un usage correct et systématique.
\item La pilule et les autres contraceptifs hormonaux ne protègent \textbf{pas} des IST.
\item Le \cle{dépistage volontaire} (gratuit et confidentiel pour le VIH au Cameroun) permet de connaître son statut, idéalement en couple.
\item La \cle{vaccination} contre l'hépatite B et le HPV est une prévention spécifique ; le vaccin anti-HPV avant les premiers rapports prévient la majorité des cancers du col de l'utérus.
\item Ne jamais partager d'objets souillés de sang ; consulter tôt en cas de signes et traiter aussi le ou les partenaires.
\item Le comportement responsable : s'informer, dialoguer, refuser les rapports forcés ou non protégés, respecter soi-même et autrui.
\end{itemize}
\end{aretenir}

\coursec{Planification familiale (notions générales)}

Dans la section précédente, nous avons vu que le préservatif est l'outil central de la prévention des IST et que le comportement responsable repose sur l'information, le dépistage et la fidélité. Mais la santé reproductive ne se limite pas à éviter les infections : elle concerne aussi la maîtrise de la fécondité. Observer la vie d'un quartier de Yaoundé ou d'un village de l'Extrême-Nord suffit pour poser le problème : deux familles voisines, l'une compte trois enfants espacés de trois ans, l'autre sept enfants nés presque chaque année. Dans la seconde, la mère est épuisée, les aînés ont quitté l'école, et les ressources du foyer sont étirées au maximum. Comment un couple peut-il choisir le nombre de ses enfants et le moment de leur naissance ? C'est toute la question de la \cle{planification familiale}.

\courssub{Définitions et intérêts de la planification familiale}

\begin{definition}[Planification familiale et contraception]
La \cle{planification familiale} (ou planning familial) est l'ensemble des moyens qui permettent à un couple de \textbf{choisir le nombre de ses enfants} et \textbf{l'espacement des naissances}, afin d'avoir les enfants qu'il désire, au moment où il le désire.

La \cle{contraception} désigne l'ensemble des \textbf{méthodes} utilisées pour empêcher temporairement ou définitivement la fécondation (rencontre du spermatozoïde et de l'ovule) ou la nidation de l'œuf. On parle aussi de méthodes \cle{contraceptives}.
\end{definition}

\begin{attention}[Ne pas confondre contraception et prévention des IST]
Erreur très fréquente au Probatoire : croire qu'une méthode contraceptive protège automatiquement des infections sexuellement transmissibles. C'est \textbf{faux}. La pilule, l'injectable, l'implant, le stérilet empêchent la grossesse mais \textbf{ne protègent pas du VIH ni des autres IST}. Seul le \cle{préservatif} (masculin ou féminin) assure la \textbf{double protection} : contre la grossesse \emph{et} contre les IST. Un adolescent ou un jeune couple exposé au risque d'IST doit donc utiliser le préservatif, éventuellement associé à une autre méthode contraceptive.
\end{attention}

\begin{aretenir}[Intérêts de la planification familiale]
\begin{itemize}
\item \textbf{Santé de la mère} : les grossesses trop rapprochées, trop précoces (avant 18 ans) ou trop nombreuses épuisent l'organisme et augmentent les risques d'hémorragies, d'anémie et de décès maternel.
\item \textbf{Santé des enfants} : un enfant né au moins deux ans après le précédent bénéficie d'un allaitement plus long et de meilleurs soins ; sa survie est meilleure.
\item \textbf{Scolarisation} : espacer les naissances permet aux filles et aux garçons de poursuivre leurs études ; la mère elle-même peut achever sa formation.
\item \textbf{Équilibre économique du foyer} : les ressources familiales (nourriture, frais de scolarité, soins) sont mieux réparties entre les enfants.
\end{itemize}
\end{aretenir}

\begin{savaistu}[Espacer les naissances sauve des vies]
Les données de santé publique, notamment les enquêtes démographiques et de santé menées au Cameroun (EDS), montrent que les enfants nés moins de deux ans après leur aîné présentent un risque de mortalité infantile nettement plus élevé. Un \textbf{espacement des naissances d'au moins deux ans} réduit sensiblement la mortalité maternelle et infantile. C'est pourquoi les agents de santé des centres de santé intégrés et les associations comme le planning familial sensibilisent les couples, y compris dans les zones rurales.
\end{savaistu}

\courssub{Les grandes familles de méthodes contraceptives}

\courssub{Méthodes naturelles}

Les méthodes dites \cle{naturelles} n'utilisent ni produit ni matériel : elles reposent sur la connaissance du cycle féminin ou sur le comportement.

\begin{itemize}
\item \textbf{Abstinence périodique (méthode du calendrier ou d'Ogino)} : le couple évite les rapports sexuels pendant la période féconde du cycle, estimée autour de l'ovulation (vers le 14\textsuperscript{e} jour d'un cycle de 28 jours). \emph{Limite majeure} : la date d'ovulation varie d'un cycle à l'autre, surtout chez l'adolescente dont les cycles sont irréguliers ; la méthode est donc \textbf{peu fiable}.
\item \textbf{Le retrait (coït interrompu)} : l'homme retire son pénis avant l'éjaculation. \emph{Limite majeure} : des spermatozoïdes peuvent être libérés avant l'éjaculation (liquide pré-éjaculatoire), et le geste exige un contrôle difficile ; la méthode est \textbf{peu fiable}.
\end{itemize}

\begin{attention}[Calendrier et retrait : déconseillés aux adolescents]
La méthode du calendrier et le retrait entraînent de nombreux échecs, donc des grossesses non désirées. Ils sont \textbf{déconseillés aux adolescents}, dont les cycles sont irréguliers et l'expérience limitée. Pour un adolescent, la seule conduite totalement sûre reste l'\cle{abstinence} ; à défaut, le préservatif correctement utilisé.
\end{attention}

\courssub{Méthodes de barrière}

Les méthodes de \cle{barrière} empêchent physiquement la rencontre entre les spermatozoïdes et l'ovule.

\begin{itemize}
\item \textbf{Préservatif masculin} : gaine fine (latex) placée sur le pénis en érection, qui recueille le sperme. Bien utilisé à chaque rapport, il est efficace et constitue la \textbf{seule méthode protégeant à la fois de la grossesse et des IST} (dont le VIH).
\item \textbf{Préservatif féminin} : gaine placée dans le vagin avant le rapport ; il offre la même double protection et peut être mis en place par la femme elle-même.
\end{itemize}

\courssub{Méthodes hormonales}

Les méthodes \cle{hormonales} apportent des hormones (analogues des œstrogènes et/ou de la progestérone) qui \textbf{bloquent l'ovulation} : sans ovule libéré, il ne peut y avoir de fécondation. Elles épaississent aussi la glaire cervicale, ce qui freine le passage des spermatozoïdes.

\begin{itemize}
\item \textbf{La pilule} : comprimé à prendre chaque jour, à heure régulière. Très efficace si elle est prise sans oubli ; un oubli réduit fortement l'efficacité.
\item \textbf{Les injectables} : injection d'hormones par un agent de santé, dont l'effet dure un à trois mois selon le produit.
\item \textbf{Les implants} : petites tiges souples placées sous la peau du bras par un personnel formé ; elles diffusent l'hormone pendant trois à cinq ans.
\end{itemize}

Ces méthodes sont \textbf{très efficaces contre la grossesse}, mais rappelons-le : \textbf{aucune ne protège des IST}.

\courssub{Méthodes de longue durée et méthodes définitives}

\begin{itemize}
\item \textbf{Le stérilet ou DIU} (dispositif intra-utérin) : petit objet en plastique et cuivre (ou diffusant une hormone) placé dans l'utérus par un personnel de santé qualifié. Il empêche la fécondation et/ou la nidation pendant plusieurs années (5 à 10 ans selon le modèle). C'est une méthode \textbf{réversible de longue durée} : son retrait rend la fertilité.
\item \textbf{La ligature des trompes} : intervention chirurgicale qui sectionne ou obstrue les trompes de Fallope chez la femme ; l'ovule ne peut plus rencontrer les spermatozoïdes. Méthode \textbf{définitive}.
\item \textbf{La vasectomie} : section des canaux déférents chez l'homme ; le sperme ne contient plus de spermatozoïdes. Méthode \textbf{définitive}.
\end{itemize}

Les méthodes définitives sont réservées aux \textbf{adultes ayant complété leur famille} et ayant mûrement réfléchi à leur choix, car elles sont irréversibles en pratique.

\courssub{Classification et efficacité comparée des méthodes}

\begin{popfigure}
\begin{center}
\begin{tikzpicture}[
  grow=right, level distance=3.8cm,
  level 1/.style={sibling distance=3.0cm},
  level 2/.style={sibling distance=1.2cm, level distance=3.5cm},
  every node/.style={font=\small},
  racine/.style={draw=popdark, fill=popblueL, rounded corners, thick, align=center, text width=2.8cm},
  cat/.style={draw=popblue, fill=popblue!15, rounded corners, thick, align=center, text width=2.4cm},
  met/.style={draw=popmuted, fill=popcream, rounded corners, align=center, text width=2.9cm, font=\footnotesize},
  edge from parent/.style={draw=popmuted, thick}
]
\node[racine]{Méthodes contraceptives}
  child { node[cat]{Définitives}
    child { node[met]{Vasectomie (homme)} }
    child { node[met]{Ligature des trompes (femme)} }
  }
  child { node[cat]{Hormonales}
    child { node[met]{Implants (3--5 ans)} }
    child { node[met]{Injectables (1--3 mois)} }
    child { node[met]{Pilule (quotidienne)} }
  }
  child { node[cat]{Barrières}
    child { node[met]{Préservatif féminin} }
    child { node[met]{Préservatif masculin} }
  }
  child { node[cat]{Naturelles}
    child { node[met]{Retrait (peu fiable)} }
    child { node[met]{Calendrier (peu fiable)} }
  }
  child { node[cat, draw=popgreen, fill=popgreenL]{Longue durée (réversible)}
    child { node[met]{Stérilet / DIU (5--10 ans)} }
  };
\end{tikzpicture}
\end{center}
\legende{Arbre de classification des méthodes contraceptives. Cinq grandes familles : naturelles, barrières, hormonales, longue durée réversible, définitives. Seuls les préservatifs (barrières) protègent aussi des IST.}
\end{popfigure}

\begin{center}
\rowcolors{1}{popcream}{white}
\begin{tabularx}{\linewidth}{|l|X|c|c|}
\hline
\rowcolor{popblueL} \textbf{Méthode} & \textbf{Principe} & \textbf{Efficacité indicative} & \textbf{Protège des IST} \\
\hline
Calendrier (naturelle) & Éviter les rapports pendant la période féconde & Faible & Non \\
\hline
Retrait (naturelle) & Retirer le pénis avant l'éjaculation & Faible & Non \\
\hline
Préservatif masculin (barrière) & Gaine retenant le sperme & Bonne si usage correct à chaque rapport & \textbf{Oui} \\
\hline
Préservatif féminin (barrière) & Gaine tapissant le vagin & Bonne si usage correct & \textbf{Oui} \\
\hline
Pilule (hormonale) & Bloque l'ovulation & Très bonne si prise sans oubli & Non \\
\hline
Injectable (hormonal) & Bloque l'ovulation (1--3 mois) & Très bonne & Non \\
\hline
Implant (hormonal) & Bloque l'ovulation (3--5 ans) & Excellente & Non \\
\hline
Stérilet / DIU (longue durée) & Empêche fécondation et/ou nidation & Excellente (5--10 ans) & Non \\
\hline
Ligature des trompes (définitive) & Obstrue les trompes de Fallope & Totale (irréversible) & Non \\
\hline
Vasectomie (définitive) & Sectionne les canaux déférents & Totale (irréversible) & Non \\
\hline
\end{tabularx}
\end{center}

\begin{aretenir}[Notion d'efficacité]
L'\cle{efficacité} d'une méthode mesure sa capacité à empêcher la grossesse en usage réel. Elle dépend de deux choses : la qualité intrinsèque de la méthode \emph{et} la rigueur de son utilisation. Une pilule oubliée ou un préservatif mal mis perdent une partie de leur efficacité. Règle d'or à retenir : \textbf{aucune méthode hormonale ou naturelle ne protège des IST} ; seul le préservatif assure la double protection.
\end{aretenir}

\courssub{Où s'informer et comment choisir}

Le choix d'une méthode ne se fait pas seul, sur conseil d'un ami ou d'une rumeur. Il se fait :

\begin{itemize}
\item auprès du \textbf{personnel de santé qualifié} : médecins, infirmiers, sages-femmes des hôpitaux et des centres de santé intégrés ;
\item dans les services de \textbf{planning familial}, qui offrent information, conseil et méthodes adaptées, souvent à faible coût ;
\item dans le cadre du \textbf{dialogue au sein du couple} : choisir ensemble le nombre d'enfants et la méthode renforce la responsabilité partagée et le respect mutuel.
\end{itemize}

Pour un adolescent, le message de santé publique est clair : l'\textbf{abstinence} reste la protection la plus sûre ; si un jeune est exposé, le \textbf{préservatif} correctement utilisé à chaque rapport est indispensable, et toute question doit être posée à un agent de santé, sans honte ni crainte.

\begin{methode}[Répondre à une question de type Probatoire sur une méthode contraceptive]
Face à une situation (par exemple : « Un couple désire espacer ses naissances de trois ans ; proposez une méthode et justifiez »), procède en quatre étapes :
\begin{enumerate}
\item \textbf{Identifier} la méthode proposée ou choisie et la \textbf{classer} : naturelle, barrière, hormonale, longue durée ou définitive.
\item \textbf{Donner son principe} : bloquer l'ovulation, empêcher la rencontre des gamètes, empêcher la nidation, etc.
\item \textbf{Citer un avantage} adapté à la situation (efficacité, durée, double protection, réversibilité).
\item \textbf{Citer une limite} : oublis possibles, absence de protection contre les IST, irréversibilité, faible fiabilité.
\end{enumerate}
Une réponse complète contient toujours ces quatre éléments : classe, principe, avantage, limite.
\end{methode}

\begin{exoresolu}[Choisir et justifier une méthode]
\textbf{Énoncé.} Marie, 24 ans, mariée, a deux enfants de 4 ans et 18 mois. Son mari et elle souhaitent attendre environ trois ans avant une troisième grossesse. Ils sont fidèles l'un à l'autre et dépistés négatifs au VIH. (a) Proposez deux méthodes adaptées à leur situation en précisant leur classe. (b) Justifiez chaque choix par un avantage et une limite. (c) Pourquoi le retrait serait-il un mauvais choix ?
\tcblower
\textbf{Solution.}
(a) \emph{Premier choix : l'implant} (méthode hormonale de longue durée). \emph{Deuxième choix : le stérilet (DIU)} (méthode réversible de longue durée).

(b) \emph{Implant} : principe — il diffuse une hormone qui bloque l'ovulation pendant 3 à 5 ans. Avantage — très efficace, couvre toute la période de trois ans souhaitée sans geste quotidien, et réversible (le retrait rend la fertilité). Limite — il ne protège pas des IST (peu important ici car le couple est fidèle et dépisté) et nécessite une pose par un personnel formé. \emph{Stérilet} : principe — placé dans l'utérus, il empêche la fécondation et/ou la nidation pendant 5 à 10 ans. Avantage — efficacité excellente, aucune contrainte quotidienne, réversible. Limite — pose et retrait par un personnel de santé qualifié, aucune protection contre les IST.

(c) Le retrait est une méthode naturelle \textbf{peu fiable} : des spermatozoïdes peuvent être émis avant l'éjaculation (liquide pré-éjaculatoire), et le contrôle du geste est difficile. Le risque de grossesse non désirée avant le délai de trois ans serait donc élevé. Il est inadapté à un projet d'espacement précis des naissances.
\end{exoresolu}

\begin{exercice}[À toi de jouer]
Un couple de jeunes adultes, non dépistés, souhaite à la fois éviter une grossesse et se protéger des IST. (1) Quelle méthode conseilles-tu et pourquoi ? (2) Classe cette méthode et donne son principe. (3) Explique pourquoi la pilule seule ne répondrait pas à leur besoin.
\end{exercice}

\begin{corrige}[Exercice : double protection]
(1) Le \textbf{préservatif} (masculin ou féminin), utilisé correctement à chaque rapport, car c'est la seule méthode assurant la \textbf{double protection} : grossesse et IST. (2) C'est une méthode de \textbf{barrière} : la gaine retient le sperme et empêche la rencontre des gamètes, tout en bloquant le contact avec les sécrétions porteuses d'agents infectieux. (3) La pilule bloque bien l'ovulation (efficace contre la grossesse) mais \textbf{ne protège pas des IST} ; or le couple n'est pas dépisté, donc le risque infectieux demeure. La pilule seule ne répond donc qu'à la moitié du besoin.
\end{corrige}

\begin{aretenir}[L'essentiel de la section]
\begin{itemize}
\item La planification familiale permet de choisir le nombre d'enfants et l'espacement des naissances ; elle protège la santé de la mère et des enfants, favorise la scolarisation et l'équilibre du foyer.
\item Cinq familles de méthodes : naturelles (peu fiables), barrières (préservatifs, seule double protection), hormonales (pilule, injectables, implants), longue durée réversible (stérilet/DIU) et définitives (ligature, vasectomie).
\item Aucune méthode hormonale ou naturelle ne protège des IST.
\item On s'informe auprès du personnel de santé qualifié et des services de planning familial, dans le dialogue du couple.
\end{itemize}
\end{aretenir}

\newpage

\coursec{Activité d'intégration}

\coursec{Amélioration de la santé reproductive}

\courssub{Activité d'intégration : Le Conseil municipal des jeunes de Nkol-Efoulan}

\textbf{Objectifs de l'activité}\par

À l'issue de cette activité d'intégration, tu seras capable de :

\begin{itemize}
\item analyser des données chiffrées de santé reproductive issues d'un contexte camerounais réaliste ;
\item identifier et hiérarchiser les \cle{causes} individuelles, sociales et économiques des grossesses précoces et des IST ;
\item construire un tableau argumenté des \cle{conséquences} des grossesses précoces pour la mère, l'enfant et la communauté ;
\item \cle{justifier scientifiquement} l'utilisation des moyens de prévention des IST et des méthodes de planification familiale ;
\item proposer et défendre oralement un \cle{plan d'action} réaliste, en adoptant un comportement responsable et citoyen en matière de santé reproductive.
\end{itemize}

\textbf{Situation}\par

\textbf{Conseil municipal des jeunes de l'arrondissement de Nkol-Efoulan (fictif, inspiré de la région du Centre).}

Tu es membre du club santé de ton lycée. Le maire de l'arrondissement de Nkol-Efoulan, commune rurale de \num{42 000} habitants, sollicite les jeunes du lycée local pour préparer le \emph{Conseil municipal des jeunes} consacré à la santé reproductive. Il transmet au club la fiche de données suivante, compilée par le centre de santé intégré (CSI) de l'arrondissement et l'inspection de l'enseignement secondaire.

\begin{center}
\begin{tabularx}{\linewidth}{|l|X|}
\hline
\rowcolor{popblueL} \textbf{Indicateur} & \textbf{Donnée (année écoulée)} \\
\hline
Population de l'arrondissement & \num{42 000} habitants, dont \num{3 800} adolescentes de 15 à 19 ans \\
\hline
Grossesses chez les 15--19 ans & 312 grossesses déclarées (dont 41 chez les moins de 16 ans) \\
\hline
Abandon scolaire des jeunes mères & 78\,\% des jeunes mères scolarisées ont quitté l'école \\
\hline
Consultations prénatales & seulement 35\,\% des jeunes mères ont effectué au moins 3 consultations \\
\hline
IST dépistées au CSI & 214 cas (chlamydiose, gonococcie, syphilis, trichomonose) ; 61\,\% des cas ont moins de 25 ans \\
\hline
Dépistage VIH volontaire & 9\,\% des 15--24 ans ont fait un test dans l'année \\
\hline
Accès à l'information & aucune séance d'éducation à la santé reproductive n'a eu lieu dans l'arrondissement ; le poste de radio communautaire couvre 80\,\% des villages \\
\hline
Contexte économique & 55\,\% des familles vivent de l'agriculture de subsistance ; le CSI le plus proche est à plus de 10 km pour 6 villages sur 14 \\
\hline
\end{tabularx}
\end{center}

Le maire dispose d'un budget jeunesse de 3 millions de FCFA et attend des élèves un \textbf{plan d'action en 5 mesures concrètes}, chacune justifiée scientifiquement. Ta classe, organisée en groupes, joue le rôle des experts ; le reste de la classe et l'enseignant jouent le rôle du conseil municipal qui vote le plan.

\textbf{Consignes}\par

\begin{enumerate}
\item \textbf{Analyse des données.} Calcule le taux de grossesses chez les 15--19 ans de l'arrondissement (pour 100 adolescentes) et la proportion des cas d'IST touchant les moins de 25 ans. Compare le taux de grossesses obtenu à la valeur nationale estimée (environ 24\,\% des adolescentes de 15--19 ans ont déjà commencé une vie de mère au Cameroun). Que constates-tu ?
\item \textbf{Recherche des causes.} À partir de la fiche et de tes connaissances sur la puberté et la reproduction humaine, identifie au moins \textbf{six causes probables} des grossesses précoces et des IST dans cet arrondissement, en les classant en trois catégories : causes individuelles, causes sociales, causes économiques.
\item \textbf{Tableau des conséquences.} Construis un tableau à trois colonnes (mère / enfant / communauté) présentant au moins deux conséquences par colonne, en précisant pour chacune s'il s'agit d'une conséquence sanitaire, scolaire, sociale ou économique.
\item \textbf{Plan d'action.} Propose \textbf{5 mesures concrètes} (par exemple : séances d'information, distribution de préservatifs, dépistage gratuit, sensibilisation à la planification familiale, soutien à la scolarisation des jeunes mères). Pour chaque mesure, rédige une justification scientifique en une ou deux phrases : quel mécanisme biologique ou quel risque la mesure permet-elle de prévenir ?
\item \textbf{Présentation.} Présente ton plan en 3 minutes devant le « conseil municipal ». Le barème du vote valorise : l'exactitude scientifique (5 points), la justification des moyens de prévention (5 points), la faisabilité et l'adaptation au contexte local (5 points), la qualité de l'argumentation orale (5 points).
\end{enumerate}

\textbf{Ressources à mobiliser}\par

\begin{aretenir}[Ce que tu dois réactiver pour réussir]
\begin{itemize}
\item \textbf{Puberté et reproduction :} la puberté rend possible la fécondation (ovulation chez la fille, production de spermatozoïdes chez le garçon) ; toute relation sexuelle non protégée pendant la période fertile peut entraîner une grossesse.
\item \textbf{Grossesses précoces :} causes (rapports non protégés, méconnaissance du cycle, pressions sociales, pauvreté, mariages précoces) et conséquences (risques obstétricaux, avortements clandestins, déscolarisation, précarité).
\item \textbf{IST :} agents pathogènes (VIH, \emph{Chlamydia trachomatis}, \emph{Neisseria gonorrhoeae}, \emph{Treponema pallidum}), modes de transmission (rapports non protégés, sang, mère-enfant) et conséquences (infertilité, cancers, transmission au nouveau-né).
\item \textbf{Prévention :} abstinence, fidélité, \cle{préservatif} (seule méthode protégeant à la fois des grossesses et des IST), dépistage, traitement précoce.
\item \textbf{Planification familiale :} méthodes naturelles (abstinence périodique, méthode des températures), méthodes barrières (préservatifs), méthodes hormonales (pilule, injectables, implants), dispositif intra-utérin ; rôle de l'espacement des naissances.
\end{itemize}
\end{aretenir}

\textbf{Démarche et corrigé commenté}\par

\begin{exoresolu}[Corrigé commenté du plan d'action de Nkol-Efoulan]
Énoncé : analyser les données de l'arrondissement, dégager causes et conséquences, puis élaborer un plan d'action en 5 mesures justifiées scientifiquement.

\tcblower

\textbf{Étape 1 : analyse quantitative des données.}

Taux de grossesses chez les 15--19 ans :
\[
\frac{312}{3800} \times 100 \approx 8{,}2 \text{ grossesses pour 100 adolescentes en un an.}
\]
Proportion des cas d'IST chez les moins de 25 ans :
\[
\frac{0{,}61 \times 214}{214} \times 100 = 61\,\% \quad (\text{soit environ } 131 \text{ jeunes}).
\]
\emph{Interprétation :} un taux annuel de 8,2\,\% est très élevé : sur les cinq années que dure la tranche 15--19 ans, une fraction considérable des adolescentes sera concernée, ce qui rejoint la valeur nationale (environ 24\,\% ont déjà commencé une vie de mère). Le fait que 61\,\% des IST touchent les moins de 25 ans confirme que les jeunes sont la population la plus exposée, alors que seulement 9\,\% d'entre eux se font dépister : beaucoup d'infections restent inconnues et continuent de se transmettre.

\textbf{Étape 2 : causes probables (classification).}

\begin{center}
\begin{tabularx}{\linewidth}{|l|X|}
\hline
\rowcolor{popblueL} \textbf{Catégorie} & \textbf{Causes identifiées} \\
\hline
Individuelles & Méconnaissance du cycle menstruel et de la période fertile ; absence d'information (aucune séance d'éducation) ; rapports non protégés ; faible recours au dépistage (9\,\%) \\
\hline
Sociales & Pression des pairs et relations précoces ; mariages ou unions précoces ; tabou de la sexualité empêchant le dialogue parents-enfants ; stigmatisation des jeunes mères (78\,\% de déscolarisation) \\
\hline
Économiques & Pauvreté (55\,\% d'agriculture de subsistance) poussant à des relations transactionnelles ; éloignement du CSI (plus de 10 km pour 6 villages) limitant l'accès aux soins, aux préservatifs et au dépistage \\
\hline
\end{tabularx}
\end{center}

\textbf{Étape 3 : tableau des conséquences.}

\begin{center}
\begin{tabularx}{\linewidth}{|X|X|X|}
\hline
\rowcolor{popblueL} \textbf{Pour la mère} & \textbf{Pour l'enfant} & \textbf{Pour la communauté} \\
\hline
Risques obstétricaux accrus (hémorragies, fistules obstétricales, éclampsie) car le bassin et l'organisme ne sont pas pleinement matures avant 18 ans (sanitaire) & Prématurité et faible poids de naissance, mortalité infantile élevée (sanitaire) & Perte de main-d'œuvre qualifiée : 78\,\% des jeunes mères quittent l'école (économique) \\
\hline
Déscolarisation et perte d'autonomie (scolaire et économique) & Transmission mère-enfant possible d'IST non dépistées (VIH, syphilis congénitale) (sanitaire) & Charge pour les familles et le CSI ; cycle de la pauvreté entretenu (social et économique) \\
\hline
Recours à des avortements clandestins à risque vital (sanitaire) & Conditions de vie précaires, faible accès aux soins (social) & Augmentation des dépenses de santé publique (économique) \\
\hline
\end{tabularx}
\end{center}

\textbf{Étape 4 : plan d'action en 5 mesures justifiées.}

\begin{enumerate}
\item \textbf{Séances mensuelles d'éducation à la santé reproductive} au lycée et dans les 14 villages (via la radio communautaire qui couvre 80\,\% des villages). \emph{Justification :} la connaissance de la puberté, du cycle menstruel et de la période fertile permet d'identifier les situations à risque ; l'absence totale d'information est la première cause individuelle relevée.
\item \textbf{Distribution gratuite de préservatifs} au CSI, au lycée (point discret tenu par le club santé) et lors des séances. \emph{Justification :} le préservatif est la seule méthode qui protège à la fois des grossesses non désirées et des IST, car il constitue une barrière mécanique empêchant le contact entre les sécrétions génitales porteuses d'agents pathogènes et les spermatozoïdes.
\item \textbf{Campagne trimestrielle de dépistage gratuit VIH/IST avec unité mobile} pour les 6 villages éloignés de plus de 10 km. \emph{Justification :} de nombreuses IST (chlamydiose, VIH au stade primaire) sont asymptomatiques ; seul le dépistage permet un traitement précoce qui guérit les IST bactériennes, évite l'infertilité et interrompt la chaîne de transmission.
\item \textbf{Sensibilisation à la planification familiale} en partenariat avec les sages-femmes du CSI (consultations jeunes confidentielles). \emph{Justification :} les méthodes contraceptives hormonales (pilule, injectables, implants) inhibent l'ovulation ; le dispositif intra-utérin empêche la nidation ; elles permettent d'espacer les naissances, réduisant les risques obstétricaux liés aux grossesses trop précoces ou trop rapprochées.
\item \textbf{Fonds de soutien à la scolarisation des jeunes mères} (frais de scolarité, garderie au lycée, dispositif de rattrapage pédagogique). \emph{Justification :} rompre le lien grossesse $\rightarrow$ déscolarisation $\rightarrow$ pauvreté ; une jeune mère scolarisée transmet de meilleures pratiques de santé à son enfant et retrouve une autonomie économique.
\end{enumerate}

\textbf{Étape 5 : présentation et vote.} Chaque groupe présente son plan en 3 minutes ; le conseil municipal vote selon le barème. Le plan ci-dessus obtient les points maximaux s'il cite correctement les mécanismes (barrière du préservatif, ovulation et contraception hormonale, asymptomaticité des IST) et s'il reste dans le budget de 3 millions de FCFA.
\end{exoresolu}

\begin{attention}[Pièges fréquents dans cette activité]
\begin{itemize}
\item Confondre \textbf{contraception} et \textbf{prévention des IST} : la pilule empêche la grossesse mais ne protège d'aucune IST ; seul le préservatif assure la double protection.
\item Affirmer qu'une IST se voit toujours : beaucoup sont asymptomatiques, d'où l'importance du dépistage.
\item Proposer des mesures irréalistes (construire un hôpital avec 3 millions de FCFA) : la faisabilité compte pour un tiers du barème.
\item Oublier les causes sociales et économiques en ne citant que la « faute » individuelle : l'analyse doit rester scientifique et sans jugement moral.
\end{itemize}
\end{attention}

\textbf{Bilan de l'activité}\par

Cette activité t'a permis de mobiliser l'ensemble des savoir-faire de la leçon dans une démarche citoyenne complète. Tu as d'abord \emph{expliqué les risques} : les données chiffrées ont montré que les grossesses précoces exposent la mère à des complications obstétricales liées à l'immaturité de son organisme, l'enfant à la prématurité et aux infections, et la communauté à un appauvrissement durable par la déscolarisation massive (78\,\%). Tu as ensuite \emph{justifié les moyens de prévention} : le préservatif comme double protection (barrière mécanique contre spermatozoïdes et agents pathogènes), le dépistage comme seul moyen de révéler les IST asymptomatiques, la planification familiale comme outil d'espacement et de choix des naissances. Enfin, en défendant ton plan devant le « conseil municipal », tu as \emph{adopté un comportement responsable} : argumenter sans stigmatiser, proposer des solutions réalistes adaptées au contexte rural camerounais, et comprendre que la santé reproductive est une affaire collective où l'information, l'accès aux soins et le soutien scolaire se complètent. Retiens l'essentiel : \cle{informer, protéger, dépister, accompagner} — quatre verbes qui sauvent des vies et des avenirs.

\newpage

\coursec{Méthodes et exercices résolus}

\coursec{Amélioration de la santé reproductive}

\courssub{Puberté et bases de la reproduction humaine}

\begin{methode}[Expliquer le lien entre puberté et capacité de reproduction]
Pour expliquer pourquoi un adolescent devient fertile à la puberté, suis cette démarche en quatre étapes :
\begin{enumerate}
\item \textbf{Définir la puberté} : période de transformation de l'organisme (vers 10--14 ans chez la fille, 12--16 ans chez le garçon) sous l'action des hormones sexuelles.
\item \textbf{Nommer les hormones et leurs effets} : chez la fille, les œstrogènes et la progestérone (sécrétés par les ovaires) provoquent le développement des caractères sexuels secondaires et déclenchent les cycles menstruels ; chez le garçon, la testostérone (sécrétée par les testicules) déclenche la spermatogenèse et l'apparition des caractères sexuels secondaires.
\item \textbf{Relier à la fécondité} : dès les premières règles (ménarche), l'ovulation devient possible, donc une grossesse peut survenir dès le premier rapport non protégé, même avant les premières règles si l'ovulation a déjà eu lieu ; dès la première production de spermatozoïdes (spermarche), le garçon peut féconder.
\item \textbf{Conclure} : la maturité sexuelle (capacité à reproduire) précède la maturité physiologique complète (bassin, croissance achevée) et psychologique — c'est tout le problème des grossesses précoces.
\end{enumerate}
\textbf{Réflexe} : distingue toujours \emph{caractères sexuels primaires} (organes génitaux présents à la naissance) et \emph{caractères sexuels secondaires} (pilosité, développement des seins, mue de la voix, modification du bassin) qui apparaissent à la puberté.
\end{methode}

\begin{exoresolu}[Questions d'une élève sur la puberté]
Une élève de Première demande : « Ma petite sœur de 12 ans vient d'avoir ses premières règles. Cela signifie-t-il qu'elle peut déjà tomber enceinte ? Et pourquoi dit-on qu'une grossesse à son âge serait dangereuse, puisque son corps sait déjà fabriquer un bébé ? »\\
Répondez à ses deux questions en mobilisant vos connaissances sur la puberté et la reproduction.
\tcblower
\textbf{1. Peut-elle déjà tomber enceinte ?}

Oui. L'apparition des premières règles (ménarche) est le signe que les ovaires, sous l'action des hormones sexuelles (œstrogènes et progestérone), ont commencé à fonctionner et que des ovulations sont possibles. Or une grossesse survient dès qu'un spermatozoïde féconde un ovule lors d'un rapport non protégé. Il faut même savoir que la toute première ovulation \emph{précède} les premières règles : une fille peut donc théoriquement concevoir avant même d'avoir eu ses règles. Dès la puberté, tout rapport non protégé expose à une grossesse.

\textbf{2. Pourquoi une grossesse à 12 ans serait-elle dangereuse ?}

Il faut distinguer la \textbf{maturité sexuelle} de la \textbf{maturité physiologique complète}. À 12 ans, l'appareil reproducteur fonctionne, mais le reste du corps n'a pas terminé sa croissance :
\begin{itemize}
\item le \textbf{bassin osseux} est encore étroit : le passage du bébé lors de l'accouchement sera difficile (dystocie), avec risque de déchirures, de césarienne d'urgence, voire de décès de la mère ou de l'enfant ;
\item l'organisme de la fillette, lui-même en pleine croissance, entre en \textbf{compétition nutritionnelle} avec le fœtus : risque d'anémie sévère chez la mère, de prématurité et de faible poids de naissance (inférieur à \SI{2,5}{kg}) chez l'enfant ;
\item sur le plan \textbf{psychologique et social}, une enfant de 12 ans n'est pas préparée à assumer une maternité : risque de déscolarisation, de rejet familial et de souffrance psychique.
\end{itemize}

\textbf{Conclusion} : savoir fabriquer un bébé ne signifie pas que le corps est prêt à le faire sans danger. La capacité de reproduction apparaît à la puberté, bien avant que l'organisme soit complètement mature : c'est précisément pourquoi les grossesses précoces doivent être prévenues par l'information et, le cas échéant, la contraception.
\end{exoresolu}

\courssub{Grossesses précoces et non désirées : causes}

\begin{methode}[Analyser une situation de grossesse précoce]
Pour expliquer les risques liés aux grossesses précoces (avant 18 ans), suis toujours la même démarche en quatre étapes :
\begin{enumerate}
\item \textbf{Identifier les données du problème} : âge de la jeune fille, contexte (scolarisation, milieu familial, accès aux soins), éventuels facteurs favorisants (rapports non protégés, ignorance de la contraception, pression sociale, précocité des premiers rapports).
\item \textbf{Classer les risques en trois catégories} : risques \cle{médicaux} (pour la mère et pour l'enfant), risques \cle{sociaux} (abandon scolaire, rejet familial, précarité), risques \cle{psychologiques} (stress, isolement, dépression, troubles anxieux).
\item \textbf{Justifier chaque risque par un mécanisme biologique ou social précis} : un bassin osseux pas encore complètement développé avant 18 ans rend l'accouchement plus difficile (dystocie, risque de fistule obstétricale) ; l'organisme de l'adolescente, encore en croissance, entre en compétition avec le fœtus pour les nutriments (anémie, prématurité, faible poids de naissance).
\item \textbf{Conclure par une phrase de synthèse} reliant les risques identifiés à l'intérêt de la prévention (information, contraception, report des premières relations, maintien à l'école).
\end{enumerate}
\textbf{Réflexe} : au Probatoire, une réponse qui ne cite que « c'est dangereux » sans mécanisme précis ne vaut presque rien. Chaque risque énoncé doit être justifié par un fait biologique ou social établi.
\end{methode}

\begin{exoresolu}[Grossesse précoce dans un lycée de Douala]
Dans un lycée de Douala, une élève de 15 ans, en classe de Seconde, découvre qu'elle est enceinte de trois mois. Elle n'avait jamais reçu d'information sur la contraception et ses rapports n'étaient pas protégés.\\
1. Citez deux causes possibles de cette grossesse précoce.\\
2. Expliquez, en vous appuyant sur des mécanismes précis, trois risques médicaux auxquels cette adolescente et son enfant sont exposés.\\
3. Indiquez deux conséquences sociales probables de cette grossesse.
\tcblower
\textbf{1. Causes possibles de la grossesse précoce}

L'énoncé fournit directement deux causes :
\begin{itemize}
\item l'\textbf{ignorance des méthodes contraceptives} : l'élève « n'avait jamais reçu d'information sur la contraception », elle ne pouvait donc pas se protéger ;
\item l'\textbf{absence de protection lors des rapports sexuels} : « ses rapports n'étaient pas protégés », c'est-à-dire sans préservatif ni autre méthode contraceptive, ce qui permet la rencontre du spermatozoïde et de l'ovule (fécondation).
\end{itemize}
On peut ajouter, de façon générale, la précocité des premiers rapports sexuels, la pression de l'entourage ou des pairs, et l'absence de dialogue parents-enfants sur la sexualité.

\textbf{2. Trois risques médicaux justifiés}

\begin{itemize}
\item \textbf{Accouchement difficile (dystocie) et fistule obstétricale} : à 15 ans, le bassin osseux de l'adolescente n'a pas terminé sa croissance ; il est souvent trop étroit pour le passage du bébé. Cela expose à un accouchement prolongé, à des déchirures périnéales graves, à une césarienne d'urgence, et au risque de fistule vésico-vaginale ou recto-vaginale (incontinence définitive) en l'absence de prise en charge chirurgicale rapide. Cela met en danger la vie de la mère et de l'enfant.
\item \textbf{Anémie et retard de croissance intra-utérin} : l'organisme de l'adolescente est lui-même encore en croissance et a de forts besoins en fer, acide folique et nutriments. Le fœtus puise dans les mêmes réserves : il y a \emph{compétition nutritionnelle}. Il en résulte une anémie ferriprive chez la mère et un risque élevé de prématurité et de faible poids de naissance (moins de \SI{2,5}{kg}) chez l'enfant, facteur majeur de mortalité néonatale.
\item \textbf{Troubles hypertensifs de la grossesse (prééclampsie, éclampsie)} : les grossesses précoces s'accompagnent plus fréquemment d'hypertension artérielle gravidique et d'éclampsie (crises convulsives), qui constituent l'une des principales causes de mortalité maternelle chez les adolescentes, surtout lorsque le suivi prénatal est insuffisant ou tardif.
\end{itemize}

\textbf{3. Deux conséquences sociales probables}

\begin{itemize}
\item l'\textbf{abandon scolaire} : la jeune mère doit souvent interrompre ses études, ce qui compromet durablement son avenir professionnel, son autonomie financière et sa capacité à prendre en charge son enfant ;
\item le \textbf{rejet familial et la précarité} : la stigmatisation peut conduire à l'exclusion du foyer familial, plongeant la jeune mère et son enfant dans la pauvreté et l'isolement social.
\end{itemize}

\textbf{Conclusion} : cette grossesse précoce, favorisée par le manque d'information et l'absence de contraception, expose l'adolescente à des risques médicaux graves (dystocie, anémie, hypertension gravidique) et à des conséquences sociales lourdes (déscolarisation, précarité). D'où l'importance capitale de l'éducation à la santé reproductive dès le collège.
\end{exoresolu}

\courssub{Infections sexuellement transmissibles (IST) et prévention}

\begin{methode}[Justifier l'utilisation des moyens de prévention des IST]
Pour justifier l'utilisation d'un moyen de prévention des infections sexuellement transmissibles (IST), procède en quatre étapes :
\begin{enumerate}
\item \textbf{Rappeler le mode de transmission} de l'IST concernée : les IST se transmettent par les sécrétions génitales (sperme, pertes vaginales), le sang, ou le contact direct de muqueuses lors de rapports sexuels non protégés ; certaines (VIH, hépatite B, syphilis) se transmettent aussi par le sang contaminé (matériel d'injection, scarifications) et de la mère à l'enfant (grossesse, accouchement, allaitement).
\item \textbf{Identifier le mécanisme d'action du moyen de prévention} : le préservatif (masculin ou féminin) est une \cle{barrière mécanique} qui empêche le contact entre les sécrétions et les muqueuses des partenaires ; le dépistage permet de traiter précocement et de rompre la chaîne de transmission ; la vaccination (hépatite B, HPV) induit une immunité active.
\item \textbf{Relier le mécanisme à l'efficacité} : en bloquant le passage des agents pathogènes (VIH, \emph{Neisseria gonorrhoeae}, \emph{Chlamydia trachomatis}, \emph{Treponema pallidum}, virus de l'hépatite B, HPV...), le préservatif correctement utilisé à chaque rapport réduit fortement le risque de transmission sexuelle.
\item \textbf{Préciser les limites et les conditions d'efficacité} : le préservatif doit être utilisé correctement, du début à la fin du rapport, à chaque rapport ; il ne protège que s'il est intact, non périmé et bien mis. L'abstinence et la fidélité mutuelle entre partenaires dépistés non infectés sont d'autres stratégies complémentaires.
\end{enumerate}
\textbf{Réflexe} : cite toujours le \emph{double avantage} du préservatif : il protège à la fois des IST \textbf{et} des grossesses non désirées. C'est la seule méthode à double protection.
\end{methode}

\begin{exoresolu}[Campagne de prévention dans un quartier de Yaoundé]
Lors d'une campagne de santé dans un quartier de Yaoundé, une infirmière distribue des préservatifs à des jeunes de 16 à 20 ans. Un jeune homme déclare : « Le préservatif ne sert à rien, il suffit de choisir un partenaire qui a l'air en bonne santé. »\\
1. Expliquez pourquoi l'apparence physique ne permet pas de détecter une IST. Donnez un exemple précis.\\
2. Justifiez, par son mécanisme d'action, l'efficacité du préservatif dans la prévention des IST.\\
3. Citez deux autres moyens de prévention des IST.
\tcblower
\textbf{1. L'apparence physique ne révèle pas une IST}

De nombreuses IST sont \textbf{asymptomatiques}, c'est-à-dire que la personne infectée ne présente aucun signe visible pendant des mois ou des années, tout en étant contagieuse. \textbf{Exemple} : une personne infectée par le VIH peut paraître en parfaite santé pendant la phase asymptomatique (qui dure souvent plusieurs années sans traitement) tout en pouvant transmettre le virus. De même, l'infection à \emph{Chlamydia trachomatis} est asymptomatique chez la majorité des femmes infectées et chez une proportion importante d'hommes. Seul un \textbf{test de dépistage} en laboratoire (sérologie, PCR, examen direct) permet de savoir si une personne est infectée. L'affirmation du jeune homme est donc fausse et dangereuse.

\textbf{2. Mécanisme d'action du préservatif}

Le préservatif masculin est une gaine fine (généralement en latex, ou en polyuréthane/polyisoprène en cas d'allergie) qui recouvre le pénis en érection et recueille le sperme et les sécrétions. Il constitue une \textbf{barrière mécanique étanche} qui empêche tout contact direct entre les sécrétions génitales (sperme, pertes vaginales) et les muqueuses des deux partenaires. Or les agents pathogènes responsables des IST se transmettent précisément par ces sécrétions et ces contacts de muqueuses. En bloquant leur passage, le préservatif, utilisé correctement du début à la fin de chaque rapport, réduit très fortement le risque de transmission. Il présente en outre un double avantage : il prévient aussi les grossesses non désirées.

\textbf{3. Deux autres moyens de prévention}

\begin{itemize}
\item le \textbf{dépistage régulier} (et celui du partenaire) suivi d'un traitement précoce en cas de positivité, qui rompt la chaîne de transmission et évite les complications (stérilité, cancers) ;
\item l'\textbf{abstinence} ou la \textbf{fidélité mutuelle} entre partenaires non infectés et dépistés ; on peut ajouter la vaccination contre l'hépatite B et le papillomavirus humain (HPV), et le non-partage d'objets tranchants souillés de sang (rasoirs, aiguilles, matériel de scarification).
\end{itemize}

\textbf{Conclusion} : une IST ne se voit pas sur le visage ; seuls le dépistage et l'usage systématique et correct du préservatif, barrière mécanique bloquant les agents pathogènes, protègent efficacement les adolescents.
\end{exoresolu}

\courssub{Interpréter des données épidémiologiques sur les IST}

\begin{methode}[Exploiter un tableau ou un graphique de données sanitaires]
Face à un tableau de données (prévalence d'une IST, nombre de grossesses précoces...), la démarche est toujours la même :
\begin{enumerate}
\item \textbf{Lire soigneusement le titre, les légendes et les unités} du tableau ou du graphique (que mesure-t-on ? chez qui ? en quelle unité ? quelle population ? quelle période ?).
\item \textbf{Dégager les faits} : comparer les valeurs (la plus grande, la plus faible, évolution au cours du temps), sans les interpréter tout de suite. Un fait se formule avec les chiffres : « la prévalence passe de $x$ \% à $y$ \% ».
\item \textbf{Calculer si nécessaire} : variation absolue ($y - x$), variation relative $\dfrac{y-x}{x}\times 100$ en pourcentage.
\item \textbf{Interpréter} : relier les faits observés à une cause biologique ou comportementale plausible (absence de protection, campagne de dépistage efficace, amélioration de l'accès aux soins...).
\item \textbf{Conclure} par une phrase qui répond exactement à la question posée.
\end{enumerate}
\textbf{Réflexe} : ne jamais écrire « on voit que ça augmente » sans donner les valeurs chiffrées qui le prouvent. Au Probatoire, l'exploitation de document sans calcul ni citation de chiffres est sanctionnée par une note nulle ou très faible.
\end{methode}

\begin{exoresolu}[Évolution des dépistages VIH chez les jeunes]
Le tableau ci-dessous indique le nombre de jeunes de 15 à 24 ans dépistés positifs au VIH dans un centre de santé de la région du Centre, selon qu'ils déclarent utiliser systématiquement le préservatif ou non (étude portant sur \num{2000} jeunes dépistés dans chaque groupe).
\begin{center}
\begin{tabularx}{\linewidth}{|l|X|X|}\hline
\rowcolor{popblueL} Groupe & Nombre de dépistés & Nombre de tests positifs \\ \hline
Utilisation systématique du préservatif & \num{2000} & 20 \\ \hline
Pas d'utilisation systématique & \num{2000} & 160 \\ \hline
\end{tabularx}
\end{center}
1. Calculez la prévalence (pourcentage de tests positifs) dans chaque groupe.\\
2. Comparez les deux prévalences.\\
3. Que concluez-vous sur le rôle du préservatif ?
\tcblower
\textbf{1. Calcul des prévalences}

La prévalence est le pourcentage de personnes testées positives dans un groupe :
\[
\text{Prévalence} = \dfrac{\text{nombre de tests positifs}}{\text{nombre de dépistés}} \times 100.
\]
\begin{align*}
\text{Groupe avec préservatif systématique : } & \dfrac{20}{2000} \times 100 = 1\,\%.\\
\text{Groupe sans préservatif systématique : } & \dfrac{160}{2000} \times 100 = 8\,\%.
\end{align*}

\textbf{2. Comparaison}

La prévalence du VIH est de $1\,\%$ chez les utilisateurs systématiques du préservatif contre $8\,\%$ chez les non-utilisateurs, soit \textbf{8 fois plus} (ou une différence de $8 - 1 = 7$ points de pourcentage).

\textbf{3. Conclusion}

À exposition comparable, l'utilisation systématique du préservatif est associée à une prévalence du VIH huit fois plus faible. Cela confirme que le préservatif, en empêchant le contact entre sécrétions et muqueuses, est un moyen de prévention très efficace de la transmission sexuelle du VIH. Il faut donc encourager son usage correct et systématique chez les jeunes.
\end{exoresolu}

\courssub{Adopter un comportement responsable en matière de santé reproductive}

\begin{methode}[Rédiger un conseil ou un plan d'action responsable]
Lorsqu'un sujet demande de « proposer un comportement responsable » ou de « conseiller un adolescent », structure ta réponse ainsi :
\begin{enumerate}
\item \textbf{Situer le problème} en une phrase (risque d'IST, de grossesse non désirée, pression de l'entourage, méconnaissance des méthodes...).
\item \textbf{Proposer des attitudes concrètes et réalistes}, classées en trois registres :
\begin{itemize}
\item \emph{Information} : se renseigner auprès des centres de santé (CSPC, hôpitaux), des infirmiers scolaires, des campagnes de sensibilisation (ex. Journée mondiale de lutte contre le sida, vacances sans sida), des numéros verts ;
\item \emph{Prévention} : abstinence ou report des premiers rapports, usage correct et systématique du préservatif (double protection), dépistage régulier (VIH, hépatites, syphilis), vaccination (hépatite B, HPV) ;
\item \emph{Responsabilité} : refuser les rapports forcés ou sous pression (viol, chantage), dialoguer avec son partenaire sur la protection et le consentement, consulter rapidement en cas de signe d'IST (écoulement, brûlures mictionnelles, ulcération, bubon) sans automédication, respecter le choix de l'autre.
\end{itemize}
\item \textbf{Justifier chaque attitude} par son bénéfice (protection contre les IST et les grossesses précoces, préservation de la scolarité, de la fertilité future et de la santé).
\item \textbf{Conclure} en rappelant que la santé reproductive est une responsabilité partagée entre les deux partenaires.
\end{enumerate}
\textbf{Réflexe} : un conseil sans justification est une réponse incomplète ; chaque proposition doit être reliée à un risque évité ou à un droit protégé.
\end{methode}

\begin{exoresolu}[Conseiller un ami]
Ton ami de 17 ans te confie qu'il envisage d'avoir ses premiers rapports sexuels avec sa petite amie, mais qu'il hésite à utiliser un préservatif car « ce n'est pas naturel ». Rédige, en quatre à six lignes structurées, le conseil argumenté que tu lui donnerais pour l'aider à adopter un comportement responsable.
\tcblower
\textbf{Conseil argumenté}

Cher ami, je comprends ton hésitation, mais sache que le préservatif est le seul moyen qui te protège à la fois des infections sexuellement transmissibles, dont le VIH/sida, et d'une grossesse non désirée qui pourrait bouleverser vos études à tous les deux. Une IST ne se devine pas à l'apparence : une personne peut être infectée et contagieuse sans aucun signe visible. Utilisé correctement, du début à la fin de chaque rapport, le préservatif constitue une barrière mécanique qui bloque le passage des agents pathogènes. Je te conseille aussi d'aller ensemble vous informer dans un centre de santé (CSPC, hôpital de district) et, pourquoi pas, de faire un dépistage : c'est une preuve de respect mutuel et de maturité, pas de méfiance. Enfin, n'oublie jamais qu'un rapport doit être librement consenti par les deux partenaires, sans aucune pression ni contrainte. Être responsable, c'est protéger sa santé, celle de l'autre et son avenir.
\end{exoresolu}

\courssub{Planification familiale (notions générales)}

\begin{methode}[Reconnaître et classer les méthodes contraceptives]
Pour répondre à une question sur la planification familiale (notions générales) :
\begin{enumerate}
\item \textbf{Définir} la planification familiale : ensemble des méthodes permettant à un couple ou à une personne de choisir le nombre de ses enfants et d'espacer les naissances, dans le respect de la santé et des droits reproductifs.
\item \textbf{Classer les méthodes} en deux grandes catégories :
\begin{itemize}
\item méthodes \cle{traditionnelles / naturelles} : abstinence périodique (méthode du calendrier / Ogino-Knaus, basée sur la période fertile du cycle menstruel), retrait (coït interrompu), méthode de l'allaitement et de l'aménorrhée (MAMA) — peu fiables (indice de Pearl élevé) ;
\item méthodes \cle{modernes} : préservatifs (masculin et féminin), pilule contraceptive (œstro-progestative ou microprogestative), injectables (progestatifs), implant sous-cutané, dispositif intra-utérin (DIU ou stérilet, au cuivre ou hormonal), stérilisation définitive (ligature des trompes, vasectomie).
\end{itemize}
\item \textbf{Préciser le mode d'action} de chaque méthode citée : barrière mécanique bloquant les spermatozoïdes (préservatifs), blocage de l'ovulation par des hormones (pilule, injectable, implant), action spermicide et modification de l'endomètre empêchant la nidation (DIU au cuivre), épaississement de la glaire cervicale et atrophie endométriale (DIU hormonal).
\item \textbf{Dégager l'intérêt} : espacement des naissances (meilleure santé de la mère et des enfants, réduction de la mortalité maternelle et infantile), prévention des grossesses non désirées et des avortements à risque (clandestins), poursuite des études et autonomisation des filles/femmes, amélioration du niveau de vie du ménage.
\end{enumerate}
\textbf{Réflexe} : rappelle toujours que seul le préservatif (masculin ou féminin) protège aussi des IST ; la pilule, l'implant, l'injectable ou le DIU évitent la grossesse mais \emph{ne protègent pas} des infections. La double protection (contraception hormonale + préservatif) est la stratégie idéale pour les adolescents sexuellement actifs.
\end{methode}

\begin{exoresolu}[Séance d'information sur la planification familiale]
Lors d'une séance d'information dans un centre de santé de Bafoussam, une animatrice présente aux jeunes couples la pilule contraceptive, le préservatif et la méthode du calendrier.\\
1. Définissez la planification familiale.\\
2. Pour chacune des trois méthodes citées, indiquez son mode d'action et sa fiabilité.\\
3. Un couple demande : « La pilule suffit-elle à nous protéger des IST ? » Répondez en justifiant.
\tcblower
\textbf{1. Définition}

La planification familiale est l'ensemble des méthodes qui permettent à un couple ou à une personne de choisir librement et en toute connaissance de cause le nombre de ses enfants et d'espacer les naissances, afin de préserver la santé de la mère et des enfants et d'améliorer les conditions de vie de la famille.

\textbf{2. Mode d'action et fiabilité des trois méthodes}

\begin{itemize}
\item \textbf{La pilule contraceptive (œstro-progestative)} : elle contient des hormones de synthèse (analogues des œstrogènes et de la progestérone) qui inhibent le pic de LH, bloquant ainsi l'ovulation : sans libération d'ovule, la fécondation est impossible. Elle épaissit aussi la glaire cervicale. Prise correctement chaque jour à heure fixe, elle est très fiable (indice de Pearl < 1), mais un oubli de plus de 12 h (pilule combinée) ou 3 h (microprogestative) réduit fortement son efficacité.
\item \textbf{Le préservatif masculin} : c'est une barrière mécanique en latex (ou polyuréthane) qui retient les spermatozoïdes au moment de l'éjaculation et empêche leur rencontre avec l'ovule. Utilisé correctement à chaque rapport (mis avant tout contact génital, retiré après éjaculation), il est fiable (indice de Pearl 2 à 15 selon l'usage) et c'est la \textbf{seule méthode qui protège aussi des IST}.
\item \textbf{La méthode du calendrier} (abstinence périodique, méthode d'Ogino) : elle consiste à éviter les rapports pendant la période fertile du cycle, autour de l'ovulation (vers le 14\textsuperscript{e} jour d'un cycle de 28 jours). Elle est \textbf{peu fiable} (indice de Pearl > 20), car la date d'ovulation varie d'un cycle à l'autre (stress, fatigue, maladie) et les spermatozoïdes survivent 3 à 5 jours dans les voies génitales féminines, élargissant la fenêtre de fécondité.
\end{itemize}

\textbf{3. La pilule protège-t-elle des IST ?}

Non. La pilule agit uniquement sur le plan hormonal en bloquant l'ovulation et en modifiant la glaire cervicale ; elle n'empêche en rien le contact direct entre les sécrétions génitales (sperme, pertes vaginales) et les muqueuses des partenaires, par lequel se transmettent les agents pathogènes des IST (VIH, gonocoque, chlamydiae, tréponème, virus de l'hépatite B, HPV...). Seul le préservatif, barrière mécanique, assure cette protection. Un couple exposé au risque d'IST doit donc utiliser le préservatif, éventuellement en complément d'une autre méthode contraceptive (double protection).

\textbf{Conclusion} : le choix d'une méthode de planification familiale doit se faire en connaissant son mode d'action, son efficacité et ses limites ; pour une double protection (grossesse et IST), le préservatif reste indispensable.
\end{exoresolu}

\begin{attention}[Les 5 erreurs qui coûtent des points au Probatoire]
\begin{enumerate}
\item \textbf{Confondre contraception et prévention des IST.} La pilule, l'implant, l'injectable ou le DIU évitent la grossesse mais ne protègent \emph{pas} des IST ; seul le préservatif (masculin ou féminin) assure la double protection. C'est la confusion la plus fréquente et la plus sanctionnée.
\item \textbf{Affirmer qu'on reconnaît une personne infectée à son apparence.} Faux : la plupart des IST (VIH, chlamydiae, HPV) sont asymptomatiques pendant longtemps ; seul un test de dépistage est fiable. Ne jamais écrire « un partenaire sain en apparence ne présente pas de risque ».
\item \textbf{Citer un risque sans le justifier.} Écrire « la grossesse précoce est dangereuse » sans mécanisme (bassin immature $\Rightarrow$ dystocie/fistule ; compétition nutritionnelle $\Rightarrow$ anémie et prématurité ; immaturité vasculaire $\Rightarrow$ prééclampsie) fait perdre la moitié des points : chaque risque énoncé doit être expliqué biologiquement ou socialement.
\item \textbf{Oublier les chiffres dans l'exploitation d'un tableau.} Une interprétation de données sans valeurs chiffrées (prévalences calculées, comparaison « 8 fois plus », variation en points de pourcentage) est considérée comme non traitée. Calcule toujours les pourcentages et cite-les explicitement dans la conclusion.
\item \textbf{Employer un vocabulaire vague ou incorrect.} Écris « infections sexuellement transmissibles (IST) » et non « maladies vénériennes » ou « maladies sales » ; « préservatif » et non « capote » (langage familier) ; distingue rigoureusement \emph{fécondation} (rencontre spermatozoïde--ovule dans la trompe), \emph{nidation} (implantation du blastocyste dans l'endomètre utérin) et \emph{ovulation} (libération de l'ovocyte II par l'ovaire). Le vocabulaire scientifique précis est un critère explicite de notation.
\end{enumerate}
\end{attention}

\newpage

\coursec{Exercices}

\courssub{Je vérifie mes connaissances}

\begin{exercice}[QCM : la puberté et la reproduction]
Pour chaque question, une seule réponse est exacte. Recopie la lettre de la bonne réponse et justifie ton choix en une phrase.

\textbf{1.} La puberté est déclenchée par :
\begin{itemize}
\item a) la production d'hormones sexuelles par les gonades sous le contrôle de l'hypothalamus et de l'hypophyse ;
\item b) l'apparition des premières règles chez la fille ;
\item c) la croissance rapide du squelette ;
\item d) le développement des caractères sexuels secondaires.
\end{itemize}

\textbf{2.} Chez la fille, la fécondation a normalement lieu :
\begin{itemize}
\item a) dans l'utérus ; \quad b) dans le vagin ; \quad c) dans le tiers externe de la trompe de Fallope ; \quad d) dans l'ovaire.
\end{itemize}

\textbf{3.} Les premières règles (ménarche) surviennent en moyenne :
\begin{itemize}
\item a) entre 8 et 10 ans ; \quad b) entre 12 et 14 ans ; \quad c) entre 16 et 18 ans ; \quad d) après 20 ans.
\end{itemize}

\textbf{4.} La nidation correspond :
\begin{itemize}
\item a) à la rencontre entre l'ovule et le spermatozoïde ;
\item b) à l'implantation de l'embryon dans la muqueuse utérine ;
\item c) à la libération de l'ovule par l'ovaire ;
\item d) à l'expulsion du fœtus lors de l'accouchement.
\end{itemize}
\end{exercice}

\begin{exercice}[Vrai ou faux : les IST]
Réponds par \textbf{vrai} ou \textbf{faux}, puis justifie chaque réponse en une ou deux phrases.

\textbf{1.} Le VIH se transmet par les poignées de main, les embrassades et le partage des repas.

\textbf{2.} Le préservatif, correctement utilisé, est le seul moyen de contraception qui protège aussi des IST.

\textbf{3.} Une personne séropositive sans symptôme apparent ne peut pas transmettre le VIH.

\textbf{4.} La syphilis et la gonococcie sont des IST d'origine bactérienne qui se soignent par des antibiotiques.

\textbf{5.} On peut dépister le VIH gratuitement et anonymement dans de nombreux centres de santé au Cameroun.
\end{exercice}

\begin{exercice}[Définitions]
Définis en deux ou trois phrases chacun des termes suivants, en employant le vocabulaire scientifique approprié :
\begin{enumerate}
\item puberté ;
\item grossesse précoce (ou grossesse en milieu scolaire) ;
\item infection sexuellement transmissible (IST) ;
\item planification familiale ;
\item contraception.
\end{enumerate}
\end{exercice}

\begin{exercice}[Calculs directs : statistiques de santé scolaire]
Dans un lycée de la région du Centre comptant \num{1200} élèves, dont \num{640} filles, le service de santé scolaire a recensé au cours de l'année 32 grossesses chez les élèves.

\textbf{1.} Calcule le pourcentage de filles de ce lycée concernées par une grossesse cette année-là (arrondis au dixième).

\textbf{2.} Sachant que 18 de ces 32 élèves ont abandonné leurs études, calcule la proportion de déscolarisation parmi les élèves enceintes, en pourcentage.

\textbf{3.} L'année précédente, on comptait 40 grossesses. Calcule le taux d'évolution (en \%) du nombre de grossesses entre les deux années et commente le résultat en une phrase.
\end{exercice}

\courssub{Je m'entraîne}

\begin{exercice}[Classement des méthodes contraceptives]
Voici une liste de méthodes contraceptives : pilule, préservatif masculin, stérilet, retrait, implant sous-cutané, ligature des trompes, préservatif féminin, pilule du lendemain.

\textbf{1.} Recopie et complète le tableau ci-dessous en rangeant chaque méthode dans la catégorie qui lui convient.

\begin{center}
\begin{tabularx}{\linewidth}{|l|X|}
\hline
\rowcolor{popblueL} \textbf{Catégorie} & \textbf{Méthodes} \\
\hline
Méthode de barrière (mécanique) & \\
\hline
Méthode hormonale & \\
\hline
Méthode mécanique intra-utérine & \\
\hline
Méthode définitive (chirurgicale) & \\
\hline
Méthode naturelle / non fiable & \\
\hline
\end{tabularx}
\end{center}

\textbf{2.} Indique, parmi toutes ces méthodes, celles qui protègent également des IST. Justifie ta réponse.

\textbf{3.} Explique pourquoi le retrait et la méthode des températures sont considérés comme peu fiables.
\end{exercice}

\begin{exercice}[Schéma à légender : l'appareil reproducteur féminin]
Le schéma ci-dessous représente l'appareil reproducteur féminin vu de face.

\begin{popfigure}
\begin{center}
\begin{tikzpicture}[scale=1.0]
% Utérus
\draw[thick, poppink] (0,0.2) .. controls (-0.9,0.4) and (-1.0,1.6) .. (0,2.0)
      .. controls (1.0,1.6) and (0.9,0.4) .. (0,0.2) -- cycle;
% Col et vagin
\draw[thick, poppink] (0,0.2) -- (0,-0.4);
\draw[thick, poppink] (-0.25,-0.4) -- (-0.25,-1.4);
\draw[thick, poppink] (0.25,-0.4) -- (0.25,-1.4);
\draw[thick, poppink] (-0.25,-1.4) .. controls (0,-1.7) .. (0.25,-1.4);
% Trompes
\draw[thick, poppurple] (-0.75,1.55) .. controls (-1.6,1.9) and (-2.2,1.5) .. (-2.6,0.9);
\draw[thick, poppurple] (0.75,1.55) .. controls (1.6,1.9) and (2.2,1.5) .. (2.6,0.9);
% Ovaires
\draw[thick, poporange, fill=poporangeL] (-2.7,0.55) ellipse (0.35 and 0.25);
\draw[thick, poporange, fill=poporangeL] (2.7,0.55) ellipse (0.35 and 0.25);
% Ovule et trajet
\fill[popdark] (-2.7,0.55) circle (0.07);
\draw[-{Stealth}, thick, popdark, dashed] (-2.7,0.55) .. controls (-2.5,1.3) and (-1.5,1.8) .. (-0.6,1.5);
\draw[-{Stealth}, thick, popdark, dashed] (-0.6,1.5) .. controls (-0.3,1.3) and (-0.15,1.0) .. (0,0.7);
% Étiquettes à compléter
\node[left] at (-3.15,0.55) {\textbf{1}};
\node[above] at (-1.7,1.95) {\textbf{2}};
\node[right] at (1.05,1.1) {\textbf{3}};
\node[right] at (0.3,-0.4) {\textbf{4}};
\node[right] at (0.3,-1.2) {\textbf{5}};
\end{tikzpicture}
\end{center}
\legende{Schéma simplifié de l'appareil reproducteur féminin (vue de face). Les flèches en pointillés indiquent le trajet de l'ovule.}
\end{popfigure}

\textbf{1.} Légende le schéma en attribuant à chaque numéro (1 à 5) le nom de l'organe correspondant, choisi parmi : vagin, utérus, ovaire, trompe de Fallope, col de l'utérus.

\textbf{2.} Décris en trois phrases le trajet de l'ovule, de sa libération jusqu'à la nidation, en nommant les organes traversés.

\textbf{3.} Précise le lieu exact où se produit normalement la fécondation et explique ce qui se passe si aucun spermatozoïde ne rencontre l'ovule.
\end{exercice}

\begin{exercice}[Étude de cas : une décision difficile]
Amina, 16 ans, élève en classe de Seconde dans un lycée de Douala, découvre qu'elle est enceinte de deux mois. Paniquée à l'idée de la réaction de ses parents et de son exclusion de l'établissement, une voisine lui propose un avortement clandestin pratiqué à domicile. Amina hésite.

\textbf{1.} Cite trois risques graves auxquels Amina s'exposerait en acceptant un avortement clandestin (risques médicaux et risque juridique).

\textbf{2.} Propose trois attitudes responsables qu'Amina devrait adopter à la place, en précisant à qui elle peut s'adresser (personnel de santé, famille, structures d'accueil).

\textbf{3.} Explique en quelques lignes pourquoi la poursuite de la scolarité d'Amina est un droit et une chance pour son avenir.
\end{exercice}

\begin{exercice}[Expliquer pour convaincre]
Ton jeune frère, élève en classe de 3\textsuperscript{e}, affirme : « Les préservatifs, c'est pour les adultes mariés ; à notre âge, il suffit de faire attention. »

Rédige un texte argumenté de 10 à 15 lignes destiné à lui expliquer pourquoi le préservatif est le moyen de prévention le plus adapté aux adolescents sexuellement actifs. Ton texte devra mobiliser au moins trois arguments scientifiques (double protection, accessibilité, absence d'effet hormonal) et se terminer par un conseil de comportement responsable.
\end{exercice}

\courssub{Je me prépare au Probatoire}

\begin{exercice}[Type Probatoire : les grossesses en milieu scolaire]
Le tableau ci-dessous présente le nombre de grossesses recensées parmi les élèves filles d'un groupement de lycées de la région du Littoral sur trois années scolaires.

\begin{center}
\begin{tabularx}{\linewidth}{|l|X|X|X|}
\hline
\rowcolor{popblueL} \textbf{Année scolaire} & \textbf{2021-2022} & \textbf{2022-2023} & \textbf{2023-2024} \\
\hline
Effectif des filles & \num{4500} & \num{4700} & \num{5000} \\
\hline
Nombre de grossesses & 135 & 141 & 160 \\
\hline
Déscolarisations liées & 81 & 92 & 112 \\
\hline
\end{tabularx}
\end{center}

\textbf{1.} Calcule, pour chacune des trois années, le taux de grossesses pour 1000 filles (arrondis à l'unité). Présente tes résultats dans un tableau.

\textbf{2.} Compare l'évolution du taux de grossesses à celle de l'effectif des filles. Que constates-tu ?

\textbf{3.} Calcule le pourcentage moyen de déscolarisation parmi les élèves enceintes sur la période 2021-2024.

\textbf{4.} Dégage deux causes des grossesses précoces en milieu scolaire et deux conséquences (une sur la santé de la mère, une sur son avenir scolaire et social).

\textbf{5.} Propose deux mesures concrètes qu'un club de santé scolaire pourrait mettre en œuvre pour réduire ce phénomène.
\end{exercice}

\begin{exercice}[Type Probatoire : prévention des IST chez les jeunes]
Lors d'une campagne de dépistage organisée dans un quartier de Yaoundé, \num{2000} jeunes de 15 à 24 ans se sont portés volontaires pour un test VIH. Les résultats sont les suivants : 60 jeunes se sont révélés séropositifs ; parmi eux, 45 ignoraient leur statut sérologique. Par ailleurs, une enquête réalisée auprès des mêmes jeunes montre que seulement 35 \% déclarent utiliser systématiquement le préservatif lors des rapports sexuels.

\textbf{1.} Calcule la prévalence du VIH (en \%) dans cet échantillon de jeunes.

\textbf{2.} Calcule la proportion de séropositifs qui ignoraient leur état. Explique en deux phrases pourquoi cette situation favorise la propagation de l'épidémie.

\textbf{3.} Cite les trois modes de transmission du VIH et, pour chacun, un moyen de prévention adapté.

\textbf{4.} Nomme deux autres IST que le VIH, en précisant pour chacune l'agent pathogène responsable (virus ou bactérie) et un symptôme caractéristique.

\textbf{5.} Justifie, en un paragraphe argumenté, l'affirmation suivante : « Le dépistage régulier et l'usage systématique du préservatif sont les deux piliers de la lutte contre les IST chez les adolescents. »
\end{exercice}

\begin{exercice}[Type Probatoire : sujet de synthèse sur la santé reproductive]
\textbf{Document 1.} Au Cameroun, selon les enquêtes démographiques et de santé, environ une adolescente de 15 à 19 ans sur quatre a déjà commencé sa vie de maternité (enceinte ou déjà mère), avec des proportions plus élevées dans les régions du Nord et en milieu rural.

\textbf{Document 2.} Les complications liées à la grossesse et à l'accouchement constituent l'une des principales causes de décès des adolescentes de 15 à 19 ans dans le monde. Le corps d'une adolescente n'est pas toujours complètement prêt à mener une grossesse à terme : bassin parfois insuffisamment développé, risques d'anémie, d'hypertension gravidique et de fistule obstétricale.

\textbf{Document 3.} La planification familiale désigne l'ensemble des méthodes permettant aux individus et aux couples de choisir librement le nombre de leurs enfants et l'espacement des naissances. Elle comprend des méthodes modernes (pilule, implant, injectable, stérilet, préservatifs) et des méthodes traditionnelles, de fiabilité variable.

À partir de l'étude critique de ces documents et de tes connaissances :

\textbf{1.} Définis les expressions suivantes : santé reproductive ; grossesse précoce ; planification familiale.

\textbf{2.} Explique, en t'appuyant sur le document 2, pourquoi une grossesse chez une adolescente de 15 ans est dite « à haut risque » (trois arguments attendus).

\textbf{3.} À partir du document 1, dégage deux facteurs qui peuvent expliquer la fréquence des maternités précoces dans certaines régions du Cameroun.

\textbf{4.} Classe les méthodes citées dans le document 3 selon leur mode d'action (barrière, hormonale, intra-utérine) et indique celle(s) qui protège(nt) aussi des IST.

\textbf{5.} Rédige un texte de 15 à 20 lignes intitulé « Adopter un comportement responsable en matière de santé reproductive », destiné à sensibiliser les élèves de ta classe. Tu y présenteras les risques des grossesses précoces et des IST, les moyens de prévention disponibles, et le rôle de l'information, de la famille et des services de santé.
\end{exercice}

\newpage

\coursec{Corrigés des exercices}

\begin{corrige}[Exercice 1 — QCM : la puberté et la reproduction]

\textbf{1. Réponse a.} La puberté est déclenchée par la mise en activité de l'axe hypothalamo-hypophysaire : l'hypothalamus sécrète la GnRH, qui stimule l'hypophyse antérieure (FSH et LH), laquelle active les gonades (ovaires ou testicules) productrices d'hormones sexuelles (œstrogènes, progestérone, testostérone). Les réponses b, c et d sont des \emph{conséquences} de la puberté (apparition des caractères sexuels secondaires, croissance, ménarche/spermarche), non sa cause initiale.

\textbf{2. Réponse c.} La fécondation (rencontre de l'ovule et du spermatozoïde) a normalement lieu dans le tiers externe de la trompe de Fallope (ampoule), près du pavillon. L'utérus est le lieu de la nidation, l'ovaire le lieu de la production de l'ovule (ovulation).

\textbf{3. Réponse b.} Les premières règles (ménarche) surviennent en moyenne entre 12 et 14 ans, avec des variations individuelles normales (généralement entre 11 et 15 ans) selon la génétique, la nutrition et l'environnement.

\textbf{4. Réponse b.} La nidation est l'implantation de l'embryon (au stade blastocyste) dans la muqueuse utérine (endomètre), environ 6 à 7 jours après la fécondation. La réponse a définit la fécondation, la c l'ovulation, la d l'accouchement.

\textbf{Points de vigilance :} ne pas confondre \cle{fécondation} (dans la trompe) et \cle{nidation} (dans l'utérus) ; distinguer la cause hormonale de la puberté (axe hypothalamo-hypophysaire) de ses manifestations visibles (caractères sexuels secondaires) ; situer précisément l'ovulation (ovaire) et la fécondation (trompe).
\end{corrige}

\begin{corrige}[Exercice 2 — Vrai ou faux : les IST]

\textbf{1. FAUX.} Le VIH ne se transmet que par trois voies : sexuelle (rapports non protégés), sanguine (sang contaminé, matériel d'injection non stérile) et mère-enfant (grossesse, accouchement, allaitement). Les contacts de la vie quotidienne (poignées de main, embrassades, repas partagés, piqûres de moustiques, piscines, toilettes) ne transmettent jamais le virus.

\textbf{2. VRAI.} Le préservatif (masculin ou féminin) est une méthode de barrière qui empêche à la fois le passage des spermatozoïdes (contraception) et celui des agents pathogènes (prévention des IST, dont le VIH). Les autres méthodes contraceptives (pilule, implant, stérilet, injectable) ne protègent \emph{pas} des IST.

\textbf{3. FAUX.} Une personne séropositive asymptomatique (phase de latence clinique) est contagieuse : le virus est présent dans son sang, son sperme, ses sécrétions vaginales et le lait maternel même en l'absence de tout symptôme. C'est pourquoi le dépistage est essentiel pour connaître son statut.

\textbf{4. VRAI.} La syphilis (bactérie \emph{Treponema pallidum}) et la gonococcie (bactérie \emph{Neisseria gonorrhoeae}) sont des IST bactériennes, guérissables par antibiothérapie adaptée si elles sont traitées précocement. À l'inverse, les IST virales (VIH, herpès, hépatite B, HPV) ne se « guérissent » pas mais se contrôlent.

\textbf{5. VRAI.} Au Cameroun, de nombreux centres de santé, hôpitaux (ex. : Centre Jamot à Yaoundé, Centre de prise en charge du VIH à l'Hôpital Laquintinie de Douala) et ONG proposent un dépistage du VIH gratuit, confidentiel et souvent anonyme, notamment lors des campagnes de sensibilisation (Journée mondiale du sida, vacances sans sida).

\textbf{Point de vigilance :} la justification est exigée pour chaque item au Probatoire ; une réponse « vrai » ou « faux » seule ne rapporte pas tous les points. Citez les voies de transmission, le type d'agent pathogène ou le mode d'action du préservatif.
\end{corrige}

\begin{corrige}[Exercice 3 — Définitions]

\textbf{1. Puberté :} période de transition entre l'enfance et l'âge adulte, au cours de laquelle les gonades, activées par l'axe hypothalamo-hypophysaire (GnRH $\to$ FSH/LH), produisent des hormones sexuelles responsables de l'apparition des caractères sexuels secondaires et de l'acquisition de la capacité de se reproduire (gamétogenèse).

\textbf{2. Grossesse précoce :} grossesse survenant chez une adolescente, généralement avant 19 ans, souvent non désirée et survenant en milieu scolaire ; elle est dite « à haut risque » car le corps de l'adolescente (bassin, utérus, maturité physiologique globale) n'est pas toujours complètement mature pour mener une grossesse sans complications majeures.

\textbf{3. Infection sexuellement transmissible (IST) :} infection causée par un agent pathogène (virus, bactérie, parasite, champignon) transmis principalement lors de rapports sexuels non protégés (vaginaux, anaux, oraux) ; exemples : VIH/Sida, syphilis, gonococcie, chlamydiose, hépatite B, herpès génital, condylomes (HPV).

\textbf{4. Planification familiale :} ensemble des méthodes et des services permettant aux individus et aux couples de choisir librement et en toute connaissance de cause le nombre de leurs enfants et l'espacement des naissances, dans le respect de leurs droits reproductifs.

\textbf{5. Contraception :} ensemble des moyens utilisés pour empêcher temporairement (pilule, préservatif, stérilet, implant, injectable, pilule du lendemain) ou définitivement (ligature des trompes, vasectomie) la survenue d'une grossesse.
\end{corrige}

\begin{corrige}[Exercice 4 — Calculs directs : statistiques de santé scolaire]

\textbf{1.} Pourcentage de filles concernées (base = effectif total des filles) :
\[
\frac{32}{640}\times 100 = 5\ \%.
\]
5 \% des filles du lycée ont été concernées par une grossesse cette année-là.

\textbf{2.} Proportion de déscolarisation parmi les élèves enceintes (base = nombre de grossesses) :
\[
\frac{18}{32}\times 100 = 56{,}25\ \%.
\]
Plus d'une élève enceinte sur deux (56,25 \%) a abandonné ses études.

\textbf{3.} Taux d'évolution entre l'année précédente (40 grossesses) et cette année (32) :
\[
\frac{32-40}{40}\times 100 = \frac{-8}{40}\times 100 = -20\ \%.
\]
Le nombre de grossesses a \textbf{diminué de 20 \%} d'une année sur l'autre, ce qui traduit une amélioration (peut-être liée aux actions de sensibilisation du club de santé), même si le phénomène reste préoccupant.

\textbf{Points de vigilance :} bien prendre l'effectif des \emph{filles} (640) et non l'effectif total (1200) pour la question 1 ; pour un taux d'évolution, diviser par la valeur de \emph{départ} (l'année de référence, ici 40) ; exprimer le résultat en pourcentage avec le signe + ou -.
\end{corrige}

\begin{corrige}[Exercice 5 — Classement des méthodes contraceptives]

\textbf{1.} Tableau complété :

\begin{center}
\begin{tabularx}{\linewidth}{|l|X|}
\hline
\rowcolor{popblueL} \textbf{Catégorie} & \textbf{Méthodes} \\
\hline
Méthode de barrière (mécanique) & préservatif masculin, préservatif féminin \\
\hline
Méthode hormonale & pilule combinée ou microprogestative, implant sous-cutané, injectable, pilule du lendemain (contraception hormonale d'urgence) \\
\hline
Méthode mécanique intra-utérine & stérilet (au cuivre ou hormonal) \\
\hline
Méthode définitive (chirurgicale) & ligature des trompes (femme), vasectomie (homme) \\
\hline
Méthode naturelle / peu fiable & retrait (coït interrompu), méthode des températures (Ogino-Knaus), méthode du calendrier \\
\hline
\end{tabularx}
\end{center}

\textbf{2.} Seuls le \cle{préservatif masculin} et le \cle{préservatif féminin} protègent également des IST : ils constituent une barrière physique (latex ou polyuréthane) qui empêche le contact direct entre les muqueuses et les sécrétions génitales des partenaires, bloquant ainsi le passage des agents pathogènes (VIH, gonocoques, tréponèmes, HPV, herpès...).

\textbf{3.} Le retrait est peu fiable car des spermatozoïdes peuvent être libérés avant l'éjaculation dans les sécrétions pré-éjaculatoires (liquide de Cowper) et parce qu'il exige une maîtrise parfaite du moment de l'éjaculation, rarement réalisée chez l'adolescent. La méthode des températures (et plus généralement les méthodes d'auto-observation du cycle) est peu fiable chez l'adolescente car la date d'ovulation varie d'un cycle à l'autre (stress, maladie, cycles irréguliers fréquents à la puberté) et les spermatozoïdes survivent 3 à 5 jours dans les voies génitales féminines, élargissant la fenêtre de fécondité.
\end{corrige}

\begin{corrige}[Exercice 6 — Schéma à légender : l'appareil reproducteur féminin]

\begin{popfigure}
\begin{tikzpicture}[scale=1.1, every node/.style={font=\small}]
% Vagin
\draw[popdark, thick, fill=poppinkL] (-0.6,0) -- (-0.6,1.2) -- (0.6,1.2) -- (0.6,0) -- cycle;
% Col de l'utérus
\draw[popdark, thick, fill=poppinkL] (-0.4,1.2) .. controls (-0.4,1.6) and (0.4,1.6) .. (0.4,1.2) -- cycle;
% Corps de l'utérus
\draw[popdark, thick, fill=poppinkL] (-1.0,1.6) .. controls (-1.3,3.0) and (1.3,3.0) .. (1.0,1.6) -- cycle;
% Trompes (utérines puis ampullaires)
\draw[popdark, thick] (-1.0,3.0) .. controls (-1.8,3.5) and (-2.5,3.5) .. (-2.8,3.0);
\draw[popdark, thick] (1.0,3.0) .. controls (1.8,3.5) and (2.5,3.5) .. (2.8,3.0);
% Pavillons (fimbriae)
\draw[popdark, thick] (-2.8,3.0) .. controls (-3.0,2.8) and (-3.0,2.5) .. (-2.8,2.3);
\draw[popdark, thick] (2.8,3.0) .. controls (3.0,2.8) and (3.0,2.5) .. (2.8,2.3);
% Ovaires
\draw[popdark, thick, fill=poporangeL] (-3.0,2.0) ellipse (0.5 and 0.35);
\draw[popdark, thick, fill=poporangeL] (3.0,2.0) ellipse (0.5 and 0.35);
% Ligaments (schéma simplifié)
\draw[popdark, dashed, thin] (-1.0,1.6) -- (-1.5,0.5);
\draw[popdark, dashed, thin] (1.0,1.6) -- (1.5,0.5);
% Numéros de légende
\node[popdark, font=\bfseries] at (-3.0, 2.0) {1};
\node[popdark, font=\bfseries] at (-2.2, 3.5) {2};
\node[popdark, font=\bfseries] at (0, 3.3) {3};
\node[popdark, font=\bfseries] at (0.6, 1.4) {4};
\node[popdark, font=\bfseries] at (0.8, 0.6) {5};
% Flèches discrètes pour le trajet ovule
\draw[poporange, ->, thick, >=Stealth] (-2.5, 2.2) .. controls (-2.0, 2.8) .. (-1.5, 3.1);
\draw[poporange, ->, thick, >=Stealth] (-1.0, 3.1) -- (-0.2, 2.8);
\draw[poporange, ->, thick, >=Stealth] (0, 2.5) -- (0, 1.8);
\end{tikzpicture}
\legende{Appareil reproducteur féminin (vue sagittale médiane) : \textbf{1} Ovaire (production ovules/hormones), \textbf{2} Trompe de Fallope (site fécondation, transport embryon), \textbf{3} Utérus (corps utérin, cavité utérine, nidation), \textbf{4} Col de l'utérus (orifice utérin, barrière mucus), \textbf{5} Vagin (réceptacle sperme, voie accouchement). Les flèches oranges indiquent le trajet de l'ovule : ovulation $\to$ captation par pavillon $\to$ trompe (fécondation) $\to$ utérus (nidation).}
\end{popfigure}

\textbf{1.} Légende confirmée par le schéma :
\begin{itemize}
\item \textbf{1} : ovaire ;
\item \textbf{2} : trompe de Fallope (pavillon et ampoule) ;
\item \textbf{3} : utérus (corps et fond utérin) ;
\item \textbf{4} : col de l'utérus ;
\item \textbf{5} : vagin.
\end{itemize}

\textbf{2.} Trajet de l'ovule : au moment de l'ovulation (vers le 14\textsuperscript{e} jour d'un cycle de 28 jours), l'ovaire libère un ovocyte II (ovule) qui est capté par les fimbriae du pavillon de la trompe de Fallope. L'ovule progresse ensuite dans la trompe grâce aux battements ciliaires et à la péristaltique, où il peut être fécondé par un spermatozoïde (fenêtre de fécondité : 12 à 24 h après ovulation). L'embryon ainsi formé (zygote puis blastocyste) migre vers l'utérus (3 à 4 jours) et s'implante dans la muqueuse utérine (endomètre) préparée par la progestérone : c'est la nidation (vers J20-J21 du cycle).

\textbf{3.} La fécondation se produit normalement dans le \cle{tiers externe de la trompe de Fallope (ampoule)}. Si aucun spermatozoïde ne rencontre l'ovule, celui-ci meurt au bout de 24 à 48 heures ; le corps jaune régresse, la progestérone chute, l'endomètre, non maintenu, se désagrège et s'évacue : ce sont les règles (menstruations), marquant le début d'un nouveau cycle ovarien.

\textbf{Points de vigilance :} ne pas situer la fécondation dans l'utérus (erreur fréquente) ; ne pas confondre col de l'utérus (orifice étroit reliant utérus et vagin, siège du frottis de dépistage) et vagin ; l'ovaire n'est pas directement collé à la trompe (espace péritonéal), la captation est active par les fimbriae.
\end{corrige}

\begin{corrige}[Exercice 7 — Étude de cas : une décision difficile]

\textbf{1.} Trois risques graves de l'avortement clandestin (pratiqué hors structure médicale agréée, par personne non qualifiée ou avec des méthodes dangereuses) :
\begin{itemize}
\item \textbf{Risque médical immédiat :} hémorragie grave (pouvant mener au choc hypovolémique et à la mort), infection sévère (septicémie, péritonite) due à l'introduction de germes ou matériel non stérile, perforation de l'utérus ou lésion d'organes voisins (vessie, intestin).
\item \textbf{Risque médical à long terme :} stérilité définitive (séquelles tubaires, adhérences synechiques), douleurs pelviennes chroniques, troubles menstruels, traumatisme psychologique profond (syndrome de stress post-traumatique, dépression).
\item \textbf{Risque juridique et social :} l'avortement clandestin est un délit puni par le Code pénal camerounais (articles 337 à 339), tant pour la praticienne que pour la jeune fille (peines de prison et amendes) ; stigmatisation familiale et communautaire.
\end{itemize}

\textbf{2.} Trois attitudes responsables pour Amina :
\begin{itemize}
\item se rendre sans délai dans un \textbf{centre de santé} ou un hôpital (ex. : Centre de santé de district, Hôpital de district, Hôpital régional) pour une consultation prénatale et un accompagnement médical (sage-femme, médecin) et psychologique (assistante sociale, psychologue) ; la loi camerounaise autorise l'avortement thérapeutique si la vie de la mère est en danger (certificat de deux médecins) ;
\item \textbf{en parler à un adulte de confiance} (parent, tante, grand-mère, enseignant, conseiller d'orientation, infirmier scolaire) plutôt que de rester isolée dans la peur et la honte ; le secret partagé allège la charge mentale ;
\item s'adresser aux \textbf{structures d'accueil et d'appui} spécialisées (ONG comme CAMNAFAW, Association Camerounaise pour le Bien-Être Familial ; centres d'écoute des adolescents ; services sociaux de la Mairie ou de la Délégation régionale des Affaires sociales) qui aident les élèves enceintes à préparer l'avenir (réinsertion scolaire, formation professionnelle, soutien matériel).
\end{itemize}

\textbf{3.} La poursuite de la scolarité est un droit fondamental : aucune élève ne devrait être définitivement exclue pour une grossesse (Circulaire n\textsuperscript{o} 0001/MINESEC/IGE du 2011 relative à la protection de la jeune fille mère à l'école). Des dispositions permettent la suspension temporaire pour raisons de santé (maternité) et la reprise des études après l'accouchement. C'est aussi une chance vitale : l'instruction est le meilleur atout pour l'autonomie future d'Amina et de son enfant, alors que la déscolarisation définitive enferme presque systématiquement dans la précarité économique et la dépendance.
\end{corrige}

\begin{corrige}[Exercice 8 — Expliquer pour convaincre]

\textit{Exemple de production attendue :}

Mon petit frère, tu te trompes : le préservatif n'est pas réservé aux adultes mariés, c'est au contraire le moyen de prévention le mieux adapté aux adolescents. D'abord, il assure une \textbf{double protection} : il empêche à la fois les grossesses non désirées et les infections sexuellement transmissibles, dont le VIH, car il constitue une barrière physique étanche aux spermatozoïdes et aux agents pathogènes (virus, bactéries). Aucune autre méthode contraceptive (pilule, implant, stérilet) ne protège des IST. Ensuite, il est \textbf{accessible} : peu coûteux (quelques centaines de francs CFA), vendu en pharmacie, supermarché, boutique, station-service, il est même distribué \textbf{gratuitement} dans de nombreux centres de santé, au CESC (Centre d'Éducation à la Santé Scolaire) et lors des campagnes de sensibilisation (Vacances sans Sida, Journée mondiale du sida) au Cameroun. Enfin, il est \textbf{sans effet hormonal} : contrairement à la pilule ou à l'implant, il ne perturbe pas l'organisme, ne demande ni ordonnance, ni consultation médicale, ni suivi biologique. Dire « il suffit de faire attention » ou « je me retirerai à temps » est une illusion dangereuse : on ne peut pas voir le VIH ni les autres IST, on ne peut pas prévoir avec certitude l'ovulation d'une partenaire, et le retrait échoue souvent. Le comportement vraiment responsable, c'est de s'informer, de se faire dépister régulièrement et, si l'on est sexuellement actif, d'utiliser \textbf{systématiquement} un préservatif à chaque rapport, du début à la fin. Protège-toi, protège l'autre, protège ton avenir.

\textbf{Points de vigilance :} trois arguments scientifiques minimum exigés (double protection, accessibilité/gratuité, absence d'effet hormonal/non invasif) ; ton adapté au destinataire (fraternel, bienveillant, sans jugement moralisateur) ; conclusion par un conseil de comportement responsable clair (systématicité du préservatif).
\end{corrige}

\begin{corrige}[Exercice 9 — Type Probatoire : les grossesses en milieu scolaire]

\textbf{1.} Taux de grossesses pour 1000 filles (taux brut de grossesse en milieu scolaire) :
\begin{align*}
2021\text{-}2022 : \quad & \frac{135}{\num{4500}}\times 1000 = 30 \text{ pour mille} ;\\
2022\text{-}2023 : \quad & \frac{141}{\num{4700}}\times 1000 = 30 \text{ pour mille} ;\\
2023\text{-}2024 : \quad & \frac{160}{\num{5000}}\times 1000 = 32 \text{ pour mille}.
\end{align*}

\begin{center}
\begin{tabularx}{\linewidth}{|l|X|X|X|}
\hline
\rowcolor{popblueL} \textbf{Année scolaire} & \textbf{2021-2022} & \textbf{2022-2023} & \textbf{2023-2024} \\
\hline
Effectif filles & \num{4500} & \num{4700} & \num{5000} \\
Nombre de grossesses & 135 & 141 & 160 \\
Taux de grossesses (pour 1000 filles) & 30 & 30 & 32 \\
\hline
\end{tabularx}
\end{center}

\textbf{2.} L'effectif des filles augmente régulièrement (\num{4500} $\to$ \num{4700} $\to$ \num{5000}), et le nombre brut de grossesses augmente aussi (135 $\to$ 141 $\to$ 160). Mais le \emph{taux} (pour 1000) montre que la fréquence réelle est restée stable les deux premières années (30 pour mille) puis a légèrement augmenté la troisième année (32 pour mille) : l'augmentation du nombre de grossesses n'est donc pas due \emph{seulement} à l'augmentation des effectifs ; le phénomène s'aggrave légèrement en 2023-2024.

\textbf{3.} Pourcentage moyen de déscolarisation sur les trois ans (moyenne pondérée par les effectifs) :
\[
\frac{81+92+112}{135+141+160}\times 100 = \frac{285}{436}\times 100 \approx 65{,}4\ \%.
\]
Environ deux élèves enceintes sur trois (65,4 \%) abandonnent leurs études. C'est un chiffre alarmant.

\textbf{4.} \textbf{Causes} (deux attendues minimum) : rapports sexuels précoces et non protégés (absence de préservatif) ; méconnaissance de la contraception et du cycle féminin (période fertile) ; pauvreté et relations transactionnelles (sexuelle contre argent/cadeaux) ; pression des pairs et mimétisme ; faible communication parents-enfants sur la sexualité ; mariages précoces (dans certaines régions). \textbf{Conséquences} : sur la santé, grossesse à haut risque (anémie, hypertension gravidique/éclampsie, fistule obstétricale, hémorragie de la délivrance, décès maternel) ; sur l'avenir, déscolarisation massive, précarité économique, rejet social/familial, difficulté d'insertion professionnelle, transmission intergénérationnelle de la pauvreté.

\textbf{5.} Mesures concrètes pour un club de santé scolaire (CESC) :
\begin{itemize}
\item Organiser des séances régulières d'information/débat (causeries éducatives) sur la puberté, le cycle menstruel, la contraception (démonstration préservatif sur pénis en bois) et les IST/VIH ;
\item Inviter du personnel de santé qualifié (sage-femme, infirmier, médecin du Centre de Santé de District) pour des causeries, des séances de conseil et orienter les élèves vers les centres de dépistage et de conseil (CDV) et les services de planification familiale (libre accès, confidentialité) ;
\item Mettre en place un dispositif d'écoute confidentiel (boîte à questions, permanences d'élèves formés « pairs-éducateurs », référent adulte bienveillant) pour briser le tabou et repérer les situations de détresse.
\end{itemize}

\textbf{Barème indicatif (sur 10) :} Q1 : 3 pts (1 pt par taux correct, calcul visible) ; Q2 : 2 pts (comparaison taux vs effectifs, conclusion) ; Q3 : 1,5 pt (calcul pondéré correct) ; Q4 : 2 pts (1 pt causes, 1 pt conséquences) ; Q5 : 1,5 pt (mesures réalistes, adaptées au contexte scolaire).

\textbf{Points de vigilance :} un taux « pour 1000 » exige une multiplication par 1000, pas par 100 ; pour la moyenne de déscolarisation (Q3), additionner d'abord les effectifs bruts (numérateur et dénominateur), ne pas faire la moyenne des trois pourcentages annuels ; la comparaison (Q2) doit porter sur le \emph{taux}, pas seulement sur les nombres bruts.
\end{corrige}

\begin{corrige}[Exercice 10 — Type Probatoire : prévention des IST chez les jeunes]

\textbf{1.} Prévalence du VIH dans l'échantillon :
\[
\frac{60}{\num{2000}}\times 100 = 3\ \%.
\]

\textbf{2.} Proportion de séropositifs ignorant leur état (non-dépistés) :
\[
\frac{45}{60}\times 100 = 75\ \%.
\]
Ces trois quarts de séropositifs non informés de leur statut peuvent transmettre le virus sans le savoir, lors de rapports non protégés. L'ignorance du statut sérologique entretient donc silencieusement la chaîne de transmission de l'épidémie.

\textbf{3.} Modes de transmission du VIH et moyens de prévention associés :
\begin{itemize}
\item \textbf{Voie sexuelle} (rapports vaginaux, anaux, oraux non protégés) : usage \textbf{systématique et correct du préservatif} (masculin ou féminin) à chaque rapport ;
\item \textbf{Voie sanguine} (partage de seringues/aiguilles chez usagers de drogues, matériel de scarification/tatouage/piercing non stérile, transfusion non dépistée, accidents d'exposition au sang) : ne jamais partager seringues ni objets tranchants ; exiger du matériel stérile à usage unique ; transfuser du sang dépisté (norme au Cameroun) ; PEP (prophylaxie post-exposition) en cas d'accident ;
\item \textbf{Transmission mère-enfant} (pendant la grossesse, l'accouchement, l'allaitement) : dépistage systématique du VIH dès le premier trimestre de la grossesse (offert gratuitement en consultation prénatale au Cameroun) et mise sous traitement antirétroviral (ARV) préventif de la mère et du nouveau-né (option B+ : ARV à vie pour la mère).
\end{itemize}

\textbf{4.} Deux autres IST fréquentes chez les jeunes :
\begin{itemize}
\item la \textbf{gonococcie} (ou « chaude-pisse ») : bactérie \emph{Neisseria gonorrhoeae} ; symptôme fréquent chez l'homme : écoulement purulent jaune-vert par l'urètre et brûlures intenses en urinant (souvent asymptomatique chez la femme) ;
\item la \textbf{syphilis} : bactérie \emph{Treponema pallidum} ; symptôme du stade primaire : chancre indolore (ulcération génitale unique, bords nets, fond propre) à la porte d'entrée, disparaissant spontanément mais la maladie continue d'évoluer.
\end{itemize}
(Acceptés aussi : chlamydiose \emph{Chlamydia trachomatis} — souvent asymptomatique, risque stérilité ; hépatite B — virus, ictère, fatigue ; herpès génital \emph{HSV-2} — vésicules douloureuses récidivantes ; condylomes \emph{HPV} — verrues génitales, risque cancer col utérin).

\textbf{5.} Justification argumentée : Le dépistage régulier permet à chacun de connaître son statut sérologique ; une personne dépistée positive peut être mise précocement sous traitement antirétroviral (ARV), ce qui préserve sa santé (charge virale indétectable = non transmissible, concept I=I) et réduit fortement le risque de transmission, tandis qu'une personne dépistée négative est incitée à maintenir des comportements protecteurs. L'usage \textbf{systématique du préservatif}, barrière mécanique efficace contre le VIH et la plupart des autres IST, bloque la transmission lors des rapports sexuels. Combinés, ces deux piliers — \textbf{connaître son statut} (dépistage) et \textbf{se protéger à chaque rapport} (préservatif) — brisent la chaîne de contamination, alors que les données de l'étude montrent que 75 \% des séropositifs ignoraient leur état (réservoir invisible) et que seulement 35 \% des jeunes utilisent toujours le préservatif (protection lacunaire).

\textbf{Barème indicatif (sur 10) :} Q1 : 1 pt ; Q2 : 2 pts (calcul + interprétation) ; Q3 : 3 pts (1 pt par couple voie/prévention correct) ; Q4 : 2 pts (1 pt par IST : agent + symptôme) ; Q5 : 2 pts (argumentation s'appuyant sur les données : 75\% ignorance, 35\% préservatif).

\textbf{Points de vigilance :} distinguer prévalence (proportion de malades dans l'échantillon total = 60/2000) et proportion d'ignorance (calculée sur les 60 séropositifs = 45/60) ; citer un agent pathogène \emph{et} un symptôme caractéristique pour chaque IST ; la question 5 exige un paragraphe argumenté structuré s'appuyant explicitement sur les données chiffrées de l'énoncé.
\end{corrige}

\begin{corrige}[Exercice 11 — Type Probatoire : sujet de synthèse sur la santé reproductive]

\textbf{1.} \textbf{Santé reproductive (OMS) :} état de complet bien-être physique, mental et social concernant l'appareil reproducteur et ses fonctions, à toutes les étapes de la vie, et pas seulement absence de maladie ; elle implique la capacité d'avoir une vie sexuelle satisfaisante et sans risque, la liberté de décider si, quand et à quelle fréquence se reproduire. \textbf{Grossesse précoce :} grossesse survenant chez une adolescente (généralement avant 19 ans), souvent non désirée, alors que son corps n'est pas toujours complètement mature (bassin, utérus, maturité psychologique), la plaçant en situation de « haut risque » obstétrical. \textbf{Planification familiale :} ensemble des méthodes (contraception) et des services (conseil, suivi) permettant aux individus et aux couples de choisir librement et en toute connaissance de cause le nombre de leurs enfants et l'espacement des naissances.

\textbf{2.} Une grossesse à 15 ans est « à haut risque » car (arguments tirés du document 2) : le bassin de l'adolescente est parfois insuffisamment développé (disproportion fœto-pelvienne), ce qui complique l'accouchement par voie basse et augmente le risque de césarienne ou de fistule obstétricale ; elle présente des risques accrus d'anémie ferriprive (besoins accrus) et d'hypertension gravidique (prééclampsie/éclampsie) ; elle peut souffrir de fistule obstétricale (communication anormale vessie/vagin ou rectum/vagin) suite à un travail prolongé obstructif ; et les complications de la grossesse et de l'accouchement sont l'une des principales causes de décès des adolescentes de 15 à 19 ans dans le monde (données OMS).

\textbf{3.} Facteurs dégageables du document 1 (contexte camerounais) : la précocité des rapports sexuels et des unions (mariages précoces, notamment dans les régions du Nord et de l'Extrême-Nord) ; la pauvreté et le faible accès à l'information et à la contraception en milieu rural ; des normes socioculturelles valorisant la maternité précoce et la fécondité comme preuve de féminité ; l'insuffisance de l'éducation sexuelle complète à l'école et en famille.

\textbf{4.} Classement des méthodes du document 3 :
\begin{itemize}
\item \textbf{Méthodes de barrière :} préservatif masculin, préservatif féminin — \emph{seules méthodes protégeant aussi des IST} ;
\item \textbf{Méthodes hormonales :} pilule (combinée ou microprogestative), implant sous-cutané, injectable (DMPA), pilule du lendemain ;
\item \textbf{Méthode intra-utérine :} stérilet (au cuivre ou hormonal au lévonorgestrel).
\end{itemize}

\textbf{5.} \textit{Exemple de production attendue (environ 18 lignes) :}

\textbf{Adopter un comportement responsable en matière de santé reproductive}

Chers camarades, la santé reproductive nous concerne tous, dès l'adolescence. Une grossesse précoce met en danger la vie même de la jeune mère : son corps n'est pas toujours prêt, et les complications de la grossesse sont l'une des premières causes de décès des adolescentes dans le monde. Elle brise aussi un avenir : près de deux élèves enceintes sur trois abandonnent l'école, s'enfonçant dans la précarité. Les infections sexuellement transmissibles, dont le VIH, constituent l'autre grand danger : elles se transmettent lors de rapports non protégés, souvent sans symptôme visible, et certaines sont incurables. Heureusement, des moyens de prévention existent et sont à notre portée : le préservatif, qui protège à la fois des grossesses et des IST, est accessible et gratuit dans de nombreux centres de santé ; le dépistage du VIH est gratuit, confidentiel et rapide au Cameroun ; les méthodes de planification familiale (pilule, implant, stérilet) permettent aux couples de choisir le nombre et l'espacement de leurs enfants. Mais aucun moyen technique ne suffit sans information : parlons-en en famille, écoutons le personnel de santé, participons aux clubs de santé scolaire, refusons les rumeurs et les pressions. Être responsable, c'est s'informer, se protéger, se faire dépister et respecter l'autre. Notre avenir en dépend.

\textbf{Barème indicatif (sur 20) :} Q1 : 3 pts (définitions précises) ; Q2 : 3 pts (arguments explicitement tirés du document 2) ; Q3 : 2 pts (facteurs contextuels document 1) ; Q4 : 4 pts (classement correct + mention protection IST pour barrière) ; Q5 : 8 pts (contenu scientifique 4 pts : risques + prévention + rôle info/services ; organisation/structure 2 pts ; qualité de la langue/ton 2 pts).

\textbf{Points de vigilance :} en Q2, les arguments doivent être explicitement tirés du document 2 cité ; en Q4, ne pas oublier de préciser que \emph{seuls} les préservatifs protègent des IST ; en Q5, le texte doit comporter les trois volets exigés (risques santé/avenir, moyens de prévention accessibles, rôle de l'information/famille/services de santé) et respecter la longueur demandée (15 à 20 lignes).
\end{corrige}

\newpage

\coursec{Fiche bilan}

\courssub{Fiche bilan : Amélioration de la santé reproductive}

\begin{popfigure}
\begin{center}\tcbox[colback=popink,colframe=popink,coltext=white,arc=9pt,boxsep=4pt,left=6pt,right=6pt]{\bfseries\small Santé reproductive}\end{center}
\vspace{-2pt}
\begin{tcbraster}[raster columns=3,raster equal height=rows,raster column skip=6pt,raster row skip=6pt,enhanced,blankest]
\begin{tcolorbox}[enhanced,colback=white,colframe=popblue,boxrule=0.9pt,arc=3pt,title={Puberté et reproduction},fonttitle=\bfseries\scriptsize,coltitle=white,colbacktitle=popblue,left=4pt,right=4pt,top=2pt,bottom=2pt,boxsep=1pt,toptitle=1.5pt,bottomtitle=1.5pt,halign=left]
\begin{itemize}[nosep,leftmargin=*,itemsep=1pt,font=\scriptsize]
\item Garçon : spermatozoïdes, testostérone
\item Fille : ovulation, règles, œstrogènes
\end{itemize}
\end{tcolorbox}
\begin{tcolorbox}[enhanced,colback=white,colframe=poporange,boxrule=0.9pt,arc=3pt,title={Grossesses précoces : causes},fonttitle=\bfseries\scriptsize,coltitle=white,colbacktitle=poporange,left=4pt,right=4pt,top=2pt,bottom=2pt,boxsep=1pt,toptitle=1.5pt,bottomtitle=1.5pt,halign=left]
\begin{itemize}[nosep,leftmargin=*,itemsep=1pt,font=\scriptsize]
\item Rapports non protégés
\item Méconnaissance, pressions, pauvreté
\end{itemize}
\end{tcolorbox}
\begin{tcolorbox}[enhanced,colback=white,colframe=poppurple,boxrule=0.9pt,arc=3pt,title={Conséquences},fonttitle=\bfseries\scriptsize,coltitle=white,colbacktitle=poppurple,left=4pt,right=4pt,top=2pt,bottom=2pt,boxsep=1pt,toptitle=1.5pt,bottomtitle=1.5pt,halign=left]
\begin{itemize}[nosep,leftmargin=*,itemsep=1pt,font=\scriptsize]
\item Santé : fistules, avortements à risque
\item Scolarité et vie sociale
\end{itemize}
\end{tcolorbox}
\begin{tcolorbox}[enhanced,colback=white,colframe=popgreen,boxrule=0.9pt,arc=3pt,title={IST},fonttitle=\bfseries\scriptsize,coltitle=white,colbacktitle=popgreen,left=4pt,right=4pt,top=2pt,bottom=2pt,boxsep=1pt,toptitle=1.5pt,bottomtitle=1.5pt,halign=left]
\begin{itemize}[nosep,leftmargin=*,itemsep=1pt,font=\scriptsize]
\item VIH/Sida, blennorragie, syphilis, chlamydiose
\item Transmission : sexuelle, sang, mère-enfant
\end{itemize}
\end{tcolorbox}
\begin{tcolorbox}[enhanced,colback=white,colframe=poppink,boxrule=0.9pt,arc=3pt,title={Prévention des IST},fonttitle=\bfseries\scriptsize,coltitle=white,colbacktitle=poppink,left=4pt,right=4pt,top=2pt,bottom=2pt,boxsep=1pt,toptitle=1.5pt,bottomtitle=1.5pt,halign=left]
\begin{itemize}[nosep,leftmargin=*,itemsep=1pt,font=\scriptsize]
\item Abstinence, fidélité, préservatif
\item Dépistage, traitement
\end{itemize}
\end{tcolorbox}
\begin{tcolorbox}[enhanced,colback=white,colframe=popgold,boxrule=0.9pt,arc=3pt,title={Planification familiale},fonttitle=\bfseries\scriptsize,coltitle=white,colbacktitle=popgold,left=4pt,right=4pt,top=2pt,bottom=2pt,boxsep=1pt,toptitle=1.5pt,bottomtitle=1.5pt,halign=left]
\begin{itemize}[nosep,leftmargin=*,itemsep=1pt,font=\scriptsize]
\item Méthodes naturelles
\item Méthodes modernes : pilule, DIU, implant
\end{itemize}
\end{tcolorbox}
\end{tcbraster}

\legende{Carte mentale de la leçon : les six grandes idées de la santé reproductive — puberté, causes et conséquences des grossesses précoces, IST, prévention et planification familiale.}
\end{popfigure}

\begin{bilan}
\textbf{Définitions clés à connaître par cœur}
\begin{itemize}
  \item \cle{Puberté} : période de transition entre l'enfance et l'âge adulte, marquée par la maturation des organes reproducteurs et l'apparition des caractères sexuels secondaires, sous le contrôle des hormones sexuelles (testostérone chez le garçon, œstrogènes et progestérone chez la fille).
  \item \cle{Fécondation} : union d'un spermatozoïde et d'un ovule formant une cellule-œuf ; elle est possible dès les premières règles (ovulation) chez la fille et dès la production de spermatozoïdes chez le garçon.
  \item \cle{Grossesse précoce (ou non désirée)} : grossesse survenant chez une adolescente, avant la maturité physiologique et psychosociale, souvent hors de tout projet de maternité.
  \item \cle{IST} : infection transmise principalement lors de rapports sexuels non protégés (VIH/Sida, blennorragie, syphilis, chlamydiose, herpès génital, hépatite B).
  \item \cle{Planification familiale} : ensemble de méthodes permettant aux couples de choisir le nombre de leurs enfants et l'espacement des naissances.
\end{itemize}

\textbf{Tableau récapitulatif des principales IST}
\begin{center}
\begin{tabularx}{\linewidth}{|l|X|X|}
\hline
\rowcolor{popblueL} \textbf{IST} & \textbf{Agent / transmission} & \textbf{Prévention et traitement} \\
\hline
VIH/Sida & Virus ; rapports sexuels, sang, mère-enfant & Pas de guérison ; préservatif, dépistage, antirétroviraux \\
\hline
Blennorragie (gonococcie) & Bactérie ; rapports sexuels & Préservatif ; guérit par antibiotiques \\
\hline
Syphilis & Bactérie (tréponème) ; rapports sexuels, mère-enfant & Préservatif ; guérit par antibiotiques si traitée tôt \\
\hline
Chlamydiose & Bactérie ; rapports sexuels, souvent sans symptômes & Préservatif, dépistage ; antibiotiques \\
\hline
Hépatite B & Virus ; rapports sexuels, sang & Vaccination, préservatif \\
\hline
\end{tabularx}
\end{center}

\textbf{Méthodes express}
\begin{itemize}
  \item \textbf{Expliquer les risques d'une grossesse précoce} : toujours structurer la réponse en trois points — risques pour la \emph{santé} de la mère (corps immature, hémorragies, fistules obstétricales, avortements clandestins mortels), risques pour l'\emph{enfant} (prématurité, faible poids), risques \emph{sociaux et scolaires} (abandon des études, rejet, pauvreté).
  \item \textbf{Justifier la prévention des IST} : citer la règle des trois comportements — \cle{abstinence}, \cle{fidélité} mutuelle, \cle{préservatif} — puis le dépistage et le traitement du couple.
  \item \textbf{Identifier une méthode de planification familiale} : distinguer méthodes \emph{naturelles} (abstinence périodique, méthode des températures) et méthodes \emph{modernes} (préservatif, pilule, injectable, implant, DIU/stérilet) ; seul le préservatif protège aussi des IST.
  \item \textbf{Adopter un comportement responsable} : s'informer auprès d'un centre de santé, refuser les rapports forcés ou précipités, se faire dépister, parler sans tabou avec un adulte de confiance.
\end{itemize}
\end{bilan}

\begin{aretenir}[Checklist avant le Probatoire]
\begin{itemize}
  \item[$\square$] Je sais définir la puberté et citer les hormones sexuelles mâle et femelles.
  \item[$\square$] Je sais expliquer pourquoi une fille peut tomber enceinte dès ses premières règles.
  \item[$\square$] Je connais au moins trois causes des grossesses précoces au Cameroun.
  \item[$\square$] Je sais citer trois conséquences médicales et deux conséquences sociales d'une grossesse précoce.
  \item[$\square$] Je connais au moins quatre IST et leurs modes de transmission.
  \item[$\square$] Je sais que le VIH ne se transmet pas par les poignées de main, les moustiques ou les repas partagés.
  \item[$\square$] Je sais citer les trois comportements de prévention des IST (abstinence, fidélité, préservatif).
  \item[$\square$] Je sais que seul le préservatif protège à la fois des IST et des grossesses non désirées.
  \item[$\square$] Je sais distinguer méthodes naturelles et méthodes modernes de planification familiale, avec un exemple de chaque.
  \item[$\square$] Je sais rédiger une réponse structurée (santé — enfant — société) sur les risques des grossesses précoces.
\end{itemize}
\end{aretenir}

\findecours

\end{document}
